\documentclass[12pt,a4paper]{book}

% Page setup
\usepackage[margin=1in]{geometry}
\usepackage{setspace}

% Math support
\usepackage{amsmath,amssymb,amsthm}

% Hyperlinks (table of contents navigation)
\usepackage{hyperref}
\hypersetup{
    colorlinks=true,
    linkcolor=black,
    citecolor=black,
    urlcolor=blue
}

% Font setup
\usepackage[T1]{fontenc}
\usepackage{lmodern}

% Table of contents style
\usepackage{tocloft}
\renewcommand{\cftchapleader}{\cftdotfill{\cftdotsep}}
\setlength{\cftbeforechapskip}{12pt}
\setlength{\cftbeforesecskip}{3pt}
\setlength{\cftsecindent}{1.5em}
\setlength{\cftsecnumwidth}{2.8em}
\renewcommand{\cftsecpresnum}{\hfill}
\renewcommand{\cftsecaftersnum}{\hspace*{1em}}

% Chapter numbering format
\renewcommand{\thechapter}{\arabic{chapter}}
\renewcommand{\thesection}{\thechapter.\arabic{section}}

% Part style
\usepackage{titlesec}
\titleformat{\part}[display]
  {\normalfont\Large\centering\scshape}
  {Part \thepart}{0.5em}{}
\titlespacing*{\part}{0pt}{40pt}{30pt}

% Chapter style
\titleformat{\chapter}[display]
  {\normalfont\Large\centering}
  {\thechapter}{0.5em}{\MakeUppercase}
\titlespacing*{\chapter}{0pt}{40pt}{30pt}

% Section style
\titleformat{\section}{\normalfont\large\bfseries}{\thesection}{1em}{}

% Headers and footers
\usepackage{fancyhdr}
\pagestyle{fancy}
\fancyhf{}
\fancyhead[L]{\small\leftmark}
\fancyhead[R]{\small\thepage}
\renewcommand{\headrulewidth}{0.4pt}

% Axiom environment
\newcounter{axiomcounter}[chapter]
\renewcommand{\theaxiomcounter}{\thechapter.\arabic{axiomcounter}}
\newenvironment{axiom}[1]{%
  \refstepcounter{axiomcounter}%
  \par\noindent\textbf{\textsc{Axiom \theaxiomcounter. #1.}}\quad%
}{\par}

% Definition environment
\newcounter{definitioncounter}[chapter]
\renewcommand{\thedefinitioncounter}{\thechapter.\arabic{definitioncounter}}
\newenvironment{definition}[1]{%
  \refstepcounter{definitioncounter}%
  \par\noindent\textbf{DEFINITION \thedefinitioncounter. #1}\quad%
}{\par}

% Theorem environment
\newcounter{theoremcounter}[chapter]
\renewcommand{\thetheoremcounter}{\thechapter.\arabic{theoremcounter}}
\newenvironment{theorem}[1]{%
  \refstepcounter{theoremcounter}%
  \par\noindent\textbf{THEOREM \thetheoremcounter. #1}\quad%
}{\par}

% Note environment
\newenvironment{note}{%
  \par\noindent\textit{Note:}\quad%
}{\par}

\begin{document}

% Title page
\begin{titlepage}
    \centering
    \vspace*{3cm}
    {\Huge\bfseries CALCULUS\par}
    \vspace{1cm}
    {\Large Volume I\par}
    \vspace{0.5cm}
    {\large One-Variable Calculus, with an Introduction to Linear Algebra\par}
    \vspace{2cm}
    {\Large Second Edition\par}
    \vspace{3cm}
    {\large Tom M. Apostol\par}
    \vfill
\end{titlepage}

% Table of contents
\tableofcontents
\newpage

% ====================
% Part I: Introduction
% ====================
\part{Introduction}

\chapter{Introduction}

\section{The Two Basic Concepts of Calculus}

Calculus consists of two main branches: \textbf{integral calculus} and \textbf{differential calculus}.

\begin{enumerate}
    \item \textbf{The area problem (Integral Calculus):} To assign a number which measures the area of a region bounded by a curve.
    \item \textbf{The tangent problem (Differential Calculus):} To assign a number which measures the steepness of a tangent line to a curve.
\end{enumerate}

\chapter{Set Theory}

\section{Set Theory}

\begin{definition}{Set}
A \textit{set} is a collection of objects, called \textit{elements} or \textit{members} of the set.
\end{definition}

\begin{definition}{Subset}
A set $A$ is called a \textit{subset} of a set $B$, denoted by $A \subseteq B$, if every element of $A$ also belongs to $B$.
\end{definition}

\begin{definition}{Union}
The \textit{union} of two sets $A$ and $B$, denoted by $A \cup B$, is the set of all elements that belong to $A$ or to $B$ (or to both).
\end{definition}

\begin{definition}{Intersection}
The \textit{intersection} of two sets $A$ and $B$, denoted by $A \cap B$, is the set of all elements that belong to both $A$ and $B$.
\end{definition}

\begin{definition}{Complement}
The \textit{complement} of a set $A$ relative to a set $B$, denoted by $B - A$ or $B \setminus A$, is the set of all elements that belong to $B$ but not to $A$.
\end{definition}

\chapter{Axioms for Real Numbers}

\section{A Set of Axioms for the Real-Number System}

\subsection{The Field Axioms}

\begin{axiom}{Commutative Laws}
$x + y = y + x$, \quad $xy = yx$.
\end{axiom}

\begin{axiom}{Associative Laws}
$x + (y + z) = (x + y) + z$, \quad $x(yz) = (xy)z$.
\end{axiom}

\begin{axiom}{Distributive Law}
$x(y + z) = xy + xz$.
\end{axiom}

\begin{axiom}{Existence of Identity Elements}
There exist two distinct real numbers, which we denote by $0$ and $1$, such that for every real $x$ we have $x + 0 = x$ and $1 \cdot x = x$.
\end{axiom}

\begin{axiom}{Existence of Negatives}
For every real number $x$ there is a real number $y$ such that $x + y = 0$.
\end{axiom}

\begin{axiom}{Existence of Reciprocals}
For every real number $x \neq 0$ there is a real number $y$ such that $xy = 1$.
\end{axiom}

\begin{note}
The numbers $0$ and $1$ in Axioms 5 and 6 are those of Axiom 4.
\end{note}

\begin{theorem}{}
\begin{enumerate}
    \item[(a)] If $a + b = a + c$, then $b = c$. \textit{(Cancellation law for addition)}
    \item[(b)] If $ab = ac$ and $a \neq 0$, then $b = c$. \textit{(Cancellation law for multiplication)}
\end{enumerate}
\end{theorem}

\textit{Proof.} (a) Given $a + b = a + c$. By Axiom 5 there exists a number $y$ such that $a + y = 0$. Adding $y$ to both sides, we get $y + (a + b) = y + (a + c)$. By associativity, $(y + a) + b = (y + a) + c$, so $0 + b = 0 + c$, thus $b = c$.

(b) Given $ab = ac$ and $a \neq 0$. By Axiom 6 there exists a number $z$ such that $az = 1$. Multiplying both sides by $z$, we get $z(ab) = z(ac)$. By associativity, $(za)b = (za)c$, so $1 \cdot b = 1 \cdot c$, thus $b = c$. \qed

\begin{theorem}{}
For any real numbers $a$ and $b$, there is exactly one real number $x$ such that $a + x = b$.
\end{theorem}

\textit{Proof.} By Axiom 5, there exists $y$ such that $a + y = 0$. Let $x = y + b$. Then $a + x = a + (y + b) = (a + y) + b = 0 + b = b$. For uniqueness, suppose $a + x = b$ and $a + x' = b$. Then $a + x = a + x'$, so by cancellation, $x = x'$. \qed

\begin{theorem}{}
For any real numbers $a$ and $b$ with $a \neq 0$, there is exactly one real number $x$ such that $ax = b$.
\end{theorem}

\textit{Proof.} Similar to Theorem 1.2, using Axiom 6. \qed

\begin{theorem}{}
If $a \cdot 0 = 0$ for every real number $a$.
\end{theorem}

\textit{Proof.} We have $a \cdot 0 + a \cdot 0 = a(0 + 0) = a \cdot 0 = a \cdot 0 + 0$. By cancellation, $a \cdot 0 = 0$. \qed

\begin{theorem}{}
If $ab = 0$, then $a = 0$ or $b = 0$.
\end{theorem}

\textit{Proof.} Suppose $ab = 0$ and $a \neq 0$. Then $b = 1 \cdot b = (a^{-1}a)b = a^{-1}(ab) = a^{-1} \cdot 0 = 0$. \qed

\subsection{The Order Axioms}

Let $\mathbb{R}^+$ denote the set of positive real numbers.

\begin{axiom}{Closure of $\mathbb{R}^+$}
If $x$ and $y$ are in $\mathbb{R}^+$, so are $x+y$ and $xy$.
\end{axiom}

\begin{axiom}{Trichotomy}
For every real $x \neq 0$, either $x \in \mathbb{R}^+$ or $-x \in \mathbb{R}^+$, but not both.
\end{axiom}

\begin{axiom}{Nonmembership of Zero}
$0 \notin \mathbb{R}^+$.
\end{axiom}

\begin{definition}{Inequality}
$a < b$ means $b - a \in \mathbb{R}^+$.
\end{definition}

\begin{theorem}{}
\begin{enumerate}
    \item[(a)] If $a < b$ and $b < c$, then $a < c$. \textit{(Transitivity)}
    \item[(b)] If $a < b$, then $a + c < b + c$.
    \item[(c)] If $a < b$ and $c > 0$, then $ac < bc$.
\end{enumerate}
\end{theorem}

\textit{Proof.} (a) If $b - a \in \mathbb{R}^+$ and $c - b \in \mathbb{R}^+$, then $(b - a) + (c - b) = c - a \in \mathbb{R}^+$ by closure.

(b) $(b + c) - (a + c) = b - a \in \mathbb{R}^+$.

(c) If $c > 0$, then $c \in \mathbb{R}^+$, so $bc - ac = (b - a)c \in \mathbb{R}^+$ by closure. \qed

\subsection{Integers and Rational Numbers}

\begin{definition}{Inductive Set}
A set of real numbers is called an \textit{inductive set} if (a) the number $1$ is in the set, and (b) for every $x$ in the set, the number $x+1$ is also in the set.
\end{definition}

\begin{definition}{Positive Integers ($\mathbb{P}$)}
A real number is called a \textit{positive integer} if it belongs to every inductive set. $\mathbb{P}$ is the smallest inductive set.
\end{definition}

\begin{definition}{Rational Numbers ($\mathbb{Q}$)}
Quotients of integers $a/b$ (where $b \neq 0$). The set $\mathbb{Q}$ contains $\mathbb{Z}$ as a subset and forms an ordered field.
\end{definition}

\subsection{Upper Bounds and Supremum}

\begin{definition}{Upper Bound}
Let $S$ be a nonempty set of real numbers. If there is a number $B$ such that $x \leq B$ for every $x$ in $S$, then $S$ is said to be \textit{bounded above} by $B$. The number $B$ is called an \textit{upper bound} for $S$.
\end{definition}

\begin{definition}{Maximum Element}
If an upper bound $B$ is also a member of $S$, then $B$ is called the \textit{largest member} or the \textit{maximum element} of $S$, denoted by $B = \max S$.
\end{definition}

\begin{definition}{Least Upper Bound (Supremum)}
A number $B$ is called a \textit{least upper bound} of a nonempty set $S$ if: (a) $B$ is an upper bound for $S$; (b) No number less than $B$ is an upper bound for $S$. We write $B = \sup S$.
\end{definition}

\begin{axiom}{Least-Upper-Bound Axiom (Completeness Axiom)}
Every nonempty set $S$ of real numbers which is bounded above has a supremum; that is, there is a real number $B$ such that $B = \sup S$.
\end{axiom}

\begin{definition}{Lower Bound}
A set $S$ is bounded below if there exists a number $L$ such that $L \leq x$ for all $x \in S$.
\end{definition}

\begin{definition}{Minimum Element}
If a lower bound $L$ is also a member of $S$, then $L$ is called the \textit{smallest member} or \textit{minimum element} of $S$, denoted by $L = \min S$.
\end{definition}

\begin{definition}{Greatest Lower Bound (Infimum)}
A number $L$ is called a \textit{greatest lower bound} (or \textit{infimum}) of $S$ if: (a) $L$ is a lower bound for $S$; (b) no number greater than $L$ is a lower bound for $S$. We denote it by $\inf S$.
\end{definition}

\begin{theorem}{}
Every nonempty set $S$ that is bounded below has a greatest lower bound; that is, there is a real number $L$ such that $L = \inf S$.
\end{theorem}

\textit{Proof.} Let $-S$ denote the set of negatives of numbers in $S$. Then $-S$ is nonempty and bounded above. By the completeness axiom, there is a number $B$ which is a supremum for $-S$. Then $-B = \inf S$. \qed

\subsection{The Archimedean Property}

\begin{theorem}{}
The set $\mathbb{P}$ of positive integers $1, 2, 3, \ldots$ is unbounded above.
\end{theorem}

\textit{Proof.} Assume $\mathbb{P}$ is bounded above. Since $\mathbb{P}$ is nonempty, by the completeness axiom $\mathbb{P}$ has a least upper bound, say $b$. The number $b - 1$, being less than $b$, cannot be an upper bound for $\mathbb{P}$. Hence, there is at least one positive integer $n$ such that $n > b - 1$. For this $n$ we have $n + 1 > b$. Since $n + 1$ is in $\mathbb{P}$, this contradicts the fact that $b$ is an upper bound for $\mathbb{P}$. \qed

\begin{theorem}{Archimedean Property}
If $x > 0$ and if $y$ is an arbitrary real number, there exists a positive integer $n$ such that $nx > y$.
\end{theorem}

\textit{Proof.} If $nx \leq y$ for all $n$, then $n \leq y/x$ for all $n$, contradicting the unboundedness of $\mathbb{P}$. \qed

\begin{theorem}{}
If three real numbers $a$, $x$, and $y$ satisfy the inequalities
\[
a \leq x \leq a + \frac{y}{n}
\]
for every integer $n \geq 1$, then $x = a$.
\end{theorem}

\textit{Proof.} If $x > a$, by the Archimedean property there is a positive integer $n$ satisfying $n(x - a) > y$, contradicting $x \leq a + y/n$. \qed

\chapter{Induction and Summation}

\section{Mathematical Induction, Summation Notation, and Related Topics}

\begin{theorem}{Principle of Mathematical Induction}
Let $P(n)$ be a proposition about the positive integer $n$. If:
\begin{enumerate}
    \item[(a)] $P(1)$ is true, and
    \item[(b)] Whenever $P(k)$ is true for some positive integer $k$, $P(k+1)$ is also true,
\end{enumerate}
then $P(n)$ is true for all positive integers $n$.
\end{theorem}

\textit{Proof.} Let $S$ be the set of positive integers for which $P(n)$ is true. By (a), $1 \in S$. By (b), if $k \in S$, then $k+1 \in S$. Thus $S$ is an inductive set, so $\mathbb{P} \subseteq S$. Since $S \subseteq \mathbb{P}$ by definition, $S = \mathbb{P}$. \qed

\begin{definition}{Summation Notation}
If $a_1, a_2, \ldots, a_n$ are real numbers, we define:
\[
\sum_{k=1}^{n} a_k = a_1 + a_2 + \cdots + a_n
\]
\end{definition}

\begin{theorem}{Properties of Summation}
\begin{enumerate}
    \item[(a)] $\displaystyle\sum_{k=1}^{n} (a_k + b_k) = \sum_{k=1}^{n} a_k + \sum_{k=1}^{n} b_k$
    \item[(b)] $\displaystyle\sum_{k=1}^{n} ca_k = c\sum_{k=1}^{n} a_k$ for any constant $c$
    \item[(c)] $\displaystyle\sum_{k=1}^{n} c = nc$ for any constant $c$
\end{enumerate}
\end{theorem}

\textit{Proof.} Follows directly from the commutative, associative, and distributive laws. \qed

\begin{definition}{Absolute Value}
For any real number $x$, the \textit{absolute value} of $x$, denoted by $|x|$, is defined by:
\[
|x| = \begin{cases} x & \text{if } x \geq 0 \\ -x & \text{if } x < 0 \end{cases}
\]
\end{definition}

\begin{theorem}{Properties of Absolute Value}
\begin{enumerate}
    \item[(a)] $|x| \geq 0$ for all $x$, and $|x| = 0$ if and only if $x = 0$
    \item[(b)] $|-x| = |x|$ for all $x$
    \item[(c)] $|xy| = |x||y|$ for all $x, y$
    \item[(d)] $|x/y| = |x|/|y|$ for all $x$ and $y \neq 0$
\end{enumerate}
\end{theorem}

\textit{Proof.} Immediate from the definition. \qed

\begin{theorem}{Triangle Inequality}
For all real numbers $x$ and $y$:
\[
|x + y| \leq |x| + |y|
\]
\end{theorem}

\textit{Proof.} We have $-|x| \leq x \leq |x|$ and $-|y| \leq y \leq |y|$. Adding these inequalities gives $-(|x| + |y|) \leq x + y \leq |x| + |y|$. This is equivalent to $|x + y| \leq |x| + |y|$. \qed

\begin{theorem}{}
For all real numbers $x$ and $y$:
\[
||x| - |y|| \leq |x - y|
\]
\end{theorem}

\textit{Proof.} We have $|x| = |(x - y) + y| \leq |x - y| + |y|$, so $|x| - |y| \leq |x - y|$. Interchanging $x$ and $y$ gives $|y| - |x| \leq |y - x| = |x - y|$. Combining these gives $||x| - |y|| \leq |x - y|$. \qed

\begin{theorem}{Bernoulli's Inequality}
If $x > -1$ and $n$ is a positive integer, then
\[
(1 + x)^n \geq 1 + nx
\]
with equality if and only if $n = 1$ or $x = 0$.
\end{theorem}

\textit{Proof.} By induction. For $n = 1$, equality holds. Assume $(1 + x)^k \geq 1 + kx$. Then $(1 + x)^{k+1} = (1 + x)^k(1 + x) \geq (1 + kx)(1 + x) = 1 + (k+1)x + kx^2 \geq 1 + (k+1)x$. Equality holds only if $k = 0$ or $x = 0$. \qed

\begin{theorem}{Cauchy-Schwarz Inequality}
For any real numbers $a_1, \ldots, a_n$ and $b_1, \ldots, b_n$:
\[
\left(\sum_{k=1}^{n} a_k b_k\right)^2 \leq \left(\sum_{k=1}^{n} a_k^2\right)\left(\sum_{k=1}^{n} b_k^2\right)
\]
\end{theorem}

\textit{Proof.} Consider $\sum_{k=1}^{n}(a_k x + b_k)^2 = x^2\sum a_k^2 + 2x\sum a_k b_k + \sum b_k^2 \geq 0$ for all $x$. This quadratic in $x$ is nonnegative, so its discriminant is nonpositive: $4(\sum a_k b_k)^2 - 4(\sum a_k^2)(\sum b_k^2) \leq 0$. \qed

% ====================
% Part II: Integral Calculus
% ====================
\part{Integral Calculus}

\chapter{The Integral}

\section{Functions}

\begin{definition}{Function}
A \textit{function} $f$ from a set $A$ to a set $B$ is a set of ordered pairs $(x, y)$ with $x \in A$ and $y \in B$ such that no two distinct pairs have the same first element. If $(x, y) \in f$, we write $y = f(x)$. The set $A$ is called the \textit{domain} of $f$, and the set of all $y$ such that $y = f(x)$ for some $x$ is called the \textit{range} of $f$.
\end{definition}

\section{The Concept of Area as a Set Function}

\begin{definition}{Area Function}
An \textit{area function} is a nonnegative set function $a$ defined on a collection of sets in the plane, satisfying the following properties:
\begin{enumerate}
    \item[(1)] \textit{Additive property:} If $S$ and $T$ are in the collection and $S \cap T = \emptyset$, then $a(S \cup T) = a(S) + a(T)$.
    \item[(2)] \textit{Comparison property:} If $S$ and $T$ are in the collection with $S \subseteq T$, then $a(S) \leq a(T)$.
\end{enumerate}
\end{definition}

\section{Intervals and Ordinate Sets}

\begin{definition}{Interval}
Let $a < b$. The set of all $x$ satisfying $a < x < b$ is called the \textit{open interval} $(a, b)$. The set of all $x$ satisfying $a \leq x \leq b$ is called the \textit{closed interval} $[a, b]$. The sets $[a, b)$ and $(a, b]$ are called \textit{half-open intervals}.
\end{definition}

\begin{definition}{Ordinate Set}
Let $f$ be a nonnegative function defined on $[a, b]$. The \textit{ordinate set} of $f$ over $[a, b]$ is the set:
\[
S = \{(x, y) \mid a \leq x \leq b, \ 0 \leq y \leq f(x)\}
\]
\end{definition}

\section{Partitions and Step Functions}

\begin{definition}{Partition}
A \textit{partition} $P$ of an interval $[a, b]$ is a set of points $P = \{x_0, x_1, \ldots, x_n\}$ such that $a = x_0 < x_1 < \cdots < x_n = b$. The $k$-th subinterval is $[x_{k-1}, x_k]$ and has length $\Delta x_k = x_k - x_{k-1}$.
\end{definition}

\begin{definition}{Step Function}
A function $s$ defined on $[a, b]$ is called a \textit{step function} if there exists a partition $P = \{x_0, x_1, \ldots, x_n\}$ of $[a, b]$ such that $s$ is constant on each open subinterval $(x_{k-1}, x_k)$.
\end{definition}

\section{The Integral of Step Functions}

\begin{definition}{Integral of a Step Function}
Let $s$ be a step function defined on $[a, b]$ which takes the constant value $s_k$ on the $k$-th open subinterval $(x_{k-1}, x_k)$ of a partition $P$. The \textit{integral} of $s$ from $a$ to $b$ is defined by:
\[
\int_a^b s(x)\, dx = \sum_{k=1}^{n} s_k \Delta x_k
\]
\end{definition}

\begin{theorem}{}
The integral of a step function is independent of the choice of partition.
\end{theorem}

\textit{Proof.} Suppose $s$ is constant on the open subintervals of partition $P$ and also on the open subintervals of a refinement $P'$. The values $s_k$ assigned to each subinterval are the same whether we use $P$ or $P'$, because $s$ is constant on any subinterval where it is defined. The sum $\sum s_k \Delta x_k$ gives the same value for both partitions. \qed

\begin{theorem}{Additivity with Respect to the Interval}
If $s$ is a step function on $[a, b]$ and $a < c < b$, then:
\[
\int_a^b s(x)\, dx = \int_a^c s(x)\, dx + \int_c^b s(x)\, dx
\]
\end{theorem}

\textit{Proof.} We can choose a partition that includes $c$ as a partition point. Then the sum $\sum_{k=1}^{n} s_k \Delta x_k$ splits naturally into two parts: those subintervals in $[a, c]$ and those in $[c, b]$. \qed

\begin{theorem}{Homogeneous Property}
For any step function $s$ and any constant $c$:
\[
\int_a^b c \cdot s(x)\, dx = c \int_a^b s(x)\, dx
\]
\end{theorem}

\textit{Proof.} If $s(x) = s_k$ on $(x_{k-1}, x_k)$, then $c \cdot s(x) = c \cdot s_k$ on the same interval. Therefore:
\[
\int_a^b c \cdot s(x)\, dx = \sum_{k=1}^{n} c \cdot s_k \Delta x_k = c \sum_{k=1}^{n} s_k \Delta x_k = c \int_a^b s(x)\, dx
\] \qed

\begin{theorem}{Additive Property}
If $s$ and $t$ are step functions on $[a, b]$, then:
\[
\int_a^b [s(x) + t(x)]\, dx = \int_a^b s(x)\, dx + \int_a^b t(x)\, dx
\]
\end{theorem}

\textit{Proof.} Let $P$ be a partition fine enough so that both $s$ and $t$ are constant on each open subinterval. If $s(x) = s_k$ and $t(x) = t_k$ on $(x_{k-1}, x_k)$, then $s(x) + t(x) = s_k + t_k$ on this interval. Therefore:
\[
\int_a^b [s(x) + t(x)]\, dx = \sum_{k=1}^{n} (s_k + t_k) \Delta x_k = \sum_{k=1}^{n} s_k \Delta x_k + \sum_{k=1}^{n} t_k \Delta x_k = \int_a^b s(x)\, dx + \int_a^b t(x)\, dx
\] \qed

\begin{theorem}{Comparison Theorem for Step Functions}
If $s$ and $t$ are step functions on $[a, b]$ with $s(x) \leq t(x)$ for all $x \in [a, b]$, then:
\[
\int_a^b s(x)\, dx \leq \int_a^b t(x)\, dx
\]
\end{theorem}

\textit{Proof.} Since $s(x) \leq t(x)$ on each subinterval, we have $s_k \leq t_k$ for each $k$. Therefore $s_k \Delta x_k \leq t_k \Delta x_k$, and summing gives the result. \qed

\section{Upper and Lower Integrals}

\begin{definition}{Upper and Lower Step Functions}
Let $f$ be a bounded function on $[a, b]$. A step function $s$ is said to be \textit{below} $f$ (or a \textit{lower step function}) if $s(x) \leq f(x)$ for all $x \in [a, b]$. A step function $t$ is said to be \textit{above} $f$ (or an \textit{upper step function}) if $f(x) \leq t(x)$ for all $x \in [a, b]$.
\end{definition}

\begin{definition}{Lower and Upper Integrals}
Let $f$ be a bounded function on $[a, b]$. The \textit{lower integral} of $f$ is defined by:
\[
I(f) = \sup\left\{\int_a^b s(x)\, dx \mid s \text{ is a step function below } f\right\}
\]
The \textit{upper integral} of $f$ is defined by:
\[
\bar{I}(f) = \inf\left\{\int_a^b t(x)\, dx \mid t \text{ is a step function above } f\right\}
\]
\end{definition}

\begin{theorem}{}
For any bounded function $f$ on $[a, b]$, we have $I(f) \leq \bar{I}(f)$.
\end{theorem}

\textit{Proof.} If $s$ is below $f$ and $t$ is above $f$, then $s(x) \leq f(x) \leq t(x)$ for all $x$, so $s(x) \leq t(x)$. By the comparison theorem for step functions, $\int_a^b s \leq \int_a^b t$. Taking supremum over all $s$ below $f$ gives $I(f) \leq \int_a^b t$ for any $t$ above $f$. Taking infimum over all $t$ above $f$ gives $I(f) \leq \bar{I}(f)$. \qed

\section{The Definition of the Integral for General Functions}

\begin{definition}{Integrable Function}
A bounded function $f$ defined on $[a, b]$ is said to be \textit{integrable} on $[a, b]$ if $I(f) = \bar{I}(f)$. In this case, the common value is called the \textit{integral} of $f$ from $a$ to $b$, denoted by $\int_a^b f(x)\, dx$.
\end{definition}

\begin{definition}{Area of an Ordinate Set}
If $f$ is nonnegative and integrable on $[a, b]$, the \textit{area} of the ordinate set of $f$ over $[a, b]$ is defined to be $\int_a^b f(x)\, dx$.
\end{definition}

\section{Monotonic and Piecewise Monotonic Functions}

\begin{definition}{Monotonic Function}
A function $f$ defined on an interval is said to be \textit{increasing} on the interval if $f(x) \leq f(y)$ whenever $x < y$. It is said to be \textit{decreasing} if $f(x) \geq f(y)$ whenever $x < y$. A function is \textit{monotonic} if it is either increasing or decreasing.
\end{definition}

\begin{definition}{Piecewise Monotonic Function}
A function $f$ defined on $[a, b]$ is called \textit{piecewise monotonic} if there exists a partition $P = \{x_0, x_1, \ldots, x_n\}$ of $[a, b]$ such that $f$ is monotonic on each open subinterval $(x_{k-1}, x_k)$.
\end{definition}

\begin{theorem}{Integrability of Bounded Monotonic Functions}
If $f$ is bounded and monotonic on $[a, b]$, then $f$ is integrable on $[a, b]$.
\end{theorem}

\textit{Proof.} Assume $f$ is increasing (the proof for decreasing is similar). Let $P = \{x_0, x_1, \ldots, x_n\}$ be a partition of $[a, b]$ into $n$ equal subintervals, each of length $\Delta x = (b-a)/n$. Define step functions $s$ and $t$ as follows: On the $k$-th open subinterval $(x_{k-1}, x_k)$, let $s(x) = f(x_{k-1})$ and $t(x) = f(x_k)$. Since $f$ is increasing, $s(x) \leq f(x) \leq t(x)$ on each subinterval.

The integrals of these step functions are:
\[
\int_a^b s = \sum_{k=1}^{n} f(x_{k-1}) \Delta x \quad \text{and} \quad \int_a^b t = \sum_{k=1}^{n} f(x_k) \Delta x
\]

Their difference is:
\[
\int_a^b t - \int_a^b s = \sum_{k=1}^{n} [f(x_k) - f(x_{k-1})] \Delta x = \Delta x \sum_{k=1}^{n} [f(x_k) - f(x_{k-1})]
\]

This is a telescoping sum equal to $\Delta x [f(b) - f(a)] = \frac{b-a}{n}[f(b) - f(a)]$.

Given $\epsilon > 0$, choose $n$ large enough so that $\frac{b-a}{n}[f(b) - f(a)] < \epsilon$. Then $\int_a^b t - \int_a^b s < \epsilon$. Since $I(f)$ and $\bar{I}(f)$ lie between these two values, we have $\bar{I}(f) - I(f) < \epsilon$. Since $\epsilon$ is arbitrary, $I(f) = \bar{I}(f)$, so $f$ is integrable. \qed

\section{Calculation of the Integral}

\begin{theorem}{}
If $f$ is constant on $[a, b]$ with $f(x) = c$ for all $x$, then:
\[
\int_a^b f(x)\, dx = c(b-a)
\]
\end{theorem}

\textit{Proof.} The function $f$ itself is a step function with respect to any partition, so $\int_a^b f = c(b-a)$. \qed

\begin{theorem}{}
If $f(x) = x$ for all $x \in [a, b]$, then:
\[
\int_a^b x\, dx = \frac{b^2 - a^2}{2}
\]
\end{theorem}

\textit{Proof.} Consider the partition $P = \{x_0, x_1, \ldots, x_n\}$ where $x_k = a + \frac{k(b-a)}{n}$. Then $\Delta x = \frac{b-a}{n}$ and $x_k = a + k\Delta x$.

For the lower step function $s$ with $s(x) = x_{k-1}$ on $(x_{k-1}, x_k)$:
\[
\int_a^b s = \sum_{k=1}^{n} x_{k-1} \Delta x = \sum_{k=1}^{n} [a + (k-1)\Delta x] \Delta x = a(b-a) + (\Delta x)^2 \sum_{k=1}^{n} (k-1)
\]

Since $\sum_{k=1}^{n} (k-1) = \frac{n(n-1)}{2}$, we get:
\[
\int_a^b s = a(b-a) + \frac{(b-a)^2}{n^2} \cdot \frac{n(n-1)}{2} = a(b-a) + \frac{(b-a)^2}{2}\left(1 - \frac{1}{n}\right)
\]

Taking $n \to \infty$: $I(f) = a(b-a) + \frac{(b-a)^2}{2} = \frac{b^2 - a^2}{2}$.

Similarly, for the upper step function, $\bar{I}(f) = \frac{b^2 - a^2}{2}$. Thus $\int_a^b x\, dx = \frac{b^2 - a^2}{2}$. \qed

\begin{theorem}{}
If $p$ is a positive integer and $b > 0$, then:
\[
\int_0^b x^p\, dx = \frac{b^{p+1}}{p+1}
\]
\end{theorem}

\textit{Proof.} Use the partition $x_k = \frac{kb}{n}$ for $k = 0, 1, \ldots, n$, so $\Delta x = \frac{b}{n}$. For the lower step function with $s(x) = x_{k-1}^p$ on $(x_{k-1}, x_k)$:
\[
\int_0^b s = \sum_{k=1}^{n} \left(\frac{(k-1)b}{n}\right)^p \cdot \frac{b}{n} = \frac{b^{p+1}}{n^{p+1}} \sum_{k=1}^{n} (k-1)^p = \frac{b^{p+1}}{n^{p+1}} \sum_{j=0}^{n-1} j^p
\]

Using the formula $\sum_{j=1}^{n-1} j^p = \frac{(n-1)^{p+1}}{p+1} + O(n^p)$, as $n \to \infty$:
\[
\int_0^b s \to \frac{b^{p+1}}{p+1}
\]

Similarly for the upper integral, giving $\int_0^b x^p\, dx = \frac{b^{p+1}}{p+1}$. \qed

\section{Basic Properties of the Integral}

\begin{theorem}{Additivity Property}
If $f$ is integrable on $[a, b]$ and $a < c < b$, then $f$ is integrable on $[a, c]$ and $[c, b]$, and:
\[
\int_a^b f(x)\, dx = \int_a^c f(x)\, dx + \int_c^b f(x)\, dx
\]
\end{theorem}

\textit{Proof.} Let $s_1$ be a step function below $f$ on $[a, c]$ and $s_2$ a step function below $f$ on $[c, b]$. Then $s$ defined by $s(x) = s_1(x)$ on $[a, c)$ and $s(x) = s_2(x)$ on $[c, b]$ is a step function below $f$ on $[a, b]$, and:
\[
\int_a^b s = \int_a^c s_1 + \int_c^b s_2
\]

Taking supremum over all such $s_1$ and $s_2$ gives $I_{[a,b]}(f) \geq I_{[a,c]}(f) + I_{[c,b]}(f)$. Similarly, $\bar{I}_{[a,b]}(f) \leq \bar{I}_{[a,c]}(f) + \bar{I}_{[c,b]}(f)$. Since $f$ is integrable on $[a, b]$, these give equality. \qed

\begin{theorem}{Homogeneous Property}
If $f$ is integrable on $[a, b]$ and $c$ is any real number, then $cf$ is integrable on $[a, b]$ and:
\[
\int_a^b cf(x)\, dx = c \int_a^b f(x)\, dx
\]
\end{theorem}

\textit{Proof.} For $c > 0$: If $s \leq f \leq t$, then $cs \leq cf \leq ct$. Also $\int_a^b cs = c\int_a^b s$ and $\int_a^b ct = c\int_a^b t$. Taking sup and inf gives the result. For $c < 0$, the inequalities reverse but the conclusion still holds. \qed

\begin{theorem}{Linearity Property}
If $f$ and $g$ are integrable on $[a, b]$, then $f + g$ is integrable on $[a, b]$ and:
\[
\int_a^b [f(x) + g(x)]\, dx = \int_a^b f(x)\, dx + \int_a^b g(x)\, dx
\]
\end{theorem}

\textit{Proof.} Given $\epsilon > 0$, choose step functions $s_1, t_1$ for $f$ and $s_2, t_2$ for $g$ such that:
\[
\int_a^b s_1 > I(f) - \frac{\epsilon}{2}, \quad \int_a^b t_1 < \bar{I}(f) + \frac{\epsilon}{2}
\]
\[
\int_a^b s_2 > I(g) - \frac{\epsilon}{2}, \quad \int_a^b t_2 < \bar{I}(g) + \frac{\epsilon}{2}
\]

Then $s_1 + s_2 \leq f + g \leq t_1 + t_2$, and:
\[
\int_a^b (t_1 + t_2) - \int_a^b (s_1 + s_2) = \left(\int_a^b t_1 - \int_a^b s_1\right) + \left(\int_a^b t_2 - \int_a^b s_2\right) < \epsilon
\]

This shows $f + g$ is integrable. Moreover, taking limits gives the linearity formula. \qed

\begin{theorem}{Comparison Theorem}
If $f$ and $g$ are integrable on $[a, b]$ with $f(x) \leq g(x)$ for all $x \in [a, b]$, then:
\[
\int_a^b f(x)\, dx \leq \int_a^b g(x)\, dx
\]
\end{theorem}

\textit{Proof.} If $s$ is a step function below $f$, then $s$ is also below $g$. Therefore:
\[
\int_a^b s \leq I(g) = \int_a^b g
\]
Taking supremum over all $s$ below $f$ gives $\int_a^b f \leq \int_a^b g$. \qed

\begin{theorem}{}
If $f$ is integrable on $[a, b]$ and $m \leq f(x) \leq M$ for all $x \in [a, b]$, then:
\[
m(b-a) \leq \int_a^b f(x)\, dx \leq M(b-a)
\]
\end{theorem}

\textit{Proof.} Apply the comparison theorem with the constant functions $m$ and $M$. \qed



\chapter{Applications of Integration}

\section{Area Between Two Curves}

\begin{definition}{Area Between Two Curves}
If $f$ and $g$ are integrable on $[a, b]$ with $f(x) \leq g(x)$ for all $x \in [a, b]$, the area $A$ of the region between the graphs of $f$ and $g$ over $[a, b]$ is:
\[
A = \int_a^b [g(x) - f(x)]\, dx
\]
\end{definition}

\section{The Trigonometric Functions}

\begin{theorem}{Integration Formulas for Sine and Cosine}
For any real numbers $a$ and $b$:
\[
\int_a^b \cos x\, dx = \sin b - \sin a
\]
\[
\int_a^b \sin x\, dx = -(\cos b - \cos a) = \cos a - \cos b
\]
\end{theorem}

\textit{Proof.} These formulas are derived from the fundamental properties of the sine and cosine functions and their derivatives. The sine and cosine functions are defined geometrically using the unit circle, and their integral formulas follow from the fact that $\frac{d}{dx}\sin x = \cos x$ and $\frac{d}{dx}\cos x = -\sin x$. \qed

\section{Polar Coordinates}

\begin{definition}{Polar Coordinates}
The \textit{polar coordinates} of a point $P$ in the plane are an ordered pair of numbers $(r, \theta)$, where $r$ is the distance from $P$ to the origin (pole), and $\theta$ is the angle that the ray $OP$ makes with the positive $x$-axis (polar axis).
\end{definition}

\begin{theorem}{Area in Polar Coordinates}
If a function $f$ is nonnegative and integrable on $[a, b]$ where $0 \leq b - a \leq 2\pi$, the area $A$ of the radial set of $f$ over $[a, b]$ is:
\[
A = \frac{1}{2} \int_a^b [f(\theta)]^2\, d\theta
\]
\end{theorem}

\textit{Proof.} The area of a circular sector of radius $r$ and angle $\Delta \theta$ is $\frac{1}{2}r^2 \Delta \theta$. For a general function $r = f(\theta)$, we partition the interval $[a, b]$ and approximate the area by summing the areas of circular sectors. Taking the limit as the partition becomes finer gives the integral formula. \qed

\section{Volume by Integration}

\begin{definition}{Volume by Cross-Sectional Area}
Let $S$ be a solid. A cross-section of $S$ perpendicular to a given axis is the intersection of $S$ with a plane perpendicular to that axis. If the cross-sectional area at position $x$ is $A(x)$, and $A$ is integrable on $[a, b]$, the volume $V$ of $S$ is:
\[
V = \int_a^b A(x)\, dx
\]
\end{definition}

\begin{theorem}{Volume of Solid of Revolution}
If $f$ is integrable on $[a, b]$ and $f(x) \geq 0$ for all $x \in [a, b]$, the volume $V$ of the solid obtained by rotating the ordinate set of $f$ about the $x$-axis is:
\[
V = \pi \int_a^b [f(x)]^2\, dx
\]
\end{theorem}

\textit{Proof.} The cross-section at position $x$ is a circle of radius $f(x)$, so $A(x) = \pi [f(x)]^2$. By the definition of volume by cross-sectional area, $V = \int_a^b A(x)\, dx = \pi \int_a^b [f(x)]^2\, dx$. \qed

\section{Work}

\begin{definition}{Work}
If a force $f(x)$ moves an object from $x = a$ to $x = b$, the \textit{work} $W$ done by the force is:
\[
W = \int_a^b f(x)\, dx
\]
provided $f$ is integrable on $[a, b]$.
\end{definition}

\section{Average Value of a Function}

\begin{definition}{Average Value}
If $f$ is integrable on $[a, b]$, the \textit{average value} (or mean value) of $f$ on $[a, b]$ is:
\[
\bar{f} = \frac{1}{b-a} \int_a^b f(x)\, dx
\]
\end{definition}

\begin{theorem}{Mean-Value Theorem for Integrals}
If $f$ is continuous on $[a, b]$, there exists a point $c \in [a, b]$ such that:
\[
\int_a^b f(x)\, dx = f(c)(b-a)
\]
Equivalently, $f(c) = \bar{f}$.
\end{theorem}

\textit{Proof.} Since $f$ is continuous on $[a, b]$, it attains a minimum $m$ and a maximum $M$ on $[a, b]$. We have:
\[
m(b-a) \leq \int_a^b f(x)\, dx \leq M(b-a)
\]
Dividing by $b-a$: $m \leq \bar{f} \leq M$. By the intermediate value theorem for continuous functions, there exists $c \in [a, b]$ with $f(c) = \bar{f}$. \qed

\section{The Integral as a Function of the Upper Limit}

\begin{definition}{Indefinite Integral}
If $f$ is integrable on $[a, b]$, the function $F$ defined by:
\[
F(x) = \int_a^x f(t)\, dt
\]
for $x \in [a, b]$ is called an \textit{indefinite integral} of $f$.
\end{definition}

\begin{theorem}{}
If $f$ is integrable on $[a, b]$ and $c \in [a, b]$, then:
\[
\int_a^b f(x)\, dx = \int_a^c f(x)\, dx + \int_c^b f(x)\, dx
\]
\end{theorem}

\textit{Proof.} This follows directly from the additivity property of the integral. \qed

\begin{theorem}{}
If $f$ is integrable on $[a, b]$ and $F(x) = \int_a^x f(t)\, dt$, then $F$ is continuous on $[a, b]$.
\end{theorem}

\textit{Proof.} Since $f$ is integrable, it is bounded, say $|f(x)| \leq M$ for all $x \in [a, b]$. For any $x, y \in [a, b]$ with $x < y$:
\[
|F(y) - F(x)| = \left|\int_x^y f(t)\, dt\right| \leq \int_x^y |f(t)|\, dt \leq M(y-x)
\]
Given $\epsilon > 0$, choose $\delta = \epsilon/M$. Then $|y-x| < \delta$ implies $|F(y) - F(x)| < \epsilon$. Thus $F$ is continuous (in fact, uniformly continuous) on $[a, b]$. \qed



\chapter{Continuous Functions}

\section{The Definition of the Limit of a Function}

\begin{definition}{Limit of a Function}
Let $f$ be a function defined on some neighborhood of a point $p$, although we do not assume that $f$ is defined at the point $p$ itself. Let $A$ be a real number. The statement
\[
\lim_{x \to p} f(x) = A
\]
means that for every neighborhood $N_\epsilon(A)$ there is a neighborhood $N_\delta(p)$ such that
\[
f(x) \in N_\epsilon(A) \quad \text{whenever} \quad x \in N_\delta(p) \quad \text{and} \quad x \neq p
\]
In other words, for every $\epsilon > 0$ there is a $\delta > 0$ such that:
\[
|f(x) - A| < \epsilon \quad \text{whenever} \quad 0 < |x - p| < \delta
\]
\end{definition}

\begin{definition}{One-Sided Limits}
If $f$ is defined on some open interval $(p, b)$, the statement
\[
\lim_{x \to p+} f(x) = A
\]
means that for every $\epsilon > 0$ there is a $\delta > 0$ such that $|f(x) - A| < \epsilon$ whenever $p < x < p + \delta$. This is called the \textit{right-hand limit}. Similarly, the \textit{left-hand limit} $\lim_{x \to p-} f(x) = A$ is defined using $p - \delta < x < p$.
\end{definition}

\begin{theorem}{Uniqueness of Limits}
If $\lim_{x \to p} f(x) = A$ and $\lim_{x \to p} f(x) = B$, then $A = B$.
\end{theorem}

\textit{Proof.} Suppose $A \neq B$. Let $\epsilon = |A - B|/2$. By the definition of limit, there exists $\delta_1 > 0$ such that $|f(x) - A| < \epsilon$ when $0 < |x - p| < \delta_1$, and $\delta_2 > 0$ such that $|f(x) - B| < \epsilon$ when $0 < |x - p| < \delta_2$. Let $\delta = \min(\delta_1, \delta_2)$. Then for $0 < |x - p| < \delta$, we have both $|f(x) - A| < \epsilon$ and $|f(x) - B| < \epsilon$. By the triangle inequality, $|A - B| \leq |A - f(x)| + |f(x) - B| < 2\epsilon = |A - B|$, a contradiction. \qed

\section{The Definition of Continuity of a Function}

\begin{definition}{Continuity at a Point}
A function $f$ is said to be \textit{continuous} at a point $p$ if:
\begin{enumerate}
    \item[(a)] $f$ is defined at $p$
    \item[(b)] $\lim_{x \to p} f(x)$ exists
    \item[(c)] $\lim_{x \to p} f(x) = f(p)$
\end{enumerate}
\end{definition}

\begin{definition}{Continuity on an Interval}
A function $f$ is said to be continuous on an interval $I$ if $f$ is continuous at every point of $I$.
\end{definition}

\begin{definition}{Neighborhood}
A \textit{neighborhood} of a point $p$ is an open interval containing $p$. Denoted by $N(p)$ or $N_\delta(p) = (p - \delta, p + \delta)$.
\end{definition}

\section{The Basic Limit Theorems}

\begin{theorem}{Basic Limit Theorems}
Let $f$ and $g$ be functions such that $\lim_{x \to p} f(x) = A$ and $\lim_{x \to p} g(x) = B$. Then:
\begin{enumerate}
    \item[(a)] $\lim_{x \to p} [f(x) + g(x)] = A + B$
    \item[(b)] $\lim_{x \to p} [f(x) - g(x)] = A - B$
    \item[(c)] $\lim_{x \to p} f(x)g(x) = AB$
    \item[(d)] $\lim_{x \to p} f(x)/g(x) = A/B$, provided $B \neq 0$
    \item[(e)] $\lim_{x \to p} cf(x) = cA$ for any constant $c$
\end{enumerate}
\end{theorem}

\textit{Proof.} We prove (a) and (c); the others follow similarly.

(a) Given $\epsilon > 0$. There exists $\delta_1 > 0$ such that $|f(x) - A| < \epsilon/2$ when $0 < |x - p| < \delta_1$, and $\delta_2 > 0$ such that $|g(x) - B| < \epsilon/2$ when $0 < |x - p| < \delta_2$. Let $\delta = \min(\delta_1, \delta_2)$. Then for $0 < |x - p| < \delta$:
\[
|(f(x) + g(x)) - (A + B)| \leq |f(x) - A| + |g(x) - B| < \frac{\epsilon}{2} + \frac{\epsilon}{2} = \epsilon
\]

(c) Given $\epsilon > 0$. Choose $\delta_1 > 0$ such that $|f(x) - A| < \epsilon/(2|B| + 1)$ when $0 < |x - p| < \delta_1$. Choose $\delta_2 > 0$ such that $|g(x) - B| < \epsilon/(2|A| + 1)$ when $0 < |x - p| < \delta_2$. Choose $\delta_3 > 0$ such that $|g(x) - B| < 1$ (so $|g(x)| < |B| + 1$) when $0 < |x - p| < \delta_3$. Let $\delta = \min(\delta_1, \delta_2, \delta_3)$. Then:
\[
|f(x)g(x) - AB| = |f(x)g(x) - Ag(x) + Ag(x) - AB| \leq |g(x)||f(x) - A| + |A||g(x) - B| < (|B|+1)\frac{\epsilon}{2|B|+1} + |A|\frac{\epsilon}{2|A|+1} < \epsilon
\] \qed

\begin{theorem}{Comparison Theorem for Limits}
Suppose $f(x) \leq g(x) \leq h(x)$ for all $x$ in some neighborhood of $p$ (except possibly at $p$ itself). If $\lim_{x \to p} f(x) = \lim_{x \to p} h(x) = L$, then $\lim_{x \to p} g(x) = L$.
\end{theorem}

\textit{Proof.} Given $\epsilon > 0$. There exists $\delta_1 > 0$ such that $L - \epsilon < f(x) < L + \epsilon$ when $0 < |x - p| < \delta_1$, and $\delta_2 > 0$ such that $L - \epsilon < h(x) < L + \epsilon$ when $0 < |x - p| < \delta_2$. Let $\delta = \min(\delta_1, \delta_2)$. Then for $0 < |x - p| < \delta$:
\[
L - \epsilon < f(x) \leq g(x) \leq h(x) < L + \epsilon
\]
Thus $|g(x) - L| < \epsilon$. \qed

\section{Composite Functions and Continuity}

\begin{definition}{Composite Function}
If $u$ and $v$ are functions, the \textit{composition} of $u$ and $v$, denoted by $u \circ v$, is the function defined by:
\[
(u \circ v)(x) = u(v(x))
\]
The domain of $u \circ v$ consists of all $x$ in the domain of $v$ such that $v(x)$ is in the domain of $u$.
\end{definition}

\begin{theorem}{Continuity of Composite Functions}
If $v$ is continuous at $p$ and $u$ is continuous at $v(p)$, then the composite function $f = u \circ v$ is continuous at $p$.
\end{theorem}

\textit{Proof.} Given $\epsilon > 0$. Since $u$ is continuous at $v(p)$, there exists $\delta_1 > 0$ such that $|u(y) - u(v(p))| < \epsilon$ whenever $|y - v(p)| < \delta_1$. Since $v$ is continuous at $p$, there exists $\delta > 0$ such that $|v(x) - v(p)| < \delta_1$ whenever $|x - p| < \delta$. Therefore, $|f(x) - f(p)| = |u(v(x)) - u(v(p))| < \epsilon$ whenever $|x - p| < \delta$. \qed

\section{Bolzano's Theorem for Continuous Functions}

\begin{theorem}{Bolzano's Theorem}
Let $f$ be continuous at each point of a closed interval $[a, b]$ and suppose $f(a)$ and $f(b)$ have opposite signs. Then there is at least one $c$ in the open interval $(a, b)$ such that $f(c) = 0$.
\end{theorem}

\textit{Proof.} Assume $f(a) < 0$ and $f(b) > 0$ (the other case is similar). Let $S = \{x \in [a, b] \mid f(x) \leq 0\}$. Then $S$ is nonempty (contains $a$) and bounded above by $b$. By the completeness axiom, $S$ has a supremum, call it $c$. We claim $f(c) = 0$.

Suppose $f(c) > 0$. By continuity, there exists $\delta > 0$ such that $f(x) > 0$ for all $x \in (c - \delta, c + \delta) \cap [a, b]$. Then $c - \delta$ is an upper bound for $S$, contradicting that $c$ is the least upper bound.

Suppose $f(c) < 0$. By continuity, there exists $\delta > 0$ such that $f(x) < 0$ for all $x \in (c - \delta, c + \delta) \cap [a, b]$. Then there are points in $S$ greater than $c$, contradicting that $c$ is an upper bound.

Therefore $f(c) = 0$. \qed

\begin{theorem}{Intermediate-Value Theorem for Continuous Functions}
Let $f$ be continuous at each point of a closed interval $[a, b]$. Then $f$ takes on every value between $f(a)$ and $f(b)$ somewhere in the interval.
\end{theorem}

\textit{Proof.} Let $k$ be any value between $f(a)$ and $f(b)$. Define $g(x) = f(x) - k$. Then $g$ is continuous on $[a, b]$ and $g(a)$ and $g(b)$ have opposite signs. By Bolzano's theorem, there exists $c \in (a, b)$ such that $g(c) = 0$, i.e., $f(c) = k$. \qed

\section{The Process of Inversion}

\begin{definition}{Strictly Monotonic Function}
A function $f$ is said to be \textit{strictly increasing} on an interval $I$ if $f(x) < f(y)$ whenever $x < y$ in $I$. It is \textit{strictly decreasing} if $f(x) > f(y)$ whenever $x < y$ in $I$. A function is \textit{strictly monotonic} if it is either strictly increasing or strictly decreasing.
\end{definition}

\begin{theorem}{Existence of Inverse Functions}
Let $f$ be strictly increasing and continuous on an interval $[a, b]$. Let $c = f(a)$ and $d = f(b)$. Then:
\begin{enumerate}
    \item[(a)] For every $y$ in $[c, d]$ there is exactly one $x$ in $[a, b]$ such that $y = f(x)$.
    \item[(b)] The inverse function $g$ defined by $g(y) = x$ is strictly increasing and continuous on $[c, d]$.
\end{enumerate}
\end{theorem}

\textit{Proof.} (a) Existence follows from the intermediate-value theorem. Uniqueness follows from strict monotonicity: if $f(x_1) = f(x_2) = y$ with $x_1 \neq x_2$, then $f$ is not strictly monotonic.

(b) To show $g$ is strictly increasing: if $y_1 < y_2$, let $x_1 = g(y_1)$ and $x_2 = g(y_2)$. If $x_1 \geq x_2$, then $f(x_1) \geq f(x_2)$ (since $f$ is increasing), so $y_1 \geq y_2$, contradiction.

To show $g$ is continuous: Given $y_0 \in [c, d]$ and $\epsilon > 0$. Let $x_0 = g(y_0)$. Choose $\delta'$ such that $[x_0 - \delta', x_0 + \delta'] \subseteq [a, b]$. Since $f$ is continuous and strictly increasing, $f$ maps $[x_0 - \delta', x_0 + \delta']$ onto $[f(x_0 - \delta'), f(x_0 + \delta')]$. Let $\delta = \min(y_0 - f(x_0 - \delta'), f(x_0 + \delta') - y_0)$. Then $|y - y_0| < \delta$ implies $|g(y) - g(y_0)| < \delta'$. \qed

\section{The Extreme-Value Theorem for Continuous Functions}

\begin{theorem}{Extreme-Value Theorem}
If $f$ is continuous on a closed interval $[a, b]$, then there exist points $c$ and $d$ in $[a, b]$ such that:
\[
f(c) \leq f(x) \leq f(d) \quad \text{for all } x \in [a, b]
\]
\end{theorem}

\textit{Proof.} We prove the existence of the maximum; the minimum is similar.

First, $f$ is bounded above on $[a, b]$. If not, for each $n$ there exists $x_n \in [a, b]$ with $f(x_n) > n$. The sequence $\{x_n\}$ has a convergent subsequence by Bolzano-Weierstrass, converging to some $x \in [a, b]$. By continuity, $f(x_n) \to f(x)$, contradicting $f(x_n) > n$.

Let $M = \sup\{f(x) \mid x \in [a, b]\}$. For each $n$, choose $x_n$ such that $f(x_n) > M - 1/n$. The sequence $\{x_n\}$ has a convergent subsequence $\{x_{n_k}\}$ converging to some $d \in [a, b]$. By continuity, $f(x_{n_k}) \to f(d)$. Since $M - 1/n_k < f(x_{n_k}) \leq M$, we have $f(d) = M$. \qed

\section{The Small-Span Theorem for Continuous Functions (Uniform Continuity)}

\begin{definition}{Uniform Continuity}
A function $f$ is said to be \textit{uniformly continuous} on an interval $I$ if for every $\epsilon > 0$ there exists $\delta > 0$ such that $|f(x) - f(y)| < \epsilon$ whenever $x, y \in I$ and $|x - y| < \delta$.
\end{definition}

\begin{theorem}{Small-Span Theorem (Uniform Continuity Theorem)}
If $f$ is continuous on a closed interval $[a, b]$, then $f$ is uniformly continuous on $[a, b]$.
\end{theorem}

\textit{Proof.} Suppose $f$ is not uniformly continuous on $[a, b]$. Then there exists $\epsilon > 0$ such that for every $\delta > 0$, there exist $x, y \in [a, b]$ with $|x - y| < \delta$ but $|f(x) - f(y)| \geq \epsilon$.

Taking $\delta = 1/n$, we obtain sequences $\{x_n\}$ and $\{y_n\}$ in $[a, b]$ such that $|x_n - y_n| < 1/n$ but $|f(x_n) - f(y_n)| \geq \epsilon$. By Bolzano-Weierstrass, $\{x_n\}$ has a convergent subsequence $\{x_{n_k}\}$ converging to some $p \in [a, b]$. Since $|x_{n_k} - y_{n_k}| < 1/n_k$, we also have $y_{n_k} \to p$. By continuity of $f$:
\[
\lim_{k \to \infty} f(x_{n_k}) = f(p) = \lim_{k \to \infty} f(y_{n_k})
\]
This contradicts $|f(x_{n_k}) - f(y_{n_k})| \geq \epsilon$ for all $k$. \qed

\section{The Integrability Theorem for Continuous Functions}

\begin{theorem}{Integrability of Continuous Functions}
If $f$ is continuous on $[a, b]$, then $f$ is integrable on $[a, b]$.
\end{theorem}

\textit{Proof.} Since $f$ is continuous on the closed interval $[a, b]$, it is bounded and uniformly continuous.

Given $\epsilon > 0$, by uniform continuity there exists $\delta > 0$ such that $|f(x) - f(y)| < \epsilon/(b-a)$ whenever $|x - y| < \delta$. Choose a partition $P = \{x_0, x_1, \ldots, x_n\}$ with $\Delta x_k < \delta$ for all $k$.

On each subinterval $[x_{k-1}, x_k]$, let $m_k = \inf f$ and $M_k = \sup f$. By continuity, $f$ attains these values, so there exist $c_k, d_k \in [x_{k-1}, x_k]$ with $f(c_k) = m_k$ and $f(d_k) = M_k$. Since $|c_k - d_k| \leq \Delta x_k < \delta$, we have $M_k - m_k = |f(d_k) - f(c_k)| < \epsilon/(b-a)$.

Therefore:
\[
\sum_{k=1}^{n} (M_k - m_k)\Delta x_k < \frac{\epsilon}{b-a}\sum_{k=1}^{n} \Delta x_k = \epsilon
\]
This shows the upper and lower integrals differ by less than $\epsilon$, so $f$ is integrable. \qed

\section{Mean-Value Theorems for Integrals of Continuous Functions}

\begin{theorem}{First Mean-Value Theorem for Integrals}
If $f$ is continuous on $[a, b]$, then for some $c \in [a, b]$:
\[
\int_a^b f(x)\, dx = f(c)(b-a)
\]
\end{theorem}

\textit{Proof.} By the extreme-value theorem, $f$ attains minimum $m$ and maximum $M$ on $[a, b]$. Then:
\[
m(b-a) \leq \int_a^b f(x)\, dx \leq M(b-a)
\]
So $\frac{1}{b-a}\int_a^b f(x)\, dx$ is a value between $m$ and $M$. By the intermediate-value theorem, there exists $c \in [a, b]$ such that $f(c) = \frac{1}{b-a}\int_a^b f(x)\, dx$. \qed

\begin{theorem}{Weighted Mean-Value Theorem for Integrals}
If $f$ and $g$ are continuous on $[a, b]$, and if $g$ never changes sign on $[a, b]$, then for some $c \in [a, b]$:
\[
\int_a^b f(x)g(x)\, dx = f(c)\int_a^b g(x)\, dx
\]
\end{theorem}

\textit{Proof.} Assume $g(x) \geq 0$ on $[a, b]$. Let $m = \min f$ and $M = \max f$ on $[a, b]$. Then:
\[
mg(x) \leq f(x)g(x) \leq Mg(x) \quad \text{for all } x \in [a, b]
\]
Integrating: $m\int_a^b g \leq \int_a^b fg \leq M\int_a^b g$. If $\int_a^b g = 0$, the result is trivial. Otherwise, $\frac{\int_a^b fg}{\int_a^b g}$ is a value between $m$ and $M$, so by the intermediate-value theorem, it equals $f(c)$ for some $c \in [a, b]$. \qed



% ====================
% Part III: Differential Calculus
% ====================
\part{Differential Calculus}

\chapter{Differential Calculus}

\section{The Derivative of a Function}

\begin{definition}{Derivative}
Let $f$ be a function defined at least on some open interval containing the point $p$. The \textit{derivative} of $f$ at $p$, denoted by $f'(p)$, is defined by the limit:
\[
f'(p) = \lim_{h \to 0} \frac{f(p+h) - f(p)}{h}
\]
whenever this limit exists. The process of forming this limit is called \textit{differentiation}.
\end{definition}

\begin{definition}{Differentiability}
A function $f$ is said to be \textit{differentiable} at a point $p$ if the derivative $f'(p)$ exists. If $f$ is differentiable at every point of an interval $I$, we say $f$ is differentiable on $I$.
\end{definition}

\begin{theorem}{Differentiability Implies Continuity}
If $f$ is differentiable at $p$, then $f$ is continuous at $p$.
\end{theorem}

\textit{Proof.} We need to show $\lim_{x \to p} f(x) = f(p)$. For $x \neq p$:
\[
f(x) - f(p) = \frac{f(x) - f(p)}{x - p} \cdot (x - p)
\]
Taking limits as $x \to p$:
\[
\lim_{x \to p} [f(x) - f(p)] = f'(p) \cdot 0 = 0
\]
Thus $\lim_{x \to p} f(x) = f(p)$. \qed

\section{The Algebra of Derivatives}

\begin{theorem}{Sum Rule}
If $f$ and $g$ are differentiable at $p$, then $f + g$ is differentiable at $p$ and:
\[
(f + g)'(p) = f'(p) + g'(p)
\]
\end{theorem}

\textit{Proof.} Using the limit definition:
\[
\lim_{h \to 0} \frac{(f+g)(p+h) - (f+g)(p)}{h} = \lim_{h \to 0} \left[\frac{f(p+h) - f(p)}{h} + \frac{g(p+h) - g(p)}{h}\right] = f'(p) + g'(p)
\] \qed

\begin{theorem}{Constant Multiple Rule}
If $f$ is differentiable at $p$ and $c$ is a constant, then $cf$ is differentiable at $p$ and:
\[
(cf)'(p) = c \cdot f'(p)
\]
\end{theorem}

\textit{Proof.} $\lim_{h \to 0} \frac{cf(p+h) - cf(p)}{h} = c \cdot \lim_{h \to 0} \frac{f(p+h) - f(p)}{h} = c \cdot f'(p)$. \qed

\begin{theorem}{Product Rule}
If $f$ and $g$ are differentiable at $p$, then $fg$ is differentiable at $p$ and:
\[
(fg)'(p) = f'(p)g(p) + f(p)g'(p)
\]
\end{theorem}

\textit{Proof.} We add and subtract $f(p)g(p+h)$:
\begin{align*}
\frac{(fg)(p+h) - (fg)(p)}{h} &= \frac{f(p+h)g(p+h) - f(p)g(p)}{h} \\
&= \frac{f(p+h)g(p+h) - f(p)g(p+h) + f(p)g(p+h) - f(p)g(p)}{h} \\
&= g(p+h)\frac{f(p+h) - f(p)}{h} + f(p)\frac{g(p+h) - g(p)}{h}
\end{align*}
As $h \to 0$: $g(p+h) \to g(p)$ by continuity, so the limit is $g(p)f'(p) + f(p)g'(p)$. \qed

\begin{theorem}{Quotient Rule}
If $f$ and $g$ are differentiable at $p$ and $g(p) \neq 0$, then $f/g$ is differentiable at $p$ and:
\[
\left(\frac{f}{g}\right)'(p) = \frac{f'(p)g(p) - f(p)g'(p)}{[g(p)]^2}
\]
\end{theorem}

\textit{Proof.} Using the product rule on $f \cdot (1/g)$, we first show $\left(\frac{1}{g}\right)'(p) = -\frac{g'(p)}{[g(p)]^2}$:
\[
\lim_{h \to 0} \frac{\frac{1}{g(p+h)} - \frac{1}{g(p)}}{h} = \lim_{h \to 0} \frac{g(p) - g(p+h)}{h \cdot g(p+h)g(p)} = -\frac{g'(p)}{[g(p)]^2}
\]
Then apply the product rule to $f \cdot (1/g)$. \qed

\section{Derivatives of Standard Functions}

\begin{theorem}{Derivative of a Constant}
If $f(x) = c$ for all $x$, then $f'(x) = 0$ for all $x$.
\end{theorem}

\textit{Proof.} $\frac{f(p+h) - f(p)}{h} = \frac{c - c}{h} = 0$, so $f'(p) = \lim_{h \to 0} 0 = 0$. \qed

\begin{theorem}{Derivative of the Identity Function}
If $f(x) = x$ for all $x$, then $f'(x) = 1$ for all $x$.
\end{theorem}

\textit{Proof.} $\frac{f(p+h) - f(p)}{h} = \frac{(p+h) - p}{h} = \frac{h}{h} = 1$, so $f'(p) = 1$. \qed

\begin{theorem}{Power Rule for Positive Integers}
If $n$ is a positive integer and $f(x) = x^n$, then $f'(x) = nx^{n-1}$.
\end{theorem}

\textit{Proof.} By the binomial theorem:
\[
\frac{(p+h)^n - p^n}{h} = \frac{np^{n-1}h + \binom{n}{2}p^{n-2}h^2 + \cdots + h^n}{h} = np^{n-1} + \binom{n}{2}p^{n-2}h + \cdots
\]
As $h \to 0$, all terms except the first approach 0, giving $np^{n-1}$. \qed

\begin{theorem}{Derivatives of Sine and Cosine}
\[
\frac{d}{dx}(\sin x) = \cos x \quad \text{and} \quad \frac{d}{dx}(\cos x) = -\sin x
\]
\end{theorem}

\textit{Proof.} Using the trigonometric identity $\sin(p+h) = \sin p \cos h + \cos p \sin h$:
\[
\frac{\sin(p+h) - \sin p}{h} = \sin p \cdot \frac{\cos h - 1}{h} + \cos p \cdot \frac{\sin h}{h}
\]
Using the limits $\lim_{h \to 0} \frac{\sin h}{h} = 1$ and $\lim_{h \to 0} \frac{\cos h - 1}{h} = 0$, we get $\cos p$. Similarly for cosine. \qed

\section{The Chain Rule}

\begin{theorem}{Chain Rule}
If $f$ is differentiable at $p$ and $g$ is differentiable at $f(p)$, then the composite function $g \circ f$ is differentiable at $p$ and:
\[
(g \circ f)'(p) = g'(f(p)) \cdot f'(p)
\]
\end{theorem}

\textit{Proof.} Define a function $d$ by:
\[
d(y) = \begin{cases} \frac{g(y) - g(f(p))}{y - f(p)} & \text{if } y \neq f(p) \\ g'(f(p)) & \text{if } y = f(p) \end{cases}
\]
Then $d$ is continuous at $f(p)$ and $g(y) - g(f(p)) = d(y)(y - f(p))$ for all $y$. Setting $y = f(p+h)$:
\[
\frac{g(f(p+h)) - g(f(p))}{h} = d(f(p+h)) \cdot \frac{f(p+h) - f(p)}{h}
\]
As $h \to 0$: $f(p+h) \to f(p)$ by continuity of $f$, so $d(f(p+h)) \to d(f(p)) = g'(f(p))$. The second factor approaches $f'(p)$. Their product is $g'(f(p)) \cdot f'(p)$. \qed

\section{The Mean-Value Theorem for Derivatives}

\begin{theorem}{Rolle's Theorem}
If $f$ is continuous on $[a, b]$, differentiable on $(a, b)$, and $f(a) = f(b)$, then there exists $c \in (a, b)$ such that $f'(c) = 0$.
\end{theorem}

\textit{Proof.} By the extreme-value theorem, $f$ attains maximum $M$ and minimum $m$ on $[a, b]$. If both occur at endpoints, then $f$ is constant and $f'(c) = 0$ for all $c \in (a, b)$. Otherwise, $f$ has an interior extremum at some $c \in (a, b)$. We show $f'(c) = 0$.

For $h > 0$ small enough that $c + h \in (a, b)$:
- If $f(c)$ is a maximum: $\frac{f(c+h) - f(c)}{h} \leq 0$, so $f'(c) = \lim_{h \to 0+} \leq 0$
- Also $\frac{f(c-h) - f(c)}{-h} \geq 0$, so $f'(c) = \lim_{h \to 0+} \geq 0$

Thus $f'(c) = 0$. \qed

\begin{theorem}{Mean-Value Theorem for Derivatives}
If $f$ is continuous on $[a, b]$ and differentiable on $(a, b)$, then there exists $c \in (a, b)$ such that:
\[
f'(c) = \frac{f(b) - f(a)}{b - a}
\]
\end{theorem}

\textit{Proof.} Define $g(x) = f(x) - \frac{f(b) - f(a)}{b - a}(x - a)$. Then $g(a) = f(a)$ and $g(b) = f(b) - (f(b) - f(a)) = f(a) = g(a)$. By Rolle's theorem, there exists $c \in (a, b)$ with $g'(c) = 0$. Since $g'(x) = f'(x) - \frac{f(b) - f(a)}{b - a}$, we have $f'(c) = \frac{f(b) - f(a)}{b - a}$. \qed

\begin{theorem}{Cauchy's Mean-Value Theorem}
If $f$ and $g$ are continuous on $[a, b]$ and differentiable on $(a, b)$, then there exists $c \in (a, b)$ such that:
\[
f'(c)[g(b) - g(a)] = g'(c)[f(b) - f(a)]
\]
If $g'(c) \neq 0$ and $g(b) \neq g(a)$, this can be written as:
\[
\frac{f'(c)}{g'(c)} = \frac{f(b) - f(a)}{g(b) - g(a)}
\]
\end{theorem}

\textit{Proof.} Define $h(x) = f(x)[g(b) - g(a)] - g(x)[f(b) - f(a)]$. Then $h(a) = h(b)$. By Rolle's theorem, there exists $c \in (a, b)$ with $h'(c) = 0$, which gives the result. \qed

\section{Applications of Differentiation to Extreme Values}

\begin{definition}{Absolute and Relative Extrema}
A function $f$ has an \textit{absolute maximum} at $p$ if $f(p) \geq f(x)$ for all $x$ in the domain of $f$. It has a \textit{relative maximum} at $p$ if $f(p) \geq f(x)$ for all $x$ in some neighborhood of $p$. Similar definitions apply for minima.
\end{definition}

\begin{definition}{Critical Point}
A point $p$ in the domain of $f$ is called a \textit{critical point} if either $f'(p) = 0$ or $f$ is not differentiable at $p$.
\end{definition}

\begin{theorem}{}
If $f$ has a relative extremum at an interior point $p$ and $f$ is differentiable at $p$, then $f'(p) = 0$.
\end{theorem}

\textit{Proof.} If $f$ has a relative maximum at $p$, then for $h$ sufficiently small:
- For $h > 0$: $\frac{f(p+h) - f(p)}{h} \leq 0$, so $f'(p) \leq 0$
- For $h < 0$: $\frac{f(p+h) - f(p)}{h} \geq 0$, so $f'(p) \geq 0$
Thus $f'(p) = 0$. \qed

\section{Geometric Properties of Functions}

\begin{theorem}{}
If $f$ is continuous on $[a, b]$ and $f'(x) = 0$ for all $x \in (a, b)$, then $f$ is constant on $[a, b]$.
\end{theorem}

\textit{Proof.} For any $x_1, x_2 \in [a, b]$ with $x_1 < x_2$, by the mean-value theorem there exists $c \in (x_1, x_2)$ with $f(x_2) - f(x_1) = f'(c)(x_2 - x_1) = 0$. Thus $f(x_1) = f(x_2)$. \qed

\begin{theorem}{Increasing Function Theorem}
If $f$ is continuous on $[a, b]$ and $f'(x) > 0$ for all $x \in (a, b)$, then $f$ is strictly increasing on $[a, b]$.
\end{theorem}

\textit{Proof.} For $a \leq x_1 < x_2 \leq b$, by the mean-value theorem there exists $c \in (x_1, x_2)$ with $f(x_2) - f(x_1) = f'(c)(x_2 - x_1) > 0$. Thus $f(x_2) > f(x_1)$. \qed

\begin{theorem}{Concavity and the Second Derivative}
If $f''(x) > 0$ for all $x \in (a, b)$, then $f$ is concave upward on $(a, b)$. If $f''(x) < 0$ for all $x \in (a, b)$, then $f$ is concave downward on $(a, b)$.
\end{theorem}

\textit{Proof.} If $f'' > 0$, then $f'$ is increasing. For any $c \in (a, b)$, the tangent line at $c$ has slope $f'(c)$. Since $f'$ is increasing, $f'(x) < f'(c)$ for $x < c$ and $f'(x) > f'(c)$ for $x > c$. This means the graph lies above its tangent lines, i.e., $f$ is concave upward. \qed

\begin{theorem}{Second-Derivative Test for Extrema}
If $f'(p) = 0$ and $f''(p) > 0$, then $f$ has a relative minimum at $p$. If $f'(p) = 0$ and $f''(p) < 0$, then $f$ has a relative maximum at $p$.
\end{theorem}

\textit{Proof.} Suppose $f'(p) = 0$ and $f''(p) > 0$. Since $f''(p) = \lim_{h \to 0} \frac{f'(p+h) - f'(p)}{h} = \lim_{h \to 0} \frac{f'(p+h)}{h} > 0$, we have $f'(p+h)/h > 0$ for small $h$. This means $f'(p+h) < 0$ for $h < 0$ and $f'(p+h) > 0$ for $h > 0$. So $f$ decreases to the left of $p$ and increases to the right, giving a minimum. \qed

\section{Partial Derivatives}

\begin{definition}{Partial Derivative}
Let $f$ be a function of two variables. The \textit{partial derivative} of $f$ with respect to $x$ at $(a, b)$ is:
\[
\frac{\partial f}{\partial x}(a, b) = \lim_{h \to 0} \frac{f(a+h, b) - f(a, b)}{h}
\]
whenever this limit exists. Similarly, the partial derivative with respect to $y$ is:
\[
\frac{\partial f}{\partial y}(a, b) = \lim_{h \to 0} \frac{f(a, b+h) - f(a, b)}{h}
\]
\end{definition}

\begin{theorem}{Equality of Mixed Partial Derivatives}
If $f_{xy}$ and $f_{yx}$ are both continuous at a point $(a, b)$, then:
\[
f_{xy}(a, b) = f_{yx}(a, b)
\]
\end{theorem}

\textit{Proof.} For small $h, k > 0$, define $\Delta = f(a+h, b+k) - f(a+h, b) - f(a, b+k) + f(a, b)$. Let $g(x) = f(x, b+k) - f(x, b)$. Then $\Delta = g(a+h) - g(a)$. By the mean-value theorem, $\Delta = g'(c)h = [f_x(c, b+k) - f_x(c, b)]h$ for some $c \in (a, a+h)$. Applying the mean-value theorem again to $f_x$ in the second variable: $\Delta = f_{xy}(c, d)hk$ for some $d \in (b, b+k)$.

Similarly, defining $\tilde{g}(y) = f(a+h, y) - f(a, y)$, we get $\Delta = f_{yx}(\tilde{c}, \tilde{d})kh$. Thus $f_{xy}(c, d) = f_{yx}(\tilde{c}, \tilde{d})$. As $h, k \to 0$, continuity gives $f_{xy}(a, b) = f_{yx}(a, b)$. \qed



\chapter{The Relation Between Integration and Differentiation}

\section{The First Fundamental Theorem of Calculus}

\begin{theorem}{First Fundamental Theorem of Calculus}
Let $f$ be a function that is integrable on $[a, b]$ and continuous at a point $c$ in $(a, b)$. Define $A(x) = \int_a^x f(t)\, dt$ for $x \in [a, b]$. Then $A$ is differentiable at $c$ and:
\[
A'(c) = f(c)
\]
\end{theorem}

\textit{Proof.} We need to show that $\lim_{h \to 0} \frac{A(c+h) - A(c)}{h} = f(c)$. For $h > 0$:
\[
\frac{A(c+h) - A(c)}{h} = \frac{1}{h} \int_c^{c+h} f(t)\, dt
\]

By the mean-value theorem for integrals, there exists $c_h \in [c, c+h]$ such that:
\[
\int_c^{c+h} f(t)\, dt = f(c_h) \cdot h
\]
Therefore $\frac{A(c+h) - A(c)}{h} = f(c_h)$. As $h \to 0^+$, we have $c_h \to c$, and by continuity of $f$ at $c$, $f(c_h) \to f(c)$. The same argument applies for $h < 0$. Thus $A'(c) = f(c)$. \qed

\section{The Zero-Derivative Theorem}

\begin{theorem}{Zero-Derivative Theorem}
If $f$ is differentiable on an open interval $I$ and $f'(x) = 0$ for all $x \in I$, then $f$ is constant on $I$.
\end{theorem}

\textit{Proof.} Let $a, b \in I$ with $a < b$. By the mean-value theorem, there exists $c \in (a, b)$ such that $f(b) - f(a) = f'(c)(b - a) = 0$. Thus $f(a) = f(b)$ for all $a, b \in I$, so $f$ is constant. \qed

\section{Primitive Functions and the Second Fundamental Theorem}

\begin{definition}{Primitive Function}
A function $P$ is called a \textit{primitive} (or \textit{antiderivative}) of $f$ on an interval $I$ if $P$ is differentiable on $I$ and $P'(x) = f(x)$ for all $x \in I$.
\end{definition}

\begin{theorem}{}
If $P$ is a primitive of $f$ on an interval $I$, then every primitive of $f$ on $I$ has the form $P + C$ for some constant $C$.
\end{theorem}

\textit{Proof.} If $Q$ is another primitive of $f$, then $(Q - P)' = Q' - P' = f - f = 0$. By the zero-derivative theorem, $Q - P = C$ for some constant $C$. \qed

\begin{theorem}{Second Fundamental Theorem of Calculus}
Let $f$ be continuous on $[a, b]$ and let $P$ be any primitive of $f$ on $[a, b]$. Then:
\[
\int_a^b f(x)\, dx = P(b) - P(a)
\]
\end{theorem}

\textit{Proof.} Define $A(x) = \int_a^x f(t)\, dt$. By the first fundamental theorem, $A'(x) = f(x)$ for all $x \in (a, b)$. Thus $A$ is a primitive of $f$. Since $P$ is also a primitive, $P(x) = A(x) + C$ for some constant $C$. Then:
\[
P(b) - P(a) = [A(b) + C] - [A(a) + C] = A(b) - A(a) = \int_a^b f(t)\, dt - 0 = \int_a^b f(x)\, dx
\] \qed

\section{Integration Formulas}

\begin{theorem}{Integration of Rational Powers}
For any rational number $r \neq -1$ and any interval where $x^r$ is defined:
\[
\int x^r\, dx = \frac{x^{r+1}}{r+1} + C
\]
\end{theorem}

\textit{Proof.} By differentiation: $\frac{d}{dx}\left(\frac{x^{r+1}}{r+1}\right) = x^r$. \qed

\begin{theorem}{Integration of Trigonometric Functions}
\[
\int \sin x\, dx = -\cos x + C \quad \text{and} \quad \int \cos x\, dx = \sin x + C
\]
\end{theorem}

\textit{Proof.} By differentiation: $\frac{d}{dx}(-\cos x) = \sin x$ and $\frac{d}{dx}(\sin x) = \cos x$. \qed

\section{Integration by Substitution}

\begin{theorem}{Integration by Substitution}
Let $g$ have a continuous derivative $g'$ on an open interval $I$. Let $J = g(I)$ and assume $f$ is continuous on $J$. Then for any $a, b \in I$:
\[
\int_a^b f[g(x)]g'(x)\, dx = \int_{g(a)}^{g(b)} f(u)\, du
\]
\end{theorem}

\textit{Proof.} Let $P$ be a primitive of $f$ on $J$, so $P'(u) = f(u)$. Define $Q(x) = P(g(x))$. By the chain rule:
\[
Q'(x) = P'(g(x))g'(x) = f(g(x))g'(x)
\]
So $Q$ is a primitive of $f(g(x))g'(x)$. By the second fundamental theorem:
\[
\int_a^b f[g(x)]g'(x)\, dx = Q(b) - Q(a) = P(g(b)) - P(g(a)) = \int_{g(a)}^{g(b)} f(u)\, du
\] \qed

\section{Integration by Parts}

\begin{theorem}{Integration by Parts}
If $u$ and $v$ have continuous derivatives on an interval $[a, b]$, then:
\[
\int_a^b u(x)v'(x)\, dx = u(b)v(b) - u(a)v(a) - \int_a^b v(x)u'(x)\, dx
\]
In differential form, with $du = u'(x)dx$ and $dv = v'(x)dx$:
\[
\int u\, dv = uv - \int v\, du
\]
\end{theorem}

\textit{Proof.} From the product rule: $(uv)' = u'v + uv'$. Integrating both sides from $a$ to $b$:
\[
u(b)v(b) - u(a)v(a) = \int_a^b u'(x)v(x)\, dx + \int_a^b u(x)v'(x)\, dx
\]
Rearranging gives the integration by parts formula. \qed

\begin{theorem}{Second Mean-Value Theorem for Integrals}
Assume $f$ is continuous on $[a, b]$ and $g$ has a continuous derivative that never changes sign on $[a, b]$. Then for some $c \in [a, b]$:
\[
\int_a^b f(x)g(x)\, dx = f(a)\int_a^c g(x)\, dx + f(b)\int_c^b g(x)\, dx
\]
\end{theorem}

\textit{Proof.} Let $G(x) = \int_a^x g(t)\, dt$. Integration by parts gives:
\[
\int_a^b f(x)g(x)\, dx = f(b)G(b) - \int_a^b f'(x)G(x)\, dx
\]
By the weighted mean-value theorem, there exists $c \in [a, b]$ such that:
\[
\int_a^b f'(x)G(x)\, dx = G(c)\int_a^b f'(x)\, dx = G(c)[f(b) - f(a)]
\]
Substituting back:
\[
\int_a^b f(x)g(x)\, dx = f(b)G(b) - G(c)[f(b) - f(a)] = f(a)G(c) + f(b)[G(b) - G(c)]
\]
Since $G(c) = \int_a^c g$ and $G(b) - G(c) = \int_c^b g$, the result follows. \qed



\chapter{Transcendental Functions}

\section{The Natural Logarithm}

\begin{definition}{Natural Logarithm}
The \textit{natural logarithm} of $x$, denoted by $\log x$ or $\ln x$, is defined for all $x > 0$ by the integral:
\[
\log x = \int_1^x \frac{1}{t}\, dt
\]
\end{definition}

\begin{theorem}{Basic Properties of the Logarithm}
\begin{enumerate}
    \item[(a)] $\log 1 = 0$
    \item[(b)] $\log x > 0$ if $x > 1$, and $\log x < 0$ if $0 < x < 1$
    \item[(c)] $\log(ab) = \log a + \log b$ for all $a, b > 0$ \textit{(functional equation)}
    \item[(d)] $\log(a/b) = \log a - \log b$ for all $a, b > 0$
    \item[(e)] $\log(a^r) = r\log a$ for all $a > 0$ and rational $r$
\end{enumerate}
\end{theorem}

\textit{Proof.} (a) $\log 1 = \int_1^1 \frac{1}{t}\, dt = 0$.

(c) For fixed $a > 0$, define $f(x) = \log(ax) - \log a - \log x$. Then:
\[
f'(x) = \frac{1}{ax} \cdot a - \frac{1}{x} = \frac{1}{x} - \frac{1}{x} = 0
\]
So $f(x) = C$ (constant). Since $f(1) = \log a - \log a - 0 = 0$, we have $C = 0$. Thus $\log(ax) = \log a + \log x$. \qed

\begin{theorem}{Derivative and Integral of the Logarithm}
\[
\frac{d}{dx} \log x = \frac{1}{x} \quad \text{for all } x > 0
\]
\[
\int \frac{1}{x}\, dx = \log |x| + C \quad \text{for } x \neq 0
\]
\end{theorem}

\textit{Proof.} By the first fundamental theorem of calculus, since $\log x = \int_1^x \frac{1}{t}\, dt$, we have $\frac{d}{dx}\log x = \frac{1}{x}$. \qed

\section{The Exponential Function}

\begin{definition}{Exponential Function}
Since $\log x$ is strictly increasing and continuous on $(0, \infty)$ with range $(-\infty, \infty)$, it has an inverse function defined on all of $\mathbb{R}$. This inverse is called the \textit{exponential function}, denoted by $\exp(x)$ or $e^x$.
\end{definition}

\begin{theorem}{Basic Properties of the Exponential Function}
\begin{enumerate}
    \item[(a)] $\exp(0) = 1$ and $\exp(1) = e$ where $\log e = 1$
    \item[(b)] $\exp(a + b) = \exp(a) \cdot \exp(b)$ for all $a, b \in \mathbb{R}$
    \item[(c)] $\exp(a - b) = \exp(a)/\exp(b)$ for all $a, b \in \mathbb{R}$
    \item[(d)] $\exp(ra) = [\exp(a)]^r$ for all $a \in \mathbb{R}$ and rational $r$
\end{enumerate}
\end{theorem}

\textit{Proof.} These follow from the corresponding properties of the logarithm. For (b): Let $A = \exp(a)$ and $B = \exp(b)$. Then $a = \log A$ and $b = \log B$, so $a + b = \log A + \log B = \log(AB)$. Therefore $\exp(a + b) = AB = \exp(a)\exp(b)$. \qed

\begin{theorem}{Derivative and Integral of the Exponential Function}
\[
\frac{d}{dx} e^x = e^x \quad \text{for all } x \in \mathbb{R}
\]
\[
\int e^x\, dx = e^x + C
\]
\end{theorem}

\textit{Proof.} Let $y = e^x$, so $x = \log y$. By the inverse function theorem:
\[
\frac{dy}{dx} = \frac{1}{\frac{dx}{dy}} = \frac{1}{\frac{1}{y}} = y = e^x
\] \qed

\section{General Powers and Logarithms}

\begin{definition}{General Exponential}
For $a > 0$ and any real $x$, define:
\[
a^x = e^{x \log a}
\]
\end{definition}

\begin{theorem}{Derivatives and Integrals with General Bases}
For $a > 0$ and $a \neq 1$:
\[
\frac{d}{dx} a^x = a^x \log a
\]
\[
\frac{d}{dx} \log_a x = \frac{1}{x \log a}
\]
\end{theorem}

\textit{Proof.} $\frac{d}{dx} a^x = \frac{d}{dx} e^{x \log a} = e^{x \log a} \cdot \log a = a^x \log a$. \qed

\section{The Hyperbolic Functions}

\begin{definition}{Hyperbolic Functions}
For all real $x$, define:
\[
\sinh x = \frac{e^x - e^{-x}}{2}, \quad \cosh x = \frac{e^x + e^{-x}}{2}, \quad \tanh x = \frac{\sinh x}{\cosh x} = \frac{e^x - e^{-x}}{e^x + e^{-x}}
\]
\end{definition}

\begin{theorem}{Derivatives of Hyperbolic Functions}
\[
\frac{d}{dx} \sinh x = \cosh x, \quad \frac{d}{dx} \cosh x = \sinh x, \quad \frac{d}{dx} \tanh x = \text{sech}^2 x
\]
\end{theorem}

\textit{Proof.} $\frac{d}{dx} \sinh x = \frac{d}{dx}\left(\frac{e^x - e^{-x}}{2}\right) = \frac{e^x + e^{-x}}{2} = \cosh x$. Similarly for the others. \qed

\begin{theorem}{Fundamental Identity}
\[
\cosh^2 x - \sinh^2 x = 1
\]
\end{theorem}

\textit{Proof.} $\cosh^2 x - \sinh^2 x = \frac{(e^x + e^{-x})^2}{4} - \frac{(e^x - e^{-x})^2}{4} = \frac{4e^x e^{-x}}{4} = 1$. \qed

\section{Inverse Trigonometric Functions}

\begin{definition}{Inverse Sine}
The function $\sin x$ is strictly increasing and continuous on $[-\pi/2, \pi/2]$, mapping this interval onto $[-1, 1]$. Its inverse, denoted by $\arcsin x$ or $\sin^{-1} x$, is defined on $[-1, 1]$ with range $[-\pi/2, \pi/2]$.
\end{definition}

\begin{definition}{Inverse Cosine}
The function $\cos x$ is strictly decreasing and continuous on $[0, \pi]$, mapping this interval onto $[-1, 1]$. Its inverse, denoted by $\arccos x$ or $\cos^{-1} x$, is defined on $[-1, 1]$ with range $[0, \pi]$.
\end{definition}

\begin{definition}{Inverse Tangent}
The function $\tan x$ is strictly increasing and continuous on $(-\pi/2, \pi/2)$, mapping this interval onto $(-\infty, \infty)$. Its inverse, denoted by $\arctan x$ or $\tan^{-1} x$, is defined on all of $\mathbb{R}$ with range $(-\pi/2, \pi/2)$.
\end{definition}

\begin{theorem}{Derivatives of Inverse Trigonometric Functions}
\[
\frac{d}{dx} \arcsin x = \frac{1}{\sqrt{1-x^2}}, \quad \frac{d}{dx} \arccos x = -\frac{1}{\sqrt{1-x^2}} \quad \text{for } |x| < 1
\]
\[
\frac{d}{dx} \arctan x = \frac{1}{1+x^2} \quad \text{for all } x
\]
\end{theorem}

\textit{Proof.} For $\arcsin x$: Let $y = \arcsin x$, so $x = \sin y$. Then:
\[
\frac{dy}{dx} = \frac{1}{\frac{dx}{dy}} = \frac{1}{\cos y} = \frac{1}{\sqrt{1 - \sin^2 y}} = \frac{1}{\sqrt{1-x^2}}
\]
where we take the positive square root since $\cos y \geq 0$ for $y \in [-\pi/2, \pi/2]$. \qed

\begin{theorem}{Integration Formulas}
\[
\int \frac{dx}{\sqrt{1-x^2}} = \arcsin x + C = -\arccos x + C
\]
\[
\int \frac{dx}{1+x^2} = \arctan x + C
\]
\[
\int \frac{dx}{\sqrt{a^2 - x^2}} = \arcsin\frac{x}{a} + C \quad (a > 0)
\]
\[
\int \frac{dx}{a^2 + x^2} = \frac{1}{a}\arctan\frac{x}{a} + C \quad (a > 0)
\]
\end{theorem}

\textit{Proof.} These follow directly from the derivative formulas. \qed

\section{Integration by Partial Fractions}

\begin{theorem}{Partial Fraction Decomposition}
Let $P(x)/Q(x)$ be a rational function where $\deg P < \deg Q$ and $Q(x)$ factors as:
\[
Q(x) = (x-a_1)^{m_1} \cdots (x-a_k)^{m_k} (x^2+b_1x+c_1)^{n_1} \cdots (x^2+b_lx+c_l)^{n_l}
\]
Then $P(x)/Q(x)$ can be written as a sum of terms of the form:
\[
\frac{A}{(x-a)^r} \quad \text{and} \quad \frac{Bx+C}{(x^2+bx+c)^s}
\]
\end{theorem}

\begin{theorem}{Integration of Rational Functions}
Every rational function can be integrated in closed form using:
\begin{enumerate}
    \item[(a)] Partial fraction decomposition
    \item[(b)] $\displaystyle\int \frac{dx}{(x-a)^n} = \begin{cases} \frac{(x-a)^{1-n}}{1-n} & n \neq 1 \\ \log|x-a| & n = 1 \end{cases}$
    \item[(c)] For irreducible quadratics $x^2 + bx + c$, complete the square and use substitution
\end{enumerate}
\end{theorem}

\textit{Proof.} After partial fraction decomposition, each term can be integrated using standard techniques. For quadratic factors, completing the square gives $(x + b/2)^2 + (c - b^2/4)$, which leads to $\arctan$ or logarithmic forms. \qed



% ====================
% Part IV: Sequences and Series
% ====================
\part{Sequences and Series}

\chapter{Polynomial Approximations}

\section{The Taylor Polynomials Generated by a Function}

\begin{definition}{Taylor Polynomial}
Let $f$ be a function with derivatives of order $n$ at the point $a$. The polynomial:
\[
P_n(x) = f(a) + f'(a)(x-a) + \frac{f''(a)}{2!}(x-a)^2 + \cdots + \frac{f^{(n)}(a)}{n!}(x-a)^n
\]
is called the $n$th Taylor polynomial of $f$ at $a$, or generated by $f$ at $a$.
\end{definition}

\begin{theorem}{Approximation Property of Taylor Polynomials}
Let $f$ have derivatives of order $n$ at the point $a$, and let $P_n$ be the $n$th Taylor polynomial of $f$ at $a$. Then:
\[
P_n(a) = f(a), \quad P_n'(a) = f'(a), \quad P_n''(a) = f''(a), \quad \ldots, \quad P_n^{(n)}(a) = f^{(n)}(a)
\]
\end{theorem}

\textit{Proof.} Differentiating $P_n(x)$ term by term $k$ times $(0 \leq k \leq n)$ gives:
\[
P_n^{(k)}(x) = f^{(k)}(a) + \frac{f^{(k+1)}(a)}{1!}(x-a) + \cdots + \frac{f^{(n)}(a)}{(n-k)!}(x-a)^{n-k}
\]
Setting $x = a$, all terms with $(x-a)$ vanish, leaving $P_n^{(k)}(a) = f^{(k)}(a)$. \qed

\begin{theorem}{Uniqueness of the Approximating Polynomial}
Let $f$ have derivatives of order $n$ at $a$, and let $P$ be a polynomial of degree $\leq n$ such that:
\[
P(a) = f(a), \quad P'(a) = f'(a), \quad \ldots, \quad P^{(n)}(a) = f^{(n)}(a)
\]
Then $P = P_n$, the $n$th Taylor polynomial of $f$ at $a$.
\end{theorem}

\textit{Proof.} Write $P(x) = c_0 + c_1(x-a) + c_2(x-a)^2 + \cdots + c_n(x-a)^n$. Successive differentiation and evaluation at $x = a$ gives $c_k = f^{(k)}(a)/k!$ for each $k$, so $P = P_n$. \qed

\section{Calculus of Taylor Polynomials}

\begin{theorem}{Linearity Property}
If $P_n$ and $Q_n$ are the $n$th Taylor polynomials of $f$ and $g$ at $a$, then the $n$th Taylor polynomial of $\alpha f + \beta g$ is $\alpha P_n + \beta Q_n$.
\end{theorem}

\textit{Proof.} By the uniqueness theorem, since $\alpha P_n + \beta Q_n$ matches the derivatives of $\alpha f + \beta g$ up to order $n$ at $a$, it is the Taylor polynomial. \qed

\begin{theorem}{Differentiation Property}
The derivative of the $n$th Taylor polynomial of $f$ at $a$ is the $(n-1)$st Taylor polynomial of $f'$ at $a$.
\end{theorem}

\textit{Proof.} Differentiating $P_n(x)$ gives:
\[
P_n'(x) = f'(a) + \frac{f''(a)}{1!}(x-a) + \cdots + \frac{f^{(n)}(a)}{(n-1)!}(x-a)^{n-1}
\]
which is precisely the $(n-1)$st Taylor polynomial of $f'$ at $a$. \qed

\begin{theorem}{Integration Property}
If $f$ has $n+1$ derivatives at $a$, the $(n+1)$st Taylor polynomial of $P(x) = \int_a^x f(t)\, dt$ is $\int_a^x P_n(t)\, dt$.
\end{theorem}

\textit{Proof.} Let $Q(x) = \int_a^x P_n(t)\, dt$. Then $Q(a) = 0$ and $Q'(x) = P_n(x)$. By the differentiation property, $Q^{(k+1)}(a) = P_n^{(k)}(a) = f^{(k)}(a)$ for $0 \leq k \leq n$, so $Q^{(k+1)}(a) = f^{(k)}(a) = P^{(k+1)}(a)$. Since $P(a) = 0 = Q(a)$, we have $Q^{(k)}(a) = P^{(k)}(a)$ for $0 \leq k \leq n+1$, proving $Q$ is the Taylor polynomial of $P$. \qed

\section{Taylor's Formula with Remainder}

\begin{theorem}{Taylor's Formula with Lagrange Remainder}
Let $f$ have derivatives of order $n+1$ in some interval containing $a$. Then for each $x$ in this interval, there exists a point $c$ between $a$ and $x$ such that:
\[
f(x) = P_n(x) + R_n(x)
\]
where $P_n$ is the $n$th Taylor polynomial and:
\[
R_n(x) = \frac{f^{(n+1)}(c)}{(n+1)!}(x-a)^{n+1}
\]
\end{theorem}

\textit{Proof.} Fix $x \neq a$ and define for $t$ between $a$ and $x$:
\[
F(t) = f(x) - f(t) - f'(t)(x-t) - \frac{f''(t)}{2!}(x-t)^2 - \cdots - \frac{f^{(n)}(t)}{n!}(x-t)^n
\]
Then $F(a) = R_n(x)$ and $F(x) = 0$. Differentiating:
\[
F'(t) = -\frac{f^{(n+1)}(t)}{n!}(x-t)^n
\]
By the mean-value theorem, there exists $c$ between $a$ and $x$ such that:
\[
\frac{F(x) - F(a)}{x - a} = F'(c) = -\frac{f^{(n+1)}(c)}{n!}(x-c)^n
\]
Since $F(x) - F(a) = -R_n(x)$, we get $R_n(x) = \frac{f^{(n+1)}(c)}{n!}(x-c)^n(x-a)$. A more careful analysis (using Cauchy's mean-value theorem) gives the stated form. \qed

\begin{theorem}{Estimate of the Remainder}
If $f$ has derivatives of order $n+1$ in an interval $I$ containing $a$, and if $|f^{(n+1)}(x)| \leq M$ for all $x \in I$, then:
\[
|R_n(x)| \leq \frac{M}{(n+1)!}|x-a|^{n+1}
\]
for all $x \in I$.
\end{theorem}

\textit{Proof.} From Taylor's formula with $|f^{(n+1)}(c)| \leq M$:
\[
|R_n(x)| = \left|\frac{f^{(n+1)}(c)}{(n+1)!}\right||x-a|^{n+1} \leq \frac{M}{(n+1)!}|x-a|^{n+1}
\] \qed

\section{Other Forms of the Remainder in Taylor's Formula}

\begin{theorem}{Taylor's Formula with Integral Remainder}
If $f$ has continuous derivative of order $n+1$ in an interval $I$ containing $a$, then for all $x \in I$:
\[
R_n(x) = \frac{1}{n!}\int_a^x (x-t)^n f^{(n+1)}(t)\, dt
\]
\end{theorem}

\textit{Proof.} Starting from $R_n(a) = 0$ and $R_n'(x) = f'(x) - P_n'(x)$, with $R_n'(a) = 0$ for the $(n-1)$st Taylor polynomial of $f'$, successive integration by parts yields the integral form. Specifically, with $R_n^{(n+1)}(x) = f^{(n+1)}(x)$ and $R_n^{(k)}(a) = 0$ for $0 \leq k \leq n$:
\[
R_n(x) = \int_a^x R_n'(t)\, dt = \int_a^x \int_a^t R_n''(s)\, ds\, dt = \cdots = \frac{1}{n!}\int_a^x (x-t)^n f^{(n+1)}(t)\, dt
\] \qed

\begin{theorem}{Taylor's Formula with Cauchy Remainder}
Under the same hypotheses as the integral form, for each $x \in I$, there exists $c$ between $a$ and $x$ such that:
\[
R_n(x) = \frac{f^{(n+1)}(c)}{n!}(x-c)^n(x-a)
\]
\end{theorem}

\textit{Proof.} Apply the weighted mean-value theorem for integrals to the integral remainder with $g(t) = (x-t)^n$. There exists $c$ between $a$ and $x$ such that:
\[
\int_a^x (x-t)^n f^{(n+1)}(t)\, dt = f^{(n+1)}(c)(x-c)^n(x-a)
\]
Substituting into the integral form gives the result. \qed

\section{Error Estimates}

\begin{theorem}{Little-o Notation for Remainder}
If $f$ has derivatives of order $n$ at $a$, then:
\[
R_n(x) = o(|x-a|^n) \quad \text{as } x \to a
\]
That is, $\lim_{x \to a} \frac{R_n(x)}{(x-a)^n} = 0$.
\end{theorem}

\textit{Proof.} By Taylor's formula with the remainder estimate, if $f^{(n+1)}$ is bounded near $a$, then $|R_n(x)| \leq M|x-a|^{n+1}/(n+1)! = o(|x-a|^n)$. \qed

\section{Applications to Indeterminate Forms}

\begin{definition}{Indeterminate Form 0/0}
If $\lim_{x \to a} f(x) = \lim_{x \to a} g(x) = 0$, the quotient $f(x)/g(x)$ is said to assume the indeterminate form $0/0$ at $x = a$.
\end{definition}

\begin{theorem}{L'Hôpital's Rule for 0/0}
Assume $f$ and $g$ have derivatives $f'$ and $g'$ at each point of an open interval $(a, b)$, and that $g'(x) \neq 0$ for all $x \in (a, b)$. Assume also that:
\[
\lim_{x \to a^+} f(x) = \lim_{x \to a^+} g(x) = 0
\]
If $\lim_{x \to a^+} \frac{f'(x)}{g'(x)} = L$ (where $L$ may be finite or infinite), then:
\[
\lim_{x \to a^+} \frac{f(x)}{g(x)} = L
\]
\end{theorem}

\textit{Proof.} Extend $f$ and $g$ to be continuous at $a$ by defining $f(a) = g(a) = 0$. By the generalized mean-value theorem, for each $x \in (a, b)$, there exists $c \in (a, x)$ such that:
\[
\frac{f(x)}{g(x)} = \frac{f(x) - f(a)}{g(x) - g(a)} = \frac{f'(c)}{g'(c)}
\]
As $x \to a^+$, we have $c \to a^+$, so:
\[
\lim_{x \to a^+} \frac{f(x)}{g(x)} = \lim_{c \to a^+} \frac{f'(c)}{g'(c)} = L
\] \qed

\begin{theorem}{L'Hôpital's Rule for $\infty/\infty$}
Assume $f$ and $g$ have derivatives on $(a, b)$ with $g'(x) \neq 0$. If $\lim_{x \to a^+} |f(x)| = \lim_{x \to a^+} |g(x)| = \infty$ and $\lim_{x \to a^+} \frac{f'(x)}{g'(x)} = L$, then:
\[
\lim_{x \to a^+} \frac{f(x)}{g(x)} = L
\]
\end{theorem}

\textit{Proof.} We prove the case where $L$ is finite. Given $\epsilon > 0$, there exists $c \in (a, b)$ such that $\left|\frac{f'(t)}{g'(t)} - L\right| < \epsilon$ for all $t \in (a, c)$. For $a < x < y < c$, by the generalized mean-value theorem, there exists $t \in (x, y)$ with:
\[
\frac{f(y) - f(x)}{g(y) - g(x)} = \frac{f'(t)}{g'(t)}
\]
Rearranging: $\frac{f(x)}{g(x)} = \frac{f(y)}{g(x)} + \left(1 - \frac{g(y)}{g(x)}\right)\frac{f'(t)}{g'(t)}$.

As $x \to a^+$ with $y$ fixed, since $|g(x)| \to \infty$, we have $\frac{f(y)}{g(x)} \to 0$ and $\frac{g(y)}{g(x)} \to 0$. Thus for $x$ close enough to $a$:
\[
\left|\frac{f(x)}{g(x)} - \frac{f'(t)}{g'(t)}\right| < \epsilon
\]
Since $\left|\frac{f'(t)}{g'(t)} - L\right| < \epsilon$, we get $\left|\frac{f(x)}{g(x)} - L\right| < 2\epsilon$. \qed

\section{The Symbols $\infty$ and $-\infty$}

\begin{definition}{Limit to Infinity}
We say $\lim_{x \to a} f(x) = +\infty$ if for every $M > 0$, there exists $\delta > 0$ such that $f(x) > M$ whenever $0 < |x-a| < \delta$.
\end{definition}

\begin{definition}{Infinite Limits}
\[
\lim_{x \to a^-} f(x) = +\infty
\]
means that for every $M > 0$, there exists $\delta > 0$ such that $f(x) > M$ whenever $a - \delta < x < a$.

\[
\lim_{x \to a^+} f(x) = +\infty
\]
means that for every $M > 0$, there exists $\delta > 0$ such that $f(x) > M$ whenever $a < x < a + \delta$.
\end{definition}

\begin{theorem}{Behavior of $\log x$ and $e^x$ for Large $x$}
For any $\alpha > 0$:
\[
\lim_{x \to +\infty} \frac{\log x}{x^\alpha} = 0
\]
\[
\lim_{x \to +\infty} \frac{x^\alpha}{e^x} = 0
\]
\end{theorem}

\textit{Proof.} For the first limit, apply L'Hôpital's rule:
\[
\lim_{x \to +\infty} \frac{\log x}{x^\alpha} = \lim_{x \to +\infty} \frac{1/x}{\alpha x^{\alpha-1}} = \lim_{x \to +\infty} \frac{1}{\alpha x^\alpha} = 0
\]

For the second limit, apply L'Hôpital's rule $n$ times where $n > \alpha$:
\[
\lim_{x \to +\infty} \frac{x^\alpha}{e^x} = \lim_{x \to +\infty} \frac{\alpha x^{\alpha-1}}{e^x} = \cdots = \lim_{x \to +\infty} \frac{\alpha(\alpha-1)\cdots(\alpha-n+1)x^{\alpha-n}}{e^x} = 0
\]
since $\alpha - n < 0$. \qed

\begin{theorem}{Exponential vs. Polynomial Growth}
For every $n > 0$:
\[
\lim_{x \to +\infty} \frac{x^n}{e^x} = 0
\]
\[
\lim_{x \to +\infty} \frac{e^x}{x^n} = +\infty
\]
\end{theorem}

\textit{Proof.} Apply L'Hôpital's rule $n$ times to the first limit; each application reduces the power of $x$ by 1 until the numerator is constant and the denominator is $e^x$. The second limit follows from the first. \qed



\chapter{Differential Equations}

\section{Terminology and Notation}

\begin{definition}{Differential Equation}
A differential equation is an equation involving an unknown function and one or more of its derivatives. An equation of the form $y' = f(x, y)$ is called a first-order differential equation. An equation of the form $y'' = f(x, y, y')$ is called a second-order differential equation.
\end{definition}

\begin{definition}{Solution of a Differential Equation}
A function $y = Y(x)$ is said to be a solution of the differential equation $y' = f(x, y)$ on an interval $I$ if $Y$ is differentiable on $I$ and satisfies $Y'(x) = f(x, Y(x))$ for all $x \in I$.
\end{definition}

\section{First-Order Linear Differential Equations}

\begin{definition}{First-Order Linear Equation}
A first-order differential equation of the form:
\[
y' + P(x)y = Q(x)
\]
where $P$ and $Q$ are given functions, is called a first-order linear differential equation.
\end{definition}

\begin{theorem}{General Solution of First-Order Linear Equations}
Assume $P$ and $Q$ are continuous on an open interval $I$. Choose any point $a$ in $I$ and let $b$ be any real number. Then there is one and only one function $y = f(x)$ which satisfies the initial-value problem:
\[
y' + P(x)y = Q(x), \quad \text{with } f(a) = b
\]
on the interval $I$. This function is given by the formula:
\[
f(x) = be^{-A(x)} + e^{-A(x)}\int_a^x Q(t)e^{A(t)}\, dt
\]
where $A(x) = \int_a^x P(t)\, dt$.
\end{theorem}

\textit{Proof.} First we verify that the function satisfies the differential equation. Let $v(x) = e^{A(x)}$. Then $v'(x) = P(x)e^{A(x)} = P(x)v(x)$. Multiplying the differential equation by $v(x)$:
\[
v(x)y' + P(x)v(x)y = v(x)Q(x)
\]
Since $v'(x) = P(x)v(x)$, the left side is $(v(x)y)'$. Integrating from $a$ to $x$:
\[
v(x)y(x) - v(a)y(a) = \int_a^x Q(t)v(t)\, dt
\]
Since $v(a) = e^{A(a)} = e^0 = 1$ and $y(a) = b$:
\[
y(x) = \frac{b + \int_a^x Q(t)v(t)\, dt}{v(x)} = be^{-A(x)} + e^{-A(x)}\int_a^x Q(t)e^{A(t)}\, dt
\]
For uniqueness, suppose $f$ and $g$ are both solutions with $f(a) = g(a) = b$. Let $h = f - g$. Then $h(a) = 0$ and $h' + P(x)h = 0$. Multiplying by the integrating factor $e^{A(x)}$:
\[
\frac{d}{dx}[h(x)e^{A(x)}] = 0
\]
So $h(x)e^{A(x)} = h(a)e^{A(a)} = 0$, implying $h(x) = 0$ for all $x \in I$. Thus $f = g$. \qed

\section{Existence and Uniqueness Theorems}

\begin{theorem}{Existence and Uniqueness for $y' = f(x, y)$}
Let $f(x, y)$ be continuous and satisfy a Lipschitz condition in $y$ on a rectangle $R = [a, b] \times [c, d]$. That is, there exists $L > 0$ such that:
\[
|f(x, y_1) - f(x, y_2)| \leq L|y_1 - y_2|
\]
for all $(x, y_1), (x, y_2) \in R$. Then for any $(x_0, y_0)$ in the interior of $R$, there exists a unique solution to the initial-value problem $y' = f(x, y)$, $y(x_0) = y_0$ on some interval containing $x_0$.
\end{theorem}

\textit{Proof.} (Sketch) The proof uses Picard iteration. Define a sequence of functions recursively:
\[
Y_0(x) = y_0, \quad Y_{n+1}(x) = y_0 + \int_{x_0}^x f(t, Y_n(t))\, dt
\]
One shows this sequence converges uniformly to a solution. The Lipschitz condition ensures the solution is unique by showing any two solutions must coincide. \qed

\section{Second-Order Linear Differential Equations with Constant Coefficients}

\begin{definition}{Second-Order Linear Equation}
A differential equation of the form:
\[
y'' + ay' + by = R(x)
\]
where $a$ and $b$ are constants and $R$ is a given function, is called a second-order linear differential equation with constant coefficients. The equation is homogeneous if $R(x) = 0$.
\end{definition}

\begin{definition}{Characteristic Equation}
For the homogeneous equation $y'' + ay' + by = 0$, the quadratic equation:
\[
r^2 + ar + b = 0
\]
is called the characteristic equation. Its roots are denoted by $r_1$ and $r_2$.
\end{definition}

\begin{theorem}{General Solution of Homogeneous Equation}
The general solution of $y'' + ay' + by = 0$ depends on the discriminant $D = a^2 - 4b$:
\begin{enumerate}
    \item[(a)] If $D > 0$ (distinct real roots $r_1 \neq r_2$): $y = c_1e^{r_1x} + c_2e^{r_2x}$
    \item[(b)] If $D = 0$ (repeated root $r$): $y = (c_1 + c_2x)e^{rx}$
    \item[(c)] If $D < 0$ (complex roots $\alpha \pm i\beta$): $y = e^{\alpha x}(c_1\cos\beta x + c_2\sin\beta x)$
\end{enumerate}
\end{theorem}

\textit{Proof.} Substituting $y = e^{rx}$ into the differential equation gives the characteristic equation.

Case (a): If $r_1 \neq r_2$ are real, both $e^{r_1x}$ and $e^{r_2x}$ satisfy the equation. Their Wronskian $W = (r_2 - r_1)e^{(r_1+r_2)x} \neq 0$, so they are linearly independent.

Case (b): If $r$ is a repeated root, $e^{rx}$ and $xe^{rx}$ are solutions. One verifies directly that $y = xe^{rx}$ satisfies the equation.

Case (c): For complex roots $\alpha \pm i\beta$, using Euler's formula $e^{(\alpha+i\beta)x} = e^{\alpha x}(\cos\beta x + i\sin\beta x)$, the real and imaginary parts give real solutions $e^{\alpha x}\cos\beta x$ and $e^{\alpha x}\sin\beta x$. \qed

\begin{theorem}{General Solution of Nonhomogeneous Equation}
If $y_p$ is any particular solution of $y'' + ay' + by = R(x)$, and $y_1, y_2$ are linearly independent solutions of the homogeneous equation, then every solution has the form:
\[
y = c_1y_1 + c_2y_2 + y_p
\]
where $c_1, c_2$ are constants.
\end{theorem}

\textit{Proof.} Let $y$ be any solution and $y_p$ be a particular solution. Then $(y - y_p)'' + a(y - y_p)' + b(y - y_p) = R - R = 0$. So $y - y_p$ satisfies the homogeneous equation, hence $y - y_p = c_1y_1 + c_2y_2$ for some constants $c_1, c_2$. \qed

\begin{theorem}{Variation of Parameters}
If $y_1$ and $y_2$ are linearly independent solutions of $y'' + ay' + by = 0$, then a particular solution of $y'' + ay' + by = R(x)$ is:
\[
y_p = -y_1\int \frac{y_2R}{W}\, dx + y_2\int \frac{y_1R}{W}\, dx
\]
where $W = y_1y_2' - y_2y_1'$ is the Wronskian.
\end{theorem}

\textit{Proof.} Try a solution of the form $y_p = v_1y_1 + v_2y_2$ where $v_1$ and $v_2$ are functions to be determined. Imposing the condition $v_1'y_1 + v_2'y_2 = 0$, we get:
\[
y_p' = v_1y_1' + v_2y_2', \quad y_p'' = v_1'y_1' + v_2'y_2' + v_1y_1'' + v_2y_2''
\]
Substituting into the differential equation and using that $y_1, y_2$ satisfy the homogeneous equation:
\[
v_1'y_1' + v_2'y_2' = R
\]
Solving the system:
\[
\begin{cases} v_1'y_1 + v_2'y_2 = 0 \\ v_1'y_1' + v_2'y_2' = R \end{cases}
\]
gives $v_1' = -\frac{y_2R}{W}$ and $v_2' = \frac{y_1R}{W}$. Integrating yields the formula. \qed

\section{Separable First-Order Equations}

\begin{definition}{Separable Equation}
A first-order differential equation of the form $y' = f(x, y)$ is called separable if $f(x, y)$ can be written as:
\[
f(x, y) = \frac{Q(x)}{P(y)}
\]
where $P$ and $Q$ are functions of $y$ and $x$ respectively.
\end{definition}

\begin{theorem}{Solution of Separable Equations}
If $y' = Q(x)/P(y)$ with $P(y) \neq 0$, and $G$ is any primitive of $P$ (i.e., $G' = P$), then every solution satisfies:
\[
G(y) = \int Q(x)\, dx + C
\]
for some constant $C$.
\end{theorem}

\textit{Proof.} Write the equation as $P(y)\frac{dy}{dx} = Q(x)$. If $y = Y(x)$ is a solution, then by the chain rule:
\[
\frac{d}{dx}G(Y(x)) = P(Y(x))Y'(x) = Q(x)
\]
Integrating both sides: $G(Y(x)) = \int Q(x)\, dx + C$. \qed

\section{Homogeneous First-Order Equations}

\begin{definition}{Homogeneous First-Order Equation}
A first-order differential equation $y' = f(x, y)$ is called homogeneous (in the sense of differential equations) if $f(tx, ty) = f(x, y)$ for all $t \neq 0$. Equivalently, $f$ can be written as $f(x, y) = g(y/x)$.
\end{definition}

\begin{theorem}{Solution of Homogeneous Equations}
The homogeneous equation $y' = g(y/x)$ can be transformed into a separable equation by the substitution $v = y/x$ (i.e., $y = vx$). Then:
\[
\frac{dv}{g(v) - v} = \frac{dx}{x}
\]
\end{theorem}

\textit{Proof.} With $y = vx$, we have $y' = v + xv'$. The equation becomes:
\[
v + xv' = g(v)
\]
Rearranging: $xv' = g(v) - v$, or $\frac{dv}{g(v) - v} = \frac{dx}{x}$. This is separable in $v$ and $x$. \qed




\chapter{Complex Numbers}

\section{Definition and Field Properties}

\begin{definition}{Complex Number}
A complex number is an ordered pair $(a, b)$ of real numbers. Two complex numbers $(a, b)$ and $(c, d)$ are equal if and only if $a = c$ and $b = d$. The set of all complex numbers is denoted by $\mathbb{C}$.
\end{definition}

\begin{definition}{Addition and Multiplication}
For complex numbers $(a, b)$ and $(c, d)$, define:
\[
(a, b) + (c, d) = (a + c, b + d)
\]
\[
(a, b) \cdot (c, d) = (ac - bd, ad + bc)
\]
\end{definition}

\begin{theorem}{Field Properties}
The set $\mathbb{C}$ with the operations of addition and multiplication forms a field. That is:
\begin{enumerate}
    \item[(a)] Addition is commutative and associative
    \item[(b)] Multiplication is commutative and associative
    \item[(c)] Distributive law: $z_1(z_2 + z_3) = z_1z_2 + z_1z_3$
    \item[(d)] Additive identity: $(0, 0)$ satisfies $(a, b) + (0, 0) = (a, b)$
    \item[(e)] Multiplicative identity: $(1, 0)$ satisfies $(a, b) \cdot (1, 0) = (a, b)$
    \item[(f)] Additive inverses exist for all $(a, b)$
    \item[(g)] Multiplicative inverses exist for all $(a, b) \neq (0, 0)$
\end{enumerate}
\end{theorem}

\textit{Proof.} The algebraic properties follow directly from the definitions and the corresponding properties of real numbers. For multiplicative inverses, if $(a, b) \neq (0, 0)$, then:
\[
(a, b)^{-1} = \left(\frac{a}{a^2 + b^2}, \frac{-b}{a^2 + b^2}\right)
\]
Verification: $(a, b) \cdot \left(\frac{a}{a^2 + b^2}, \frac{-b}{a^2 + b^2}\right) = \left(\frac{a^2 + b^2}{a^2 + b^2}, \frac{-ab + ab}{a^2 + b^2}\right) = (1, 0)$. \qed

\section{The Imaginary Unit}

\begin{definition}{Imaginary Unit}
The imaginary unit $i$ is defined as the complex number $i = (0, 1)$.
\end{definition}

\begin{theorem}{Basic Property of $i$}
\[
i^2 = -1
\]
\end{theorem}

\textit{Proof.} $i^2 = (0, 1) \cdot (0, 1) = (0 \cdot 0 - 1 \cdot 1, 0 \cdot 1 + 1 \cdot 0) = (-1, 0) = -1$. \qed

\begin{theorem}{Standard Form}
Every complex number $(a, b)$ can be written uniquely as:
\[
(a, b) = a + bi
\]
where $a$ and $b$ are real numbers.
\end{theorem}

\textit{Proof.} $(a, b) = (a, 0) + (0, b) = (a, 0) + (b, 0) \cdot (0, 1) = a + bi$, identifying real number $a$ with complex number $(a, 0)$. \qed

\section{Geometric Interpretation}

\begin{definition}{Modulus and Argument}
For a complex number $z = x + yi$:
\begin{enumerate}
    \item[(a)] The \textit{modulus} (absolute value) is $|z| = \sqrt{x^2 + y^2}$
    \item[(b)] The \textit{argument} $\theta = \arg(z)$ satisfies $\cos\theta = x/|z|$ and $\sin\theta = y/|z|$ (for $z \neq 0$)
\end{enumerate}
\end{definition}

\begin{theorem}{Polar Form}
Every nonzero complex number can be written in polar form:
\[
z = r(\cos\theta + i\sin\theta)
\]
where $r = |z|$ and $\theta = \arg(z)$.
\end{theorem}

\textit{Proof.} If $z = x + yi$, then $x = r\cos\theta$ and $y = r\sin\theta$ where $r = \sqrt{x^2 + y^2}$. Thus $z = r(\cos\theta + i\sin\theta)$. \qed

\begin{theorem}{Multiplication in Polar Form}
If $z_1 = r_1(\cos\theta_1 + i\sin\theta_1)$ and $z_2 = r_2(\cos\theta_2 + i\sin\theta_2)$, then:
\[
z_1 \cdot z_2 = r_1r_2[\cos(\theta_1 + \theta_2) + i\sin(\theta_1 + \theta_2)]
\]
and:
\[
\frac{z_1}{z_2} = \frac{r_1}{r_2}[\cos(\theta_1 - \theta_2) + i\sin(\theta_1 - \theta_2)]
\]
\end{theorem}

\textit{Proof.} Using the angle addition formulas:
\begin{align*}
z_1 \cdot z_2 &= r_1r_2[(\cos\theta_1\cos\theta_2 - \sin\theta_1\sin\theta_2) + i(\sin\theta_1\cos\theta_2 + \cos\theta_1\sin\theta_2)] \\
&= r_1r_2[\cos(\theta_1 + \theta_2) + i\sin(\theta_1 + \theta_2)]
\end{align*}
The division formula follows similarly. \qed

\section{Exponential Form}

\begin{definition}{Complex Exponential}
For a complex number $z = x + iy$, define:
\[
e^z = e^x(\cos y + i\sin y)
\]
In particular, for pure imaginary $z = i\theta$:
\[
e^{i\theta} = \cos\theta + i\sin\theta
\]
\end{definition}

\begin{theorem}{Euler's Formula}
For any real number $\theta$:
\[
e^{i\theta} = \cos\theta + i\sin\theta
\]
In particular:
\[
\cos\theta = \frac{e^{i\theta} + e^{-i\theta}}{2}, \quad \sin\theta = \frac{e^{i\theta} - e^{-i\theta}}{2i}
\]
\end{theorem}

\textit{Proof.} This follows directly from the definition of the complex exponential. The formulas for cosine and sine are obtained by adding and subtracting $e^{i\theta}$ and $e^{-i\theta}$. \qed

\begin{theorem}{Properties of the Complex Exponential}
For all complex numbers $z, z_1, z_2$:
\begin{enumerate}
    \item[(a)] $e^{z_1 + z_2} = e^{z_1} \cdot e^{z_2}$
    \item[(b)] $e^{z_1 - z_2} = e^{z_1}/e^{z_2}$
    \item[(c)] $e^z \neq 0$ for all $z \in \mathbb{C}$
    \item[(d)] $|e^z| = e^{\text{Re}(z)}$
    \item[(e)] $e^{z + 2\pi i} = e^z$ (periodic with period $2\pi i$)
\end{enumerate}
\end{theorem}

\textit{Proof.} (a) Let $z_1 = x_1 + iy_1$ and $z_2 = x_2 + iy_2$. Then:
\begin{align*}
e^{z_1} \cdot e^{z_2} &= e^{x_1}(\cos y_1 + i\sin y_1) \cdot e^{x_2}(\cos y_2 + i\sin y_2) \\
&= e^{x_1+x_2}[\cos(y_1+y_2) + i\sin(y_1+y_2)] \\
&= e^{z_1+z_2}
\end{align*}
using the multiplication rule for polar forms.

(d) $|e^z| = |e^x(\cos y + i\sin y)| = e^x \cdot 1 = e^{\text{Re}(z)}$.

(e) $e^{z + 2\pi i} = e^x[\cos(y + 2\pi) + i\sin(y + 2\pi)] = e^x[\cos y + i\sin y] = e^z$. \qed

\section{De Moivre's Formula}

\begin{theorem}{De Moivre's Formula}
For any integer $n$ and real number $\theta$:
\[
(\cos\theta + i\sin\theta)^n = \cos(n\theta) + i\sin(n\theta)
\]
or equivalently:
\[
(e^{i\theta})^n = e^{in\theta}
\]
\end{theorem}

\textit{Proof.} By induction using the multiplication rule for polar forms. For positive $n$, repeated application gives $(\cos\theta + i\sin\theta)^n = \cos(n\theta) + i\sin(n\theta)$. For negative $n$, use the reciprocal formula. \qed

\section{$n$th Roots of Complex Numbers}

\begin{theorem}{$n$th Roots}
For a nonzero complex number $z = r(\cos\theta + i\sin\theta)$ and positive integer $n$, there are exactly $n$ distinct $n$th roots given by:
\[
z^{1/n} = r^{1/n}\left[\cos\left(\frac{\theta + 2\pi k}{n}\right) + i\sin\left(\frac{\theta + 2\pi k}{n}\right)\right]
\]
for $k = 0, 1, 2, \ldots, n-1$.
\end{theorem}

\textit{Proof.} We seek $w = \rho(\cos\phi + i\sin\phi)$ such that $w^n = z$. By De Moivre's formula:
\[
\rho^n(\cos(n\phi) + i\sin(n\phi)) = r(\cos\theta + i\sin\theta)
\]
This requires $\rho^n = r$ (so $\rho = r^{1/n}$) and $n\phi = \theta + 2\pi k$ for integer $k$. Thus $\phi = (\theta + 2\pi k)/n$. The distinct values occur for $k = 0, 1, \ldots, n-1$. \qed



\chapter{Sequences and Infinite Series}

\section{Sequences}

\begin{definition}{Sequence}
A sequence is a function whose domain is the set of positive integers. If $f$ is a sequence, the value $f(n)$ is called the $n$th term of the sequence and is denoted by $a_n$ or $f_n$. The sequence is denoted by $\{a_n\}$.
\end{definition}

\begin{definition}{Convergence of a Sequence}
A sequence $\{a_n\}$ is said to converge to a limit $L$, written $\lim_{n \to \infty} a_n = L$, if for every $\epsilon > 0$ there exists an integer $N$ such that:
\[
|a_n - L| < \epsilon \quad \text{whenever } n \geq N
\]
If no such $L$ exists, the sequence is said to diverge.
\end{definition}

\begin{theorem}{Uniqueness of Limits}
If a sequence converges, its limit is unique.
\end{theorem}

\textit{Proof.} Suppose $\lim_{n \to \infty} a_n = L_1$ and $\lim_{n \to \infty} a_n = L_2$. For any $\epsilon > 0$, there exists $N_1, N_2$ such that $|a_n - L_1| < \epsilon/2$ for $n \geq N_1$ and $|a_n - L_2| < \epsilon/2$ for $n \geq N_2$. For $n \geq \max(N_1, N_2)$:
\[
|L_1 - L_2| \leq |L_1 - a_n| + |a_n - L_2| < \epsilon/2 + \epsilon/2 = \epsilon
\]
Since $\epsilon$ is arbitrary, $|L_1 - L_2| = 0$, so $L_1 = L_2$. \qed

\begin{theorem}{Properties of Limits of Sequences}
If $\lim_{n \to \infty} a_n = A$ and $\lim_{n \to \infty} b_n = B$, then:
\begin{enumerate}
    \item[(a)] $\lim_{n \to \infty} (a_n + b_n) = A + B$
    \item[(b)] $\lim_{n \to \infty} (a_n \cdot b_n) = A \cdot B$
    \item[(c)] $\lim_{n \to \infty} (ca_n) = cA$ for any constant $c$
    \item[(d)] $\lim_{n \to \infty} \frac{a_n}{b_n} = \frac{A}{B}$, provided $B \neq 0$
\end{enumerate}
\end{theorem}

\textit{Proof.} (a) Given $\epsilon > 0$, choose $N$ such that $|a_n - A| < \epsilon/2$ and $|b_n - B| < \epsilon/2$ for $n \geq N$. Then:
\[
|(a_n + b_n) - (A + B)| \leq |a_n - A| + |b_n - B| < \epsilon/2 + \epsilon/2 = \epsilon
\]

(b) Since $\{a_n\}$ converges, it is bounded, say $|a_n| \leq M$. Given $\epsilon > 0$, choose $N$ such that for $n \geq N$: $|a_n - A| < \epsilon/(2|B| + 1)$ and $|b_n - B| < \epsilon/(2M)$. Then:
\begin{align*}
|a_nb_n - AB| &= |a_nb_n - a_nB + a_nB - AB| \\
&\leq |a_n||b_n - B| + |B||a_n - A| \\
&< M \cdot \frac{\epsilon}{2M} + |B| \cdot \frac{\epsilon}{2|B| + 1} < \epsilon
\end{align*}
\qed

\begin{theorem}{Sandwich Theorem for Sequences}
If $a_n \leq b_n \leq c_n$ for all $n \geq N$ for some $N$, and if $\lim_{n \to \infty} a_n = \lim_{n \to \infty} c_n = L$, then $\lim_{n \to \infty} b_n = L$.
\end{theorem}

\textit{Proof.} Given $\epsilon > 0$, there exists $N_1$ such that $|a_n - L| < \epsilon$ and $|c_n - L| < \epsilon$ for $n \geq N_1$. For $n \geq \max(N, N_1)$:
\[
L - \epsilon < a_n \leq b_n \leq c_n < L + \epsilon
\]
Thus $|b_n - L| < \epsilon$. \qed

\begin{theorem}{Monotone Convergence Theorem}
A monotone increasing sequence that is bounded above converges. A monotone decreasing sequence that is bounded below converges.
\end{theorem}

\textit{Proof.} Let $\{a_n\}$ be monotone increasing and bounded above. Let $L = \sup\{a_n : n \geq 1\}$. For any $\epsilon > 0$, since $L - \epsilon$ is not an upper bound, there exists $N$ such that $a_N > L - \epsilon$. Since $\{a_n\}$ is increasing, for $n \geq N$: $L - \epsilon < a_N \leq a_n \leq L < L + \epsilon$. Thus $|a_n - L| < \epsilon$ for $n \geq N$. \qed

\section{Infinite Series}

\begin{definition}{Infinite Series}
Given a sequence $\{a_n\}$, the expression $\sum_{n=1}^{\infty} a_n$ is called an infinite series. The $n$th partial sum is $s_n = a_1 + a_2 + \cdots + a_n = \sum_{k=1}^{n} a_k$. The series converges if the sequence of partial sums $\{s_n\}$ converges, and we write $\sum_{n=1}^{\infty} a_n = \lim_{n \to \infty} s_n$.
\end{definition}

\begin{theorem}{Necessary Condition for Convergence}
If $\sum_{n=1}^{\infty} a_n$ converges, then $\lim_{n \to \infty} a_n = 0$.
\end{theorem}

\textit{Proof.} Let $s_n = \sum_{k=1}^{n} a_k$. If the series converges, $\lim_{n \to \infty} s_n = S$ exists. Then $a_n = s_n - s_{n-1}$, so:
\[
\lim_{n \to \infty} a_n = \lim_{n \to \infty} s_n - \lim_{n \to \infty} s_{n-1} = S - S = 0
\] \qed

\begin{theorem}{Linear Property of Series}
If $\sum_{n=1}^{\infty} a_n = A$ and $\sum_{n=1}^{\infty} b_n = B$, then:
\[
\sum_{n=1}^{\infty} (ca_n) = cA \quad \text{and} \quad \sum_{n=1}^{\infty} (a_n + b_n) = A + B
\]
\end{theorem}

\textit{Proof.} Follows from the corresponding properties for sequences applied to the partial sums. \qed

\begin{theorem}{Comparison Test}
Suppose $0 \leq a_n \leq b_n$ for all $n \geq N$.
\begin{enumerate}
    \item[(a)] If $\sum_{n=1}^{\infty} b_n$ converges, then $\sum_{n=1}^{\infty} a_n$ converges.
    \item[(b)] If $\sum_{n=1}^{\infty} a_n$ diverges, then $\sum_{n=1}^{\infty} b_n$ diverges.
\end{enumerate}
\end{theorem}

\textit{Proof.} (a) Let $s_n = \sum_{k=1}^{n} a_k$ and $t_n = \sum_{k=1}^{n} b_k$. For $n \geq N$: $s_n \leq s_N + (t_n - t_N) \leq s_N + T$ where $T = \sum_{k=N+1}^{\infty} b_k$. Thus $\{s_n\}$ is bounded above and increasing, hence convergent.

(b) This is the contrapositive of (a). \qed

\begin{theorem}{Limit Comparison Test}
Let $\sum a_n$ and $\sum b_n$ be series with positive terms. If $\lim_{n \to \infty} \frac{a_n}{b_n} = c$ where $0 < c < \infty$, then both series converge or both diverge.
\end{theorem}

\textit{Proof.} Since $\lim_{n \to \infty} \frac{a_n}{b_n} = c$, there exists $N$ such that for $n \geq N$:
\[
\frac{c}{2} < \frac{a_n}{b_n} < \frac{3c}{2}
\]
Thus $\frac{c}{2}b_n < a_n < \frac{3c}{2}b_n$ for $n \geq N$. By the comparison test, $\sum a_n$ converges iff $\sum b_n$ converges. \qed

\begin{theorem}{Integral Test}
Let $f$ be a positive, continuous, decreasing function on $[1, \infty)$ with $f(n) = a_n$. Then $\sum_{n=1}^{\infty} a_n$ converges if and only if $\int_1^{\infty} f(x)\, dx$ converges.
\end{theorem}

\textit{Proof.} Since $f$ is decreasing, for $k \geq 2$: $f(k) \leq \int_{k-1}^{k} f(x)\, dx \leq f(k-1)$. Summing from $k = 2$ to $n$:
\[
\sum_{k=2}^{n} f(k) \leq \int_1^{n} f(x)\, dx \leq \sum_{k=1}^{n-1} f(k)
\]
Thus the partial sums $\sum_{k=1}^{n} a_k$ are bounded iff the integrals $\int_1^{n} f(x)\, dx$ are bounded. By the monotone convergence theorem, both converge or both diverge. \qed

\begin{theorem}{Ratio Test}
Let $\sum a_n$ be a series with positive terms and suppose:
\[
\lim_{n \to \infty} \frac{a_{n+1}}{a_n} = L
\]
\begin{enumerate}
    \item[(a)] If $L < 1$, the series converges.
    \item[(b)] If $L > 1$, the series diverges.
    \item[(c)] If $L = 1$, the test is inconclusive.
\end{enumerate}
\end{theorem}

\textit{Proof.} (a) If $L < 1$, choose $r$ with $L < r < 1$. There exists $N$ such that $\frac{a_{n+1}}{a_n} < r$ for $n \geq N$. Thus $a_{N+k} < a_N r^k$ for $k \geq 0$. The geometric series $\sum r^k$ converges, so by comparison, $\sum a_n$ converges.

(b) If $L > 1$, then $a_{n+1} > a_n$ for $n \geq N$, so $a_n$ does not approach 0, and the series diverges. \qed

\begin{theorem}{Root Test}
Let $\sum a_n$ be a series with nonnegative terms and suppose:
\[
\lim_{n \to \infty} \sqrt[n]{a_n} = L
\]
\begin{enumerate}
    \item[(a)] If $L < 1$, the series converges.
    \item[(b)] If $L > 1$, the series diverges.
    \item[(c)] If $L = 1$, the test is inconclusive.
\end{enumerate}
\end{theorem}

\textit{Proof.} (a) If $L < 1$, choose $r$ with $L < r < 1$. There exists $N$ such that $\sqrt[n]{a_n} < r$ for $n \geq N$. Thus $a_n < r^n$ for $n \geq N$, and comparison with the geometric series gives convergence.

(b) If $L > 1$, then $a_n > 1$ for infinitely many $n$, so $a_n$ does not approach 0. \qed

\begin{theorem}{Alternating Series Test (Leibniz Criterion)}
If $\{a_n\}$ is a decreasing sequence of positive numbers with $\lim_{n \to \infty} a_n = 0$, then the alternating series $\sum_{n=1}^{\infty} (-1)^{n+1} a_n$ converges.
\end{theorem}

\textit{Proof.} Let $s_n$ be the $n$th partial sum. Consider the even-indexed partial sums $s_{2n}$:
\[
s_{2n} = (a_1 - a_2) + (a_3 - a_4) + \cdots + (a_{2n-1} - a_{2n})
\]
Since $\{a_n\}$ is decreasing, each term in parentheses is nonnegative, so $\{s_{2n}\}$ is increasing. Also:
\[
s_{2n} = a_1 - (a_2 - a_3) - \cdots - (a_{2n-2} - a_{2n-1}) - a_{2n} \leq a_1
\]
Thus $\{s_{2n}\}$ converges. Similarly, $\{s_{2n+1}\}$ converges to the same limit since $s_{2n+1} = s_{2n} + a_{2n+1} \to L + 0 = L$. \qed

\begin{definition}{Absolute Convergence}
A series $\sum a_n$ is said to converge absolutely if $\sum |a_n|$ converges. A series that converges but does not converge absolutely is said to converge conditionally.
\end{definition}

\begin{theorem}{Absolute Convergence Implies Convergence}
If $\sum_{n=1}^{\infty} |a_n|$ converges, then $\sum_{n=1}^{\infty} a_n$ converges.
\end{theorem}

\textit{Proof.} Define $a_n^+ = \max(a_n, 0)$ and $a_n^- = \max(-a_n, 0)$. Then $0 \leq a_n^+ \leq |a_n|$ and $0 \leq a_n^- \leq |a_n|$. By comparison, both $\sum a_n^+$ and $\sum a_n^-$ converge. Since $a_n = a_n^+ - a_n^-$, we have $\sum a_n = \sum a_n^+ - \sum a_n^-$ converges. \qed

\section{Improper Integrals}

\begin{definition}{Improper Integral of the First Kind}
For a function $f$ continuous on $[a, \infty)$, the improper integral is defined as:
\[
\int_a^{\infty} f(x)\, dx = \lim_{b \to \infty} \int_a^b f(x)\, dx
\]
If the limit exists and is finite, the integral is said to converge; otherwise, it diverges.
\end{definition}

\begin{definition}{Improper Integral of the Second Kind}
If $f$ is continuous on $(a, b]$ but unbounded near $a$, define:
\[
\int_a^b f(x)\, dx = \lim_{c \to a^+} \int_c^b f(x)\, dx
\]
\end{definition}

\begin{theorem}{Comparison Test for Improper Integrals}
Suppose $0 \leq f(x) \leq g(x)$ for all $x \geq a$.
\begin{enumerate}
    \item[(a)] If $\int_a^{\infty} g(x)\, dx$ converges, then $\int_a^{\infty} f(x)\, dx$ converges.
    \item[(b)] If $\int_a^{\infty} f(x)\, dx$ diverges, then $\int_a^{\infty} g(x)\, dx$ diverges.
\end{enumerate}
\end{theorem}

\textit{Proof.} (a) Let $F(b) = \int_a^b f(x)\, dx$ and $G(b) = \int_a^b g(x)\, dx$. Since $f \leq g$, we have $F(b) \leq G(b)$. If $\int_a^{\infty} g$ converges, $G(b)$ is bounded above, so $F(b)$ is bounded above and increasing, hence converges. \qed

\begin{theorem}{$p$-Test for Improper Integrals}
The improper integral:
\[
\int_1^{\infty} \frac{1}{x^p}\, dx
\]
converges if $p > 1$ and diverges if $p \leq 1$.
\end{theorem}

\textit{Proof.} For $p \neq 1$:
\[
\int_1^{b} \frac{1}{x^p}\, dx = \left[\frac{x^{1-p}}{1-p}\right]_1^b = \frac{b^{1-p} - 1}{1-p}
\]
As $b \to \infty$: if $p > 1$, then $1 - p < 0$, so $b^{1-p} \to 0$, giving convergence to $\frac{1}{p-1}$. If $p < 1$, then $b^{1-p} \to \infty$, giving divergence.

For $p = 1$: $\int_1^b \frac{1}{x}\, dx = \log b \to \infty$ as $b \to \infty$. \qed



\chapter{Sequences and Series of Functions}

\section{Pointwise Convergence}

\begin{definition}{Pointwise Convergence}
A sequence of functions $\{f_n\}$ defined on a set $S$ is said to converge pointwise to a function $f$ on $S$ if for each $x \in S$:
\[
\lim_{n \to \infty} f_n(x) = f(x)
\]
That is, for each $x \in S$ and each $\epsilon > 0$, there exists $N$ (depending on $x$ and $\epsilon$) such that $|f_n(x) - f(x)| < \epsilon$ whenever $n \geq N$.
\end{definition}

\begin{definition}{Pointwise Convergence of Series}
A series of functions $\sum_{n=1}^{\infty} f_n(x)$ converges pointwise to $f(x)$ on $S$ if the sequence of partial sums $s_n(x) = \sum_{k=1}^{n} f_k(x)$ converges pointwise to $f(x)$ on $S$.
\end{definition}

\section{Uniform Convergence}

\begin{definition}{Uniform Convergence}
A sequence of functions $\{f_n\}$ converges uniformly to $f$ on a set $S$ if for every $\epsilon > 0$, there exists $N$ (depending only on $\epsilon$) such that:
\[
|f_n(x) - f(x)| < \epsilon \quad \text{for all } n \geq N \text{ and all } x \in S
\]
\end{definition}

\begin{theorem}{Uniform Convergence and Continuity}
If $\{f_n\}$ is a sequence of continuous functions on $[a, b]$ that converges uniformly to $f$ on $[a, b]$, then $f$ is continuous on $[a, b]$.
\end{theorem}

\textit{Proof.} Given $\epsilon > 0$, choose $N$ such that $|f_N(x) - f(x)| < \epsilon/3$ for all $x \in [a, b]$. Since $f_N$ is continuous at any $c \in [a, b]$, there exists $\delta > 0$ such that $|f_N(x) - f_N(c)| < \epsilon/3$ whenever $|x - c| < \delta$. Then:
\begin{align*}
|f(x) - f(c)| &\leq |f(x) - f_N(x)| + |f_N(x) - f_N(c)| + |f_N(c) - f(c)| \\
&< \epsilon/3 + \epsilon/3 + \epsilon/3 = \epsilon
\end{align*}
for $|x - c| < \delta$. Thus $f$ is continuous at $c$. \qed

\begin{theorem}{Uniform Convergence and Integration}
If $\{f_n\}$ is a sequence of continuous functions on $[a, b]$ that converges uniformly to $f$ on $[a, b]$, then:
\[
\lim_{n \to \infty} \int_a^b f_n(x)\, dx = \int_a^b f(x)\, dx
\]
Equivalently:
\[
\lim_{n \to \infty} \int_a^b f_n(x)\, dx = \int_a^b \lim_{n \to \infty} f_n(x)\, dx
\]
\end{theorem}

\textit{Proof.} Given $\epsilon > 0$, choose $N$ such that $|f_n(x) - f(x)| < \epsilon/(b-a)$ for all $n \geq N$ and $x \in [a, b]$. Then:
\[
\left|\int_a^b f_n(x)\, dx - \int_a^b f(x)\, dx\right| \leq \int_a^b |f_n(x) - f(x)|\, dx < \int_a^b \frac{\epsilon}{b-a}\, dx = \epsilon
\]
for $n \geq N$. Thus the limit holds. \qed

\begin{theorem}{Weierstrass M-Test}
Let $\{f_n\}$ be a sequence of functions defined on $S$, and let $\{M_n\}$ be a sequence of nonnegative constants such that $|f_n(x)| \leq M_n$ for all $x \in S$ and all $n$. If $\sum_{n=1}^{\infty} M_n$ converges, then $\sum_{n=1}^{\infty} f_n(x)$ converges uniformly on $S$.
\end{theorem}

\textit{Proof.} Since $\sum M_n$ converges, by the Cauchy criterion, for $\epsilon > 0$ there exists $N$ such that $\sum_{k=n+1}^{m} M_k < \epsilon$ for $m > n \geq N$. Then for all $x \in S$:
\[
\left|\sum_{k=n+1}^{m} f_k(x)\right| \leq \sum_{k=n+1}^{m} |f_k(x)| \leq \sum_{k=n+1}^{m} M_k < \epsilon
\]
Thus the series satisfies the uniform Cauchy criterion and converges uniformly. \qed

\section{Power Series}

\begin{definition}{Power Series}
A power series is a series of the form:
\[
\sum_{n=0}^{\infty} a_n(x - a)^n = a_0 + a_1(x-a) + a_2(x-a)^2 + \cdots
\]
where $a$ is the center and $\{a_n\}$ are the coefficients.
\end{definition}

\begin{theorem}{Radius of Convergence}
For any power series $\sum_{n=0}^{\infty} a_n(x - a)^n$, there exists a number $R$ ($0 \leq R \leq \infty$), called the radius of convergence, such that:
\begin{enumerate}
    \item[(a)] The series converges absolutely for $|x - a| < R$
    \item[(b)] The series diverges for $|x - a| > R$
    \item[(c)] The series converges uniformly on any interval $[a-r, a+r]$ where $0 < r < R$
\end{enumerate}
\end{theorem}

\textit{Proof.} Let $R = \sup\{|x-a| : \sum |a_n(x-a)^n| \text{ converges}\}$. If $|x-a| < R$, there exists $y$ with $|x-a| < |y-a| \leq R$ where $\sum |a_n(y-a)^n|$ converges. Since $|a_n(x-a)^n| = |a_n(y-a)^n| \cdot |\frac{x-a}{y-a}|^n$ and $|\frac{x-a}{y-a}| < 1$, comparison with the geometric series gives absolute convergence.

If $|x-a| > R$, then $\sum |a_n(x-a)^n|$ diverges by definition of $R$, so the series diverges.

For uniform convergence on $[a-r, a+r]$ with $r < R$, choose $y$ with $r < |y-a| < R$. Then $\sum |a_n(y-a)^n|$ converges, and for $|x-a| \leq r$: $|a_n(x-a)^n| \leq |a_n| r^n \leq |a_n(y-a)^n|$. By the Weierstrass M-test with $M_n = |a_n(y-a)^n|$, the series converges uniformly. \qed

\begin{theorem}{Ratio Test for Radius of Convergence}
If $\lim_{n \to \infty} |\frac{a_{n+1}}{a_n}| = L$ (or $\limsup_{n \to \infty} |\frac{a_{n+1}}{a_n}| = L$), then the radius of convergence is:
\[
R = \begin{cases} \frac{1}{L} & \text{if } 0 < L < \infty \\ \infty & \text{if } L = 0 \\ 0 & \text{if } L = \infty \end{cases}
\]
\end{theorem}

\textit{Proof.} Apply the ratio test to $\sum |a_n(x-a)^n|$:
\[
\lim_{n \to \infty} \left|\frac{a_{n+1}(x-a)^{n+1}}{a_n(x-a)^n}\right| = L|x-a|
\]
The series converges when $L|x-a| < 1$, i.e., $|x-a| < 1/L$. \qed

\section{Properties of Power Series}

\begin{theorem}{Continuity of Power Series}
A power series $f(x) = \sum_{n=0}^{\infty} a_n(x-a)^n$ is continuous on its interval of convergence.
\end{theorem}

\textit{Proof.} On any $[a-r, a+r]$ with $r < R$, the series converges uniformly by the previous theorem. Since each partial sum is continuous, the limit function $f$ is continuous. Since this holds for any $r < R$, $f$ is continuous on $(a-R, a+R)$. \qed

\begin{theorem}{Integration of Power Series}
If $f(x) = \sum_{n=0}^{\infty} a_n(x-a)^n$ has radius of convergence $R > 0$, then for $|x-a| < R$:
\[
\int_a^x f(t)\, dt = \sum_{n=0}^{\infty} \frac{a_n}{n+1}(x-a)^{n+1}
\]
The integrated series has the same radius of convergence $R$.
\end{theorem}

\textit{Proof.} On $[a-r, a+r]$ with $r < R$, the series converges uniformly, so:
\[
\int_a^x f(t)\, dt = \sum_{n=0}^{\infty} \int_a^x a_n(t-a)^n\, dt = \sum_{n=0}^{\infty} \frac{a_n}{n+1}(x-a)^{n+1}
\]
The radius of convergence is the same since $\limsup \sqrt[n]{|a_n|/(n+1)} = \limsup \sqrt[n]{|a_n|}$. \qed

\begin{theorem}{Differentiation of Power Series}
If $f(x) = \sum_{n=0}^{\infty} a_n(x-a)^n$ has radius of convergence $R > 0$, then $f$ is differentiable on $(a-R, a+R)$ and:
\[
f'(x) = \sum_{n=1}^{\infty} na_n(x-a)^{n-1}
\]
The differentiated series has the same radius of convergence $R$.
\end{theorem}

\textit{Proof.} For $x \in (a-R, a+R)$, choose $r$ with $|x-a| < r < R$. The differentiated series $\sum na_n(x-a)^{n-1}$ converges uniformly on $[a-r, a+r]$ by comparison with $\sum n|a_n| r^{n-1}$ which converges by the ratio test. Thus term-by-term differentiation is valid. The radius is the same since $\limsup \sqrt[n]{n|a_n|} = \limsup \sqrt[n]{|a_n|}$. \qed

\section{Taylor Series}

\begin{definition}{Taylor Series}
If $f$ has derivatives of all orders at $a$, the Taylor series of $f$ at $a$ is:
\[
\sum_{n=0}^{\infty} \frac{f^{(n)}(a)}{n!}(x-a)^n
\]
When $a = 0$, this is called the Maclaurin series.
\end{definition}

\begin{theorem}{Taylor Series Convergence}
Let $f$ have derivatives of all orders on $(a-R, a+R)$. If there exists $M$ such that $|f^{(n)}(x)| \leq M$ for all $n$ and all $x \in (a-R, a+R)$, then the Taylor series converges to $f(x)$ for all $x \in (a-R, a+R)$.
\end{theorem}

\textit{Proof.} By Taylor's theorem with remainder:
\[
f(x) = \sum_{k=0}^{n} \frac{f^{(k)}(a)}{k!}(x-a)^k + R_n(x)
\]
where $R_n(x) = \frac{f^{(n+1)}(c)}{(n+1)!}(x-a)^{n+1}$ for some $c$ between $a$ and $x$. Since $|f^{(n+1)}(c)| \leq M$:
\[
|R_n(x)| \leq \frac{M}{(n+1)!}|x-a|^{n+1} \to 0 \quad \text{as } n \to \infty
\]
for $|x-a| < R$. Thus the Taylor series converges to $f(x)$. \qed

\begin{theorem}{Uniqueness of Power Series Representation}
If $f(x) = \sum_{n=0}^{\infty} a_n(x-a)^n$ for $|x-a| < R$, then $a_n = \frac{f^{(n)}(a)}{n!}$ for all $n$.
\end{theorem}

\textit{Proof.} Setting $x = a$ gives $f(a) = a_0$. Differentiating term by term $n$ times:
\[
f^{(n)}(x) = \sum_{k=n}^{\infty} k(k-1)\cdots(k-n+1)a_k(x-a)^{k-n}
\]
Setting $x = a$, all terms with $k > n$ vanish, leaving $f^{(n)}(a) = n! a_n$. Thus $a_n = f^{(n)}(a)/n!$. \qed



% ====================
% Part V: Vector Algebra
% ====================
\part{Vector Algebra}

\chapter{Vector Algebra}

\section{The Vector Space of n-tuples of Real Numbers}

\begin{definition}{Vector}
For a positive integer $n$, an ordered $n$-tuple of real numbers $(a_1, a_2, \ldots, a_n)$ is called an $n$-dimensional vector or $n$-vector. The set of all $n$-vectors is denoted by $V_n$.
\end{definition}

\begin{definition}{Vector Operations}
For vectors $A = (a_1, \ldots, a_n)$ and $B = (b_1, \ldots, b_n)$ in $V_n$ and scalar $c \in \mathbb{R}$:
\begin{enumerate}
    \item[(a)] Addition: $A + B = (a_1 + b_1, a_2 + b_2, \ldots, a_n + b_n)$
    \item[(b)] Scalar multiplication: $cA = (ca_1, ca_2, \ldots, ca_n)$
    \item[(c)] Zero vector: $O = (0, 0, \ldots, 0)$
    \item[(d)] Negative: $-A = (-a_1, -a_2, \ldots, -a_n)$
\end{enumerate}
\end{definition}

\begin{theorem}{Properties of Vector Space}
For all vectors $A, B, C$ in $V_n$ and all scalars $c, d$:
\begin{enumerate}
    \item[(a)] $A + B = B + A$ (commutativity)
    \item[(b)] $A + (B + C) = (A + B) + C$ (associativity)
    \item[(c)] $A + O = A$ (additive identity)
    \item[(d)] $A + (-A) = O$ (additive inverse)
    \item[(e)] $c(A + B) = cA + cB$ (distributivity)
    \item[(f)] $(c + d)A = cA + dA$ (distributivity)
    \item[(g)] $(cd)A = c(dA)$ (associativity of scalar mult.)
    \item[(h)] $1A = A$ (scalar identity)
\end{enumerate}
\end{theorem}

\textit{Proof.} These follow directly from the corresponding properties of real numbers applied component-wise. \qed

\section{The Dot Product}

\begin{definition}{Dot Product}
For vectors $A = (a_1, \ldots, a_n)$ and $B = (b_1, \ldots, b_n)$ in $V_n$, the dot product is:
\[
A \cdot B = \sum_{k=1}^{n} a_kb_k = a_1b_1 + a_2b_2 + \cdots + a_nb_n
\]
\end{definition}

\begin{theorem}{Properties of Dot Product}
For all vectors $A, B, C$ in $V_n$ and scalar $c$:
\begin{enumerate}
    \item[(a)] $A \cdot B = B \cdot A$ (symmetry)
    \item[(b)] $A \cdot (B + C) = A \cdot B + A \cdot C$ (linearity)
    \item[(c)] $c(A \cdot B) = (cA) \cdot B = A \cdot (cB)$ (homogeneity)
    \item[(d)] $A \cdot A > 0$ if $A \neq O$, and $A \cdot A = 0$ if $A = O$ (positive definiteness)
\end{enumerate}
\end{theorem}

\textit{Proof.} All follow from the definition and properties of real numbers. For (d): $A \cdot A = \sum a_k^2 \geq 0$ with equality iff all $a_k = 0$. \qed

\begin{definition}{Norm and Length}
The norm (length) of vector $A$ is $\|A\| = \sqrt{A \cdot A} = \sqrt{\sum_{k=1}^{n} a_k^2}$.
\end{definition}

\begin{theorem}{Cauchy-Schwarz Inequality}
For all vectors $A, B$ in $V_n$:
\[
|A \cdot B| \leq \|A\| \|B\|
\]
with equality if and only if one vector is a scalar multiple of the other.
\end{theorem}

\textit{Proof.} If $B = O$, both sides are 0. Otherwise, for any real $t$:
\[
\|A - tB\|^2 = (A - tB) \cdot (A - tB) = \|A\|^2 - 2t(A \cdot B) + t^2\|B\|^2 \geq 0
\]
Setting $t = \frac{A \cdot B}{\|B\|^2}$:
\[
\|A\|^2 - \frac{(A \cdot B)^2}{\|B\|^2} \geq 0
\]
Thus $(A \cdot B)^2 \leq \|A\|^2 \|B\|^2$, giving $|A \cdot B| \leq \|A\| \|B\|$. Equality holds iff $A = tB$. \qed

\begin{theorem}{Triangle Inequality}
For all vectors $A, B$ in $V_n$:
\[
\|A + B\| \leq \|A\| + \|B\|
\]
\end{theorem}

\textit{Proof.} Using Cauchy-Schwarz:
\begin{align*}
\|A + B\|^2 &= (A + B) \cdot (A + B) = \|A\|^2 + 2(A \cdot B) + \|B\|^2 \\
&\leq \|A\|^2 + 2|A \cdot B| + \|B\|^2 \\
&\leq \|A\|^2 + 2\|A\|\|B\| + \|B\|^2 = (\|A\| + \|B\|)^2
\end{align*}
Taking square roots gives the result. \qed

\section{Orthogonality}

\begin{definition}{Orthogonal Vectors}
Two vectors $A$ and $B$ are orthogonal (perpendicular) if $A \cdot B = 0$.
\end{definition}

\begin{theorem}{Pythagorean Theorem}
Vectors $A$ and $B$ are orthogonal if and only if:
\[
\|A + B\|^2 = \|A\|^2 + \|B\|^2
\]
\end{theorem}

\textit{Proof.} $\|A + B\|^2 = \|A\|^2 + 2(A \cdot B) + \|B\|^2 = \|A\|^2 + \|B\|^2$ if and only if $A \cdot B = 0$. \qed

\section{Linear Independence and Bases}

\begin{definition}{Linear Combination}
A vector $X$ is a linear combination of vectors $A_1, A_2, \ldots, A_k$ if:
\[
X = c_1A_1 + c_2A_2 + \cdots + c_kA_k
\]
for some scalars $c_1, c_2, \ldots, c_k$.
\end{definition}

\begin{definition}{Linear Independence}
Vectors $A_1, A_2, \ldots, A_k$ are linearly independent if:
\[
c_1A_1 + c_2A_2 + \cdots + c_kA_k = O \implies c_1 = c_2 = \cdots = c_k = 0
\]
Otherwise, they are linearly dependent.
\end{definition}

\begin{definition}{Basis}
A set of vectors $\{E_1, E_2, \ldots, E_n\}$ is a basis for $V_n$ if:
\begin{enumerate}
    \item[(a)] They are linearly independent
    \item[(b)] Every vector in $V_n$ can be written as a linear combination of $E_1, \ldots, E_n$
\end{enumerate}
The standard basis consists of $E_k = (0, \ldots, 1, \ldots, 0)$ with 1 in the $k$th position.
\end{definition}

\begin{theorem}{Unique Representation}
If $\{E_1, \ldots, E_n\}$ is a basis for $V_n$, then every vector $X \in V_n$ can be written uniquely as:
\[
X = x_1E_1 + x_2E_2 + \cdots + x_nE_n
\]
\end{theorem}

\textit{Proof.} Since $\{E_i\}$ spans $V_n$, such a representation exists. If $X = \sum x_iE_i = \sum y_iE_i$, then $\sum (x_i - y_i)E_i = O$. By linear independence, $x_i - y_i = 0$ for all $i$, so the representation is unique. \qed

\section{Complex Vector Space $V_n(\mathbb{C})$}

\begin{definition}{Complex Vector}
The complex vector space $V_n(\mathbb{C})$ is the set of all ordered $n$-tuples of complex numbers $(a_1, a_2, \ldots, a_n)$ where each $a_k \in \mathbb{C}$.
\end{definition}

\begin{definition}{Hermitian Dot Product}
For vectors $A = (a_1, \ldots, a_n)$ and $B = (b_1, \ldots, b_n)$ in $V_n(\mathbb{C})$, the Hermitian dot product is:
\[
A \cdot B = \sum_{k=1}^{n} a_k\overline{b_k} = a_1\overline{b_1} + a_2\overline{b_2} + \cdots + a_n\overline{b_n}
\]
where $\overline{b_k}$ denotes the complex conjugate of $b_k$.
\end{definition}

\begin{theorem}{Properties of Hermitian Dot Product}
For all vectors $A, B, C$ in $V_n(\mathbb{C})$ and complex scalar $c$:
\begin{enumerate}
    \item[(a)] $A \cdot B = \overline{B \cdot A}$ (conjugate symmetry)
    \item[(b)] $A \cdot (B + C) = A \cdot B + A \cdot C$ (linearity in first argument)
    \item[(c)] $c(A \cdot B) = (cA) \cdot B$ and $A \cdot (cB) = \overline{c}(A \cdot B)$ (homogeneity)
    \item[(d)] $A \cdot A > 0$ if $A \neq O$, and $A \cdot A = 0$ if $A = O$ (positive definiteness)
\end{enumerate}
\end{theorem}

\textit{Proof.} (a) $A \cdot B = \sum a_k\overline{b_k} = \sum \overline{\overline{a_k}b_k} = \overline{\sum \overline{a_k}b_k} = \overline{B \cdot A}$.

(c) $(cA) \cdot B = \sum (ca_k)\overline{b_k} = c\sum a_k\overline{b_k} = c(A \cdot B)$. And $A \cdot (cB) = \sum a_k\overline{cb_k} = \sum a_k\overline{c}\overline{b_k} = \overline{c}(A \cdot B)$.

(d) $A \cdot A = \sum a_k\overline{a_k} = \sum |a_k|^2 \geq 0$ with equality iff all $a_k = 0$. \qed

\begin{definition}{Norm in Complex Space}
The norm of $A \in V_n(\mathbb{C})$ is $\|A\| = \sqrt{A \cdot A} = \sqrt{\sum_{k=1}^{n} |a_k|^2}$.
\end{definition}

\begin{theorem}{Cauchy-Schwarz Inequality for Complex Vectors}
For all vectors $A, B$ in $V_n(\mathbb{C})$:
\[
|A \cdot B| \leq \|A\| \|B\|
\]
\end{theorem}

\textit{Proof.} For any complex scalar $t$, $\|A - tB\|^2 = (A - tB) \cdot (A - tB) \geq 0$. Expanding:
\[
\|A\|^2 - t(A \cdot B) - \overline{t}(B \cdot A) + |t|^2\|B\|^2 \geq 0
\]
Setting $t = \frac{A \cdot B}{\|B\|^2}$ (assuming $B \neq O$) and using $A \cdot B = \overline{B \cdot A}$:
\[
\|A\|^2 - \frac{|A \cdot B|^2}{\|B\|^2} \geq 0
\]
Thus $|A \cdot B|^2 \leq \|A\|^2 \|B\|^2$. \qed



\chapter{Applications to Analytic Geometry}

\section{Lines}

\begin{definition}{Parametric Equation of a Line}
A line $L$ through point $P$ in the direction of nonzero vector $A$ consists of all points $X$ such that:
\[
X = P + tA
\]
for some real number $t$. In coordinates, if $P = (p_1, p_2, p_3)$ and $A = (a_1, a_2, a_3)$, then:
\[
x_1 = p_1 + ta_1, \quad x_2 = p_2 + ta_2, \quad x_3 = p_3 + ta_3
\]
\end{definition}

\begin{definition}{Linear Span of Two Points}
The line through two distinct points $P$ and $Q$ is given by:
\[
L(P, Q) = \{P + t(Q - P) : t \in \mathbb{R}\}
\]
\end{definition}

\begin{theorem}{Parallel Lines}
Two lines $L(P, A)$ and $L(Q, B)$ are parallel if and only if their direction vectors $A$ and $B$ are parallel (i.e., one is a scalar multiple of the other).
\end{theorem}

\textit{Proof.} The lines have the same direction if and only if $A = cB$ for some scalar $c \neq 0$. \qed

\section{Planes}

\begin{definition}{Plane Through a Point with Normal Vector}
A plane $M$ through point $P$ with nonzero normal vector $N$ consists of all points $X$ satisfying:
\[
(X - P) \cdot N = 0
\]
or equivalently $X \cdot N = P \cdot N$.
\end{definition}

\begin{theorem}{Plane Through Three Points}
Given three non-collinear points $P, Q, R$, there is exactly one plane containing all three. If $A = Q - P$ and $B = R - P$, a normal vector is $N = A \times B$.
\end{theorem}

\textit{Proof.} Since $P, Q, R$ are non-collinear, $A$ and $B$ are linearly independent, so $N = A \times B \neq O$. The plane through $P$ with normal $N$ contains $Q$ and $R$ since $(Q-P) \cdot N = A \cdot (A \times B) = 0$ and similarly for $R$. \qed

\begin{theorem}{Distance from Point to Plane}
The distance from point $Q$ to the plane $(X - P) \cdot N = 0$ is:
\[
d = \frac{|(Q - P) \cdot N|}{\|N\|}
\]
\end{theorem}

\textit{Proof.} The distance is the length of the projection of $(Q - P)$ onto the normal direction:
\[
d = \left\|\frac{(Q - P) \cdot N}{\|N\|^2}N\right\| = \frac{|(Q - P) \cdot N|}{\|N\|}
\] \qed

\section{The Cross Product}

\begin{definition}{Cross Product}
For vectors $A = (a_1, a_2, a_3)$ and $B = (b_1, b_2, b_3)$ in $V_3$, the cross product is:
\[
A \times B = (a_2b_3 - a_3b_2, a_3b_1 - a_1b_3, a_1b_2 - a_2b_1)
\]
\end{definition}

\begin{theorem}{Properties of Cross Product}
For all vectors $A, B, C$ in $V_3$ and scalars $c$:
\begin{enumerate}
    \item[(a)] $A \times B = -(B \times A)$ (anticommutativity)
    \item[(b)] $A \times (B + C) = (A \times B) + (A \times C)$ (distributivity)
    \item[(c)] $(cA) \times B = c(A \times B) = A \times (cB)$ (homogeneity)
    \item[(d)] $A \times A = O$
    \item[(e)] $A \cdot (A \times B) = 0$ and $B \cdot (A \times B) = 0$ (orthogonality)
\end{enumerate}
\end{theorem}

\textit{Proof.} (a)-(d) follow directly from the definition by component-wise verification.

(e) $A \cdot (A \times B) = a_1(a_2b_3 - a_3b_2) + a_2(a_3b_1 - a_1b_3) + a_3(a_1b_2 - a_2b_1) = 0$ after cancellation. Similarly for $B$. \qed

\begin{theorem}{Geometric Interpretation of Cross Product}
For vectors $A, B$ in $V_3$:
\[
\|A \times B\| = \|A\| \|B\| \sin\theta
\]
where $\theta$ is the angle between $A$ and $B$. The vector $A \times B$ is orthogonal to both $A$ and $B$, with direction given by the right-hand rule.
\end{theorem}

\textit{Proof.} Using $\|A \times B\|^2 + (A \cdot B)^2 = \|A\|^2\|B\|^2$ (Lagrange's identity) and $A \cdot B = \|A\|\|B\|\cos\theta$:
\[
\|A \times B\|^2 = \|A\|^2\|B\|^2(1 - \cos^2\theta) = \|A\|^2\|B\|^2\sin^2\theta
\]
Taking square roots gives the result. \qed

\begin{theorem}{Scalar Triple Product}
For vectors $A, B, C$ in $V_3$:
\[
A \cdot (B \times C) = \det\begin{pmatrix} a_1 & a_2 & a_3 \\ b_1 & b_2 & b_3 \\ c_1 & c_2 & c_3 \end{pmatrix}
\]
The absolute value $|A \cdot (B \times C)|$ equals the volume of the parallelepiped determined by $A, B, C$.
\end{theorem}

\textit{Proof.} Expanding both sides gives the same expression. Geometrically, $\|B \times C\|$ is the area of the base parallelogram, and $|A \cdot (B \times C)|/\|B \times C\| = |A||\cos\phi|$ is the height, where $\phi$ is the angle between $A$ and $B \times C$. \qed

\section{Conic Sections}

\begin{definition}{Focus-Directrix Definition}
A conic section is the set of all points $P$ such that the distance from $P$ to a fixed point $F$ (the focus) divided by the distance from $P$ to a fixed line $D$ (the directrix) equals a constant $e$ (the eccentricity):
\[
\frac{\|P - F\|}{d(P, D)} = e
\]
If $e < 1$: ellipse; if $e = 1$: parabola; if $e > 1$: hyperbola.
\end{definition}

\begin{theorem}{Standard Equation of Ellipse}
An ellipse with foci at $(c, 0)$ and $(-c, 0)$ and sum of distances $2a$ ($a > c$) has equation:
\[
\frac{x^2}{a^2} + \frac{y^2}{b^2} = 1
\]
where $b^2 = a^2 - c^2$. The eccentricity is $e = c/a$.
\end{theorem}

\textit{Proof.} By definition, $\sqrt{(x-c)^2 + y^2} + \sqrt{(x+c)^2 + y^2} = 2a$. Moving one radical and squaring:
\[
(x-c)^2 + y^2 = 4a^2 - 4a\sqrt{(x+c)^2 + y^2} + (x+c)^2 + y^2
\]
Simplifying: $-xc = a^2 - a\sqrt{(x+c)^2 + y^2}$. Rearranging and squaring again gives the result. \qed

\begin{theorem}{Standard Equation of Parabola}
A parabola with focus at $(p, 0)$ and directrix $x = -p$ has equation:
\[
y^2 = 4px
\]
\end{theorem}

\textit{Proof.} By definition: $\sqrt{(x-p)^2 + y^2} = |x + p|$. Squaring: $(x-p)^2 + y^2 = (x+p)^2$. Expanding: $x^2 - 2px + p^2 + y^2 = x^2 + 2px + p^2$. Thus $y^2 = 4px$. \qed

\begin{theorem}{Standard Equation of Hyperbola}
A hyperbola with foci at $(c, 0)$ and $(-c, 0)$ and difference of distances $2a$ ($c > a$) has equation:
\[
\frac{x^2}{a^2} - \frac{y^2}{b^2} = 1
\]
where $b^2 = c^2 - a^2$. The eccentricity is $e = c/a > 1$.
\end{theorem}

\textit{Proof.} Similar to the ellipse proof, using $|\sqrt{(x-c)^2 + y^2} - \sqrt{(x+c)^2 + y^2}| = 2a$. \qed

\begin{theorem}{Polar Equation of Conic}
With focus at the origin and directrix $x = d$, a conic with eccentricity $e$ has polar equation:
\[
r = \frac{ed}{1 + e\cos\theta}
\]
\end{theorem}

\textit{Proof.} If $P$ has polar coordinates $(r, \theta)$, then $|PF| = r$ and the distance to the directrix $x = d$ is $|d - r\cos\theta|$. By the focus-directrix property: $r = e|d - r\cos\theta|$. For $ed/(1+e\cos\theta) > 0$, we get $r = ed/(1+e\cos\theta)$. \qed



\chapter{Calculus of Vector-Valued Functions}

\section{Vector-Valued Functions}

\begin{definition}{Vector-Valued Function}
A vector-valued function $\mathbf{r}$ defined on a set $S$ assigns to each $t \in S$ a unique vector $\mathbf{r}(t)$ in $V_n$. If $n = 3$, we can write:
\[
\mathbf{r}(t) = f(t)\mathbf{i} + g(t)\mathbf{j} + h(t)\mathbf{k} = (f(t), g(t), h(t))
\]
where $f, g, h$ are real-valued functions called the components of $\mathbf{r}$.
\end{definition}

\begin{definition}{Limit of a Vector-Valued Function}
If $\mathbf{r}(t) = (f_1(t), f_2(t), \ldots, f_n(t))$, then:
\[
\lim_{t \to a} \mathbf{r}(t) = \left(\lim_{t \to a} f_1(t), \lim_{t \to a} f_2(t), \ldots, \lim_{t \to a} f_n(t)\right)
\]
provided all component limits exist.
\end{definition}

\begin{definition}{Continuity}
A vector-valued function $\mathbf{r}$ is continuous at $a$ if $\lim_{t \to a} \mathbf{r}(t) = \mathbf{r}(a)$.
\end{definition}

\section{Derivatives and Integrals}

\begin{definition}{Derivative}
The derivative of a vector-valued function $\mathbf{r}$ is defined as:
\[
\mathbf{r}'(t) = \lim_{h \to 0} \frac{\mathbf{r}(t+h) - \mathbf{r}(t)}{h}
\]
In components, if $\mathbf{r}(t) = (f_1(t), \ldots, f_n(t))$, then $\mathbf{r}'(t) = (f_1'(t), \ldots, f_n'(t))$.
\end{definition}

\begin{theorem}{Differentiation Rules}
For vector-valued functions $\mathbf{u}, \mathbf{v}$ and scalar function $f$:
\begin{enumerate}
    \item[(a)] $(\mathbf{u} + \mathbf{v})' = \mathbf{u}' + \mathbf{v}'$
    \item[(b)] $(f\mathbf{u})' = f'\mathbf{u} + f\mathbf{u}'$
    \item[(c)] $(\mathbf{u} \cdot \mathbf{v})' = \mathbf{u}' \cdot \mathbf{v} + \mathbf{u} \cdot \mathbf{v}'$
    \item[(d)] $(\mathbf{u} \times \mathbf{v})' = \mathbf{u}' \times \mathbf{v} + \mathbf{u} \times \mathbf{v}'$ (in $V_3$)
\end{enumerate}
\end{theorem}

\textit{Proof.} All follow from the component-wise definition and the corresponding rules for scalar functions. For (c):
\[
(\mathbf{u} \cdot \mathbf{v})' = \sum (u_iv_i)' = \sum (u_i'v_i + u_iv_i') = \mathbf{u}' \cdot \mathbf{v} + \mathbf{u} \cdot \mathbf{v}'
\] \qed

\begin{theorem}{Constant Length}
If $\mathbf{r}(t)$ has constant length ($\|\mathbf{r}(t)\|$ is constant), then $\mathbf{r}(t) \cdot \mathbf{r}'(t) = 0$, i.e., $\mathbf{r}(t)$ is orthogonal to $\mathbf{r}'(t)$.
\end{theorem}

\textit{Proof.} If $\|\mathbf{r}(t)\|^2 = \mathbf{r}(t) \cdot \mathbf{r}(t)$ is constant, then:
\[
0 = \frac{d}{dt}(\mathbf{r} \cdot \mathbf{r}) = \mathbf{r}' \cdot \mathbf{r} + \mathbf{r} \cdot \mathbf{r}' = 2\mathbf{r} \cdot \mathbf{r}'
\]
Thus $\mathbf{r} \cdot \mathbf{r}' = 0$. \qed

\begin{definition}{Integral}
The integral of $\mathbf{r}(t) = (f_1(t), \ldots, f_n(t))$ is:
\[
\int_a^b \mathbf{r}(t)\, dt = \left(\int_a^b f_1(t)\, dt, \ldots, \int_a^b f_n(t)\, dt\right)
\]
\end{definition}

\section{Applications to Curves}

\begin{definition}{Tangent Vector}
For a curve defined by $\mathbf{r}(t)$, the tangent vector at $t$ is $\mathbf{r}'(t)$. The unit tangent vector is:
\[
\mathbf{T}(t) = \frac{\mathbf{r}'(t)}{\|\mathbf{r}'(t)\|}
\]
\end{definition}

\begin{theorem}{Geometric Interpretation}
The vector $\mathbf{r}'(t)$ is tangent to the curve at the point $\mathbf{r}(t)$ and points in the direction of increasing $t$.
\end{theorem}

\textit{Proof.} The difference quotient $\frac{\mathbf{r}(t+h) - \mathbf{r}(t)}{h}$ lies along the secant line. As $h \to 0$, this approaches the tangent direction. \qed

\section{Arc Length}

\begin{definition}{Arc Length}
The arc length of a curve $\mathbf{r}(t)$ from $t = a$ to $t = b$ is:
\[
L = \int_a^b \|\mathbf{r}'(t)\|\, dt
\]
\end{definition}

\begin{theorem}{Arc Length is Independent of Parametrization}
The arc length of a curve depends only on the curve itself, not on the particular parametrization used.
\end{theorem}

\textit{Proof.} If $\mathbf{r}_1(u) = \mathbf{r}(t(u))$ where $t(u)$ is a reparametrization, then by the chain rule:
\[
\mathbf{r}_1'(u) = \mathbf{r}'(t(u))t'(u)
\]
So $\|\mathbf{r}_1'(u)\| = \|\mathbf{r}'(t(u))\||t'(u)|$, and:
\[
\int \|\mathbf{r}_1'(u)\|\, du = \int \|\mathbf{r}'(t)\|\, dt
\]
by substitution. \qed

\begin{definition}{Arc Length Function}
The arc length function $s(t)$ measures length from a fixed point:
\[
s(t) = \int_a^t \|\mathbf{r}'(u)\|\, du
\]
\end{definition}

\begin{theorem}{}
$\frac{ds}{dt} = \|\mathbf{r}'(t)\|$.
\end{theorem}

\textit{Proof.} By the fundamental theorem of calculus applied to each component. \qed

\section{Curvature}

\begin{definition}{Curvature}
The curvature $\kappa$ of a curve measures how fast the unit tangent vector changes direction:
\[
\kappa = \left\|\frac{d\mathbf{T}}{ds}\right\| = \frac{\|\mathbf{T}'(t)\|}{\|\mathbf{r}'(t)\|}
\]
\end{definition}

\begin{theorem}{Curvature Formula}
If $\mathbf{r}(t)$ is a smooth curve, then:
\[
\kappa = \frac{\|\mathbf{r}'(t) \times \mathbf{r}''(t)\|}{\|\mathbf{r}'(t)\|^3}
\]
\end{theorem}

\textit{Proof.} Since $\mathbf{T} = \mathbf{r}'/\|\mathbf{r}'\| = \mathbf{r}'/(ds/dt)$, we have $\mathbf{r}' = (ds/dt)\mathbf{T}$. Differentiating:
\[
\mathbf{r}'' = \frac{d^2s}{dt^2}\mathbf{T} + \frac{ds}{dt}\mathbf{T}'
\]
Since $\mathbf{T}'$ is orthogonal to $\mathbf{T}$ (constant length), taking the cross product:
\[
\mathbf{r}' \times \mathbf{r}'' = \left(\frac{ds}{dt}\right)^2 (\mathbf{T} \times \mathbf{T}')
\]
Since $\|\mathbf{T}\| = 1$ and $\mathbf{T} \perp \mathbf{T}'$: $\|\mathbf{T} \times \mathbf{T}'\| = \|\mathbf{T}'\|$. Thus:
\[
\|\mathbf{r}' \times \mathbf{r}''\| = \left(\frac{ds}{dt}\right)^2 \|\mathbf{T}'\| = \|\mathbf{r}'\|^2 \kappa \|\mathbf{r}'\| = \kappa \|\mathbf{r}'\|^3
\] \qed

\section{Polar Coordinates}

\begin{definition}{Polar Coordinates}
In polar coordinates, a point is described by $(r, \theta)$ where $r$ is the distance from the origin and $\theta$ is the angle from the positive x-axis. The relation to Cartesian coordinates is:
\[
x = r\cos\theta, \quad y = r\sin\theta
\]
\end{definition}

\begin{theorem}{Velocity in Polar Coordinates}
If $\mathbf{r}(t) = r(t)\mathbf{e}_r$ where $\mathbf{e}_r = (\cos\theta, \sin\theta)$ and $\mathbf{e}_\theta = (-\sin\theta, \cos\theta)$, then:
\[
\mathbf{v} = \frac{d\mathbf{r}}{dt} = \dot{r}\mathbf{e}_r + r\dot{\theta}\mathbf{e}_\theta
\]
where dots denote time derivatives.
\end{theorem}

\textit{Proof.} $\mathbf{r} = r\mathbf{e}_r$, so:
\[
\dot{\mathbf{r}} = \dot{r}\mathbf{e}_r + r\dot{\mathbf{e}}_r = \dot{r}\mathbf{e}_r + r\dot{\theta}\mathbf{e}_\theta
\]
since $\dot{\mathbf{e}}_r = \dot{\theta}(-\sin\theta, \cos\theta) = \dot{\theta}\mathbf{e}_\theta$. \qed

\begin{theorem}{Acceleration in Polar Coordinates}
The acceleration in polar coordinates is:
\[
\mathbf{a} = (\ddot{r} - r\dot{\theta}^2)\mathbf{e}_r + (r\ddot{\theta} + 2\dot{r}\dot{\theta})\mathbf{e}_\theta
\]
\end{theorem}

\textit{Proof.} Differentiating the velocity formula and using $\dot{\mathbf{e}}_\theta = -\dot{\theta}\mathbf{e}_r$:
\begin{align*}
\mathbf{a} &= \ddot{r}\mathbf{e}_r + \dot{r}\dot{\theta}\mathbf{e}_\theta + \dot{r}\dot{\theta}\mathbf{e}_\theta + r\ddot{\theta}\mathbf{e}_\theta + r\dot{\theta}(-\dot{\theta}\mathbf{e}_r) \\
&= (\ddot{r} - r\dot{\theta}^2)\mathbf{e}_r + (r\ddot{\theta} + 2\dot{r}\dot{\theta})\mathbf{e}_\theta
\end{align*} \qed

\section{Applications to Planetary Motion}

\begin{theorem}{Kepler's Second Law}
A particle moving under a central force (force always directed toward a fixed point) sweeps out equal areas in equal times.
\end{theorem}

\textit{Proof.} For a central force, the torque is zero, so angular momentum $\mathbf{L} = \mathbf{r} \times m\mathbf{v}$ is constant. In polar coordinates, the area swept out in time $dt$ is $dA = \frac{1}{2}r^2 d\theta$, so:
\[
\frac{dA}{dt} = \frac{1}{2}r^2\dot{\theta} = \frac{\|\mathbf{L}\|}{2m} = \text{constant}
\] \qed

\begin{theorem}{Kepler's First Law}
Planets move in elliptical orbits with the sun at one focus.
\end{theorem}

\textit{Proof.} Using Newton's law of gravitation $\mathbf{F} = -\frac{GMm}{r^2}\mathbf{e}_r$ and the acceleration formula in polar coordinates, one can derive the orbital equation:
\[
r = \frac{L^2/(GMm^2)}{1 + e\cos\theta}
\]
which is the polar equation of a conic with eccentricity $e$. For bound orbits, $e < 1$ giving an ellipse. \qed



% ====================
% Part VI: Linear Algebra
% ====================
\part{Linear Algebra}

\chapter{Linear Spaces}

\section{Definition of a Linear Space}

\begin{definition}{Linear Space (Vector Space)}
A linear space (or vector space) is a nonempty set $V$ of objects, called elements or vectors, together with two operations called addition and multiplication by scalars (real numbers), subject to the following ten axioms for all vectors $x, y, z \in V$ and all scalars $a, b \in \mathbb{R}$:
\begin{enumerate}
    \item[(1)] Closure under addition: $x + y \in V$
    \item[(2)] Commutative law: $x + y = y + x$
    \item[(3)] Associative law: $(x + y) + z = x + (y + z)$
    \item[(4)] Existence of zero element: There exists $O \in V$ such that $x + O = x$ for all $x$
    \item[(5)] Existence of negatives: For each $x \in V$, there exists $-x \in V$ such that $x + (-x) = O$
    \item[(6)] Closure under scalar multiplication: $ax \in V$
    \item[(7)] Associative law: $a(bx) = (ab)x$
    \item[(8)] Distributive law for addition in $V$: $a(x + y) = ax + ay$
    \item[(9)] Distributive law for addition of scalars: $(a + b)x = ax + bx$
    \item[(10)] Identity: $1x = x$
\end{enumerate}
\end{definition}

\section{Examples of Linear Spaces}

\begin{theorem}{Examples}
The following are examples of linear spaces:
\begin{enumerate}
    \item[(a)] The real vector space $V_n$ of $n$-tuples of real numbers
    \item[(b)] The set of all real-valued functions defined on a set $S$
    \item[(c)] The set of all continuous functions on an interval $[a, b]$, denoted $C[a, b]$
    \item[(d)] The set of all polynomials of degree $\leq n$, denoted $\mathcal{P}_n$
    \item[(e)] The set of all solutions to a homogeneous linear differential equation
\end{enumerate}
\end{theorem}

\section{Subspaces}

\begin{definition}{Subspace}
A nonempty subset $S$ of a linear space $V$ is called a subspace if $S$ is itself a linear space under the same operations as $V$. Equivalently, $S$ is a subspace if:
\begin{enumerate}
    \item[(a)] $O \in S$
    \item[(b)] $S$ is closed under addition: $x, y \in S \implies x + y \in S$
    \item[(c)] $S$ is closed under scalar multiplication: $x \in S, a \in \mathbb{R} \implies ax \in S$
\end{enumerate}
\end{definition}

\section{Linear Dependence and Independence}

\begin{definition}{Linear Dependence and Independence}
A finite set $S = \{x_1, x_2, \ldots, x_k\}$ of elements in a linear space $V$ is said to be linearly dependent if there exist scalars $c_1, c_2, \ldots, c_k$, not all zero, such that:
\[
\sum_{i=1}^{k} c_ix_i = c_1x_1 + c_2x_2 + \cdots + c_kx_k = O
\]
If no such relation exists, the set $S$ is called linearly independent.
\end{definition}

\begin{theorem}{}
A set $S = \{x_1, x_2, \ldots, x_k\}$ is linearly dependent if and only if at least one element of $S$ can be expressed as a linear combination of the others.
\end{theorem}

\textit{Proof.} If $\sum c_ix_i = O$ with some $c_j \neq 0$, then $x_j = -\sum_{i \neq j} (c_i/c_j)x_i$. Conversely, if $x_j = \sum_{i \neq j} d_ix_i$, then $\sum_{i \neq j} d_ix_i - x_j = O$ gives a nontrivial relation. \qed

\section{Bases and Dimension}

\begin{definition}{Span}
The span of a set $S$ in a linear space $V$ is the set of all finite linear combinations of elements of $S$:
\[
\text{span}(S) = \left\{\sum_{i=1}^{k} c_ix_i : x_i \in S, c_i \in \mathbb{R}, k \in \mathbb{N}\right\}
\]
\end{definition}

\begin{definition}{Basis}
A set $B$ is a basis for a linear space $V$ if:
\begin{enumerate}
    \item[(a)] $B$ is linearly independent
    \item[(b)] $\text{span}(B) = V$
\end{enumerate}
\end{definition}

\begin{theorem}{Existence of Basis}
Every finite-dimensional linear space has a basis.
\end{theorem}

\begin{definition}{Dimension}
The dimension of a finite-dimensional linear space $V$, denoted $\dim(V)$, is the number of elements in any basis for $V$.
\end{definition}

\begin{theorem}{Uniqueness of Representation}
If $\{e_1, e_2, \ldots, e_n\}$ is a basis for $V$, then every element $x \in V$ can be written uniquely as:
\[
x = \sum_{i=1}^{n} c_ie_i
\]
\end{theorem}

\textit{Proof.} Since the basis spans $V$, such a representation exists. If $x = \sum c_ie_i = \sum d_ie_i$, then $\sum (c_i - d_i)e_i = O$. By linear independence, $c_i = d_i$ for all $i$. \qed

\section{Inner Products}

\begin{definition}{Inner Product}
An inner product on a real linear space $V$ is a function $(\cdot, \cdot): V \times V \to \mathbb{R}$ satisfying:
\begin{enumerate}
    \item[(a)] Symmetry: $(x, y) = (y, x)$ for all $x, y \in V$
    \item[(b)] Linearity in first argument: $(ax + by, z) = a(x, z) + b(y, z)$ for all $x, y, z \in V$ and $a, b \in \mathbb{R}$
    \item[(c)] Positive definiteness: $(x, x) \geq 0$ for all $x \in V$, and $(x, x) = 0$ if and only if $x = O$
\end{enumerate}
A linear space with an inner product is called an inner product space.
\end{definition}

\begin{definition}{Norm from Inner Product}
In an inner product space, the norm (length) of a vector $x$ is defined as:
\[
\|x\| = \sqrt{(x, x)}
\]
\end{definition}

\begin{theorem}{Cauchy-Schwarz Inequality}
In any inner product space:
\[
|(x, y)| \leq \|x\| \|y\|
\]
for all $x, y \in V$.
\end{theorem}

\textit{Proof.} For any real $t$, $(x + ty, x + ty) \geq 0$. Expanding:
\[
(x, x) + 2t(x, y) + t^2(y, y) \geq 0
\]
This quadratic in $t$ is nonnegative for all $t$, so its discriminant is non-positive:
\[
4(x, y)^2 - 4(x, x)(y, y) \leq 0
\]
Thus $(x, y)^2 \leq \|x\|^2 \|y\|^2$. \qed

\begin{theorem}{Triangle Inequality}
In any inner product space:
\[
\|x + y\| \leq \|x\| + \|y\|
\]
\end{theorem}

\textit{Proof.} $\|x + y\|^2 = (x + y, x + y) = \|x\|^2 + 2(x, y) + \|y\|^2 \leq \|x\|^2 + 2\|x\|\|y\| + \|y\|^2 = (\|x\| + \|y\|)^2$. \qed

\section{Orthogonality}

\begin{definition}{Orthogonality}
Two elements $x$ and $y$ in an inner product space are orthogonal if $(x, y) = 0$.
\end{definition}

\begin{theorem}{Pythagorean Theorem}
If $x$ and $y$ are orthogonal, then:
\[
\|x + y\|^2 = \|x\|^2 + \|y\|^2
\]
\end{theorem}

\textit{Proof.} $\|x + y\|^2 = (x + y, x + y) = \|x\|^2 + 2(x, y) + \|y\|^2 = \|x\|^2 + \|y\|^2$. \qed

\begin{theorem}{Best Approximation}
Let $S$ be a finite-dimensional subspace of an inner product space $V$, and let $\{e_1, \ldots, e_n\}$ be an orthonormal basis for $S$. For any $x \in V$, the element $s \in S$ that minimizes $\|x - s\|$ is:
\[
s = \sum_{i=1}^{n} (x, e_i)e_i
\]
This is called the orthogonal projection of $x$ onto $S$.
\end{theorem}

\textit{Proof.} For any $s \in S$, write $s = \sum c_ie_i$. Then:
\[
\|x - s\|^2 = (x - \sum c_ie_i, x - \sum c_ie_i) = \|x\|^2 - 2\sum c_i(x, e_i) + \sum c_i^2
\]
Completing the square in $c_i$:
\[
\|x - s\|^2 = \|x\|^2 - \sum (x, e_i)^2 + \sum (c_i - (x, e_i))^2
\]
This is minimized when $c_i = (x, e_i)$. \qed



\chapter{Linear Transformations and Matrices}

\section{Linear Transformations}

\begin{definition}{Linear Transformation}
Let $V$ and $W$ be linear spaces. A function $T: V \to W$ is called a linear transformation if:
\begin{enumerate}
    \item[(a)] $T(x + y) = T(x) + T(y)$ for all $x, y \in V$ (additivity)
    \item[(b)] $T(cx) = cT(x)$ for all $x \in V$ and all scalars $c$ (homogeneity)
\end{enumerate}
Equivalently, $T$ is linear if $T(ax + by) = aT(x) + bT(y)$ for all $x, y \in V$ and all scalars $a, b$.
\end{definition}

\begin{theorem}{Properties of Linear Transformations}
If $T: V \to W$ is a linear transformation, then:
\begin{enumerate}
    \item[(a)] $T(O_V) = O_W$ (maps zero to zero)
    \item[(b)] $T(-x) = -T(x)$ for all $x \in V$
    \item[(c)] $T(\sum_{i=1}^{n} c_ix_i) = \sum_{i=1}^{n} c_iT(x_i)$ for any linear combination
\end{enumerate}
\end{theorem}

\textit{Proof.} (a) $T(O_V) = T(0 \cdot O_V) = 0 \cdot T(O_V) = O_W$.

(b) $T(-x) = T((-1)x) = (-1)T(x) = -T(x)$. \qed

\section{Null Space and Range}

\begin{definition}{Null Space and Range}
Let $T: V \to W$ be a linear transformation.
\begin{enumerate}
    \item[(a)] The null space (or kernel) of $T$ is $N(T) = \{x \in V : T(x) = O_W\}$
    \item[(b)] The range (or image) of $T$ is $R(T) = \{T(x) : x \in V\}$
\end{enumerate}
\end{definition}

\begin{theorem}{Null Space and Range are Subspaces}
Let $T: V \to W$ be a linear transformation. Then:
\begin{enumerate}
    \item[(a)] $N(T)$ is a subspace of $V$
    \item[(b)] $R(T)$ is a subspace of $W$
\end{enumerate}
\end{theorem}

\textit{Proof.} (a) If $x, y \in N(T)$, then $T(x + y) = T(x) + T(y) = O + O = O$, so $x + y \in N(T)$. If $x \in N(T)$ and $c$ is a scalar, $T(cx) = cT(x) = cO = O$, so $cx \in N(T)$.

(b) Similar proof for $R(T)$. \qed

\begin{definition}{Nullity and Rank}
For a linear transformation $T: V \to W$:
\begin{enumerate}
    \item[(a)] The nullity of $T$ is $\dim N(T)$
    \item[(b)] The rank of $T$ is $\dim R(T)$
\end{enumerate}
\end{definition}

\begin{theorem}{Rank-Nullity Theorem}
Let $T: V \to W$ be a linear transformation with $V$ finite-dimensional. Then:
\[
\dim V = \text{nullity}(T) + \text{rank}(T)
\]
\end{theorem}

\textit{Proof.} Let $\{e_1, \ldots, e_k\}$ be a basis for $N(T)$. Extend to a basis $\{e_1, \ldots, e_k, e_{k+1}, \ldots, e_n\}$ for $V$. We claim $\{T(e_{k+1}), \ldots, T(e_n)\}$ is a basis for $R(T)$.

Spanning: Any $y \in R(T)$ has $y = T(x)$ for some $x = \sum c_ie_i$. Then $y = \sum c_iT(e_i) = \sum_{i=k+1}^{n} c_iT(e_i)$ since $T(e_i) = O$ for $i \leq k$.

Independence: If $\sum_{i=k+1}^{n} c_iT(e_i) = O$, then $T(\sum_{i=k+1}^{n} c_ie_i) = O$, so $\sum_{i=k+1}^{n} c_ie_i \in N(T)$. Thus $\sum_{i=k+1}^{n} c_ie_i = \sum_{i=1}^{k} d_ie_i$, giving $\sum_{i=k+1}^{n} c_ie_i - \sum_{i=1}^{k} d_ie_i = O$. By independence of $\{e_i\}$, all $c_i = 0$.

Thus $\dim R(T) = n - k = \dim V - \dim N(T)$. \qed

\section{Matrix Representation}

\begin{definition}{Matrix of a Linear Transformation}
Let $T: V \to W$ be a linear transformation with $\dim V = n$ and $\dim W = m$. Let $\{e_1, \ldots, e_n\}$ be a basis for $V$ and $\{f_1, \ldots, f_m\}$ be a basis for $W$. The matrix of $T$ relative to these bases is the $m \times n$ matrix $A = (a_{ij})$ where:
\[
T(e_j) = \sum_{i=1}^{m} a_{ij}f_i \quad \text{for } j = 1, \ldots, n
\]
\end{definition}

\begin{theorem}{Matrix Representation}
Let $T: V \to W$ have matrix $A$ relative to bases $\{e_j\}$ of $V$ and $\{f_i\}$ of $W$. If $x = \sum_{j=1}^{n} x_je_j \in V$ and $y = T(x) = \sum_{i=1}^{m} y_if_i \in W$, then:
\[
y_i = \sum_{j=1}^{n} a_{ij}x_j \quad \text{for } i = 1, \ldots, m
\]
In matrix form: $y = Ax$.
\end{theorem}

\textit{Proof.} $T(x) = T(\sum_j x_je_j) = \sum_j x_jT(e_j) = \sum_j x_j\sum_i a_{ij}f_i = \sum_i(\sum_j a_{ij}x_j)f_i$. Comparing with $y = \sum_i y_if_i$ gives the result. \qed

\section{Composition and Matrix Multiplication}

\begin{definition}{Composition of Linear Transformations}
If $T: U \to V$ and $S: V \to W$ are linear transformations, the composition $S \circ T: U \to W$ is defined by:
\[
(S \circ T)(x) = S(T(x)) \quad \text{for all } x \in U
\]
\end{definition}

\begin{theorem}{Matrix of Composition}
Let $T: U \to V$ have matrix $B$ and $S: V \to W$ have matrix $A$ (relative to appropriate bases). Then the composition $S \circ T: U \to W$ has matrix $AB$.
\end{theorem}

\textit{Proof.} Let $\dim U = p$, $\dim V = n$, $\dim W = m$. For $x \in U$:
\[
(S \circ T)(x) = S(T(x)) = S(Bx) = A(Bx) = (AB)x
\]
Thus $S \circ T$ has matrix $AB$. \qed

\section{Invertibility}

\begin{definition}{Invertible Linear Transformation}
A linear transformation $T: V \to W$ is invertible if there exists a linear transformation $T^{-1}: W \to V$ such that:
\[
T^{-1} \circ T = I_V \quad \text{and} \quad T \circ T^{-1} = I_W
\]
where $I_V$ and $I_W$ are the identity transformations on $V$ and $W$.
\end{definition}

\begin{theorem}{Invertibility Criteria}
Let $T: V \to W$ be a linear transformation with $V$ and $W$ finite-dimensional and $\dim V = \dim W = n$. The following are equivalent:
\begin{enumerate}
    \item[(a)] $T$ is invertible
    \item[(b)] $T$ is injective (one-to-one)
    \item[(c)] $T$ is surjective (onto)
    \item[(d)] $\text{rank}(T) = n$
    \item[(e)] The matrix of $T$ (relative to any bases) is invertible
\end{enumerate}
\end{theorem}

\textit{Proof.} By the rank-nullity theorem: $\dim V = \text{nullity}(T) + \text{rank}(T) = n$.

(a) $\Leftrightarrow$ (b): Standard result for any function.

(b) $\Leftrightarrow$ (d): $T$ is injective iff $N(T) = \{O\}$ iff $\text{nullity}(T) = 0$ iff $\text{rank}(T) = n$.

(c) $\Leftrightarrow$ (d): $T$ is surjective iff $R(T) = W$ iff $\text{rank}(T) = \dim W = n$.

(d) $\Leftrightarrow$ (e): The matrix $A$ is invertible iff $Ax = 0$ has only the trivial solution, which is equivalent to $N(T) = \{O\}$. \qed

\section{Eigenvalues and Eigenvectors}

\begin{definition}{Eigenvalue and Eigenvector}
Let $T: V \to V$ be a linear transformation. A scalar $\lambda$ is called an eigenvalue of $T$ if there exists a nonzero vector $x \in V$ such that:
\[
T(x) = \lambda x
\]
Such a vector $x$ is called an eigenvector of $T$ corresponding to $\lambda$.
\end{definition}

\begin{definition}{Characteristic Polynomial}
Let $A$ be the matrix of $T: V \to V$ relative to some basis. The characteristic polynomial of $T$ (and of $A$) is:
\[
p(\lambda) = \det(A - \lambda I)
\]
\end{definition}

\begin{theorem}{Eigenvalues and Characteristic Polynomial}
The eigenvalues of a linear transformation $T$ are the roots of its characteristic polynomial.
\end{theorem}

\textit{Proof.} $\lambda$ is an eigenvalue iff $T(x) = \lambda x$ for some $x \neq O$, iff $(T - \lambda I)(x) = O$ for some $x \neq O$, iff $T - \lambda I$ is not invertible, iff $\det(A - \lambda I) = 0$. \qed

\section{Systems of Linear Equations}

\begin{definition}{System of Linear Equations}
A system of $m$ linear equations in $n$ unknowns can be written in matrix form as:
\[
Ax = b
\]
where $A$ is an $m \times n$ matrix, $x$ is an $n \times 1$ column vector of unknowns, and $b$ is an $m \times 1$ column vector.
\end{definition}

\begin{theorem}{Solution Set Structure}
Let $A$ be an $m \times n$ matrix. The solution set of the homogeneous system $Ax = O$ is the null space of the linear transformation $T(x) = Ax$, hence is a subspace of $\mathbb{R}^n$.

If $x_p$ is a particular solution of $Ax = b$, then every solution has the form $x = x_p + x_h$ where $x_h$ is a solution of $Ax = O$.
\end{theorem}

\textit{Proof.} If $Ax_p = b$ and $Ax = b$, then $A(x - x_p) = O$, so $x - x_p = x_h$ for some $x_h$ with $Ax_h = O$. Thus $x = x_p + x_h$. \qed

\section{Gaussian Elimination}

\begin{theorem}{Elementary Row Operations}
The following elementary row operations on a matrix do not change the solution set of the corresponding linear system:
\begin{enumerate}
    \item[(a)] Interchanging two rows
    \item[(b)] Multiplying a row by a nonzero scalar
    \item[(c)] Adding a multiple of one row to another row
\end{enumerate}
\end{theorem}

\begin{theorem}{Row Echelon Form}
By elementary row operations, any matrix can be transformed to row echelon form, from which the solution set of the corresponding linear system can be easily determined.
\end{theorem}



\chapter{Determinants}


\section{Introduction}

\begin{definition}{Determinant Function}
Let $A = (a_{ij})$ be an $n \times n$ matrix. The determinant of $A$, denoted by $\det A$ or $|A|$, is defined inductively as follows:
\begin{enumerate}
    \item[(a)] If $n = 1$, then $\det A = a_{11}$
    \item[(b)] If $n > 1$, then $\det A = \sum_{j=1}^{n} a_{1j} \cdot \text{cof}(a_{1j})$, where $\text{cof}(a_{1j}) = (-1)^{1+j} \cdot \det A_{1j}$ is the cofactor of $a_{1j}$, and $A_{1j}$ is the $(n-1) \times (n-1)$ submatrix obtained by deleting the first row and $j$-th column of $A$.
\end{enumerate}
\end{definition}

\section{Fundamental Properties of Determinants}

\begin{theorem}{Effect of Row Operations on Determinants}
Let $A$ be an $n \times n$ matrix.
\begin{enumerate}
    \item[(a)] If $B$ is obtained from $A$ by interchanging two rows, then $\det B = -\det A$
    \item[(b)] If $B$ is obtained from $A$ by multiplying a row by a scalar $c$, then $\det B = c \cdot \det A$
    \item[(c)] If $B$ is obtained from $A$ by adding a multiple of one row to another row, then $\det B = \det A$
\end{enumerate}
\end{theorem}

\textit{Proof.} These follow from the definition by induction on $n$. \qed

\begin{theorem}{Zero Row or Column}
If a matrix $A$ has a row or column of zeros, then $\det A = 0$.
\end{theorem}

\textit{Proof.} Expand along the zero row or column. \qed

\begin{theorem}{Equal Rows or Columns}
If a matrix $A$ has two identical rows or two identical columns, then $\det A = 0$.
\end{theorem}

\textit{Proof.} If two rows are identical, interchanging them gives the same matrix, so $\det A = -\det A$, hence $\det A = 0$. \qed

\section{The Determinant of a Transpose}

\begin{theorem}{Transpose Property}
For any square matrix $A$:
\[
\det A = \det A^t
\]
where $A^t$ is the transpose of $A$.
\end{theorem}

\textit{Proof.} We prove this by induction on $n$. For $n = 1$, $A = A^t$ so the result is trivial. Assume the theorem holds for matrices of size $(n-1) \times (n-1)$.

Expanding $\det A$ by its first-column minors and $\det A^t$ by its first-row minors, we have:
\[
\det A = \sum_{i=1}^{n} (-1)^{i+1} a_{i1} \det A_{i1}, \quad \det A^t = \sum_{j=1}^{n} (-1)^{1+j} a_{1j}^t \det A_{1j}^t
\]

But from the definition of transpose we have $a_{i1} = a_{1i}^t$ and $A_{i1} = A_{1i}^t$. Since we are assuming the theorem for determinants of size $(n-1) \times (n-1)$, we have $\det A_{i1} = \det A_{1i}^t$. Hence the foregoing sums are equal term by term, so $\det A = \det A^t$. \qed

\section{The Cofactor Matrix}

\begin{definition}{Cofactor Matrix}
Let $A$ be an $n \times n$ matrix. The cofactor matrix of $A$, denoted $\text{cof } A$, is the matrix whose $(i,j)$ entry is $\text{cof}(a_{ij})$, the cofactor of $a_{ij}$.
\[
(\text{cof } A)_{ij} = (-1)^{i+j} \det A_{ij}
\]
where $A_{ij}$ is the submatrix obtained by deleting the $i$-th row and $j$-th column of $A$.
\end{definition}

\begin{theorem}{Cofactor Expansion}
For any $n \times n$ matrix $A$ with $n \geq 2$:
\[
A(\text{cof } A)^t = (\det A)I
\]
In particular, if $\det A \neq 0$, the inverse of $A$ exists and is given by:
\[
A^{-1} = \frac{1}{\det A} (\text{cof } A)^t
\]
\end{theorem}

\textit{Proof.} The $(i,k)$ entry of $A(\text{cof } A)^t$ is $\sum_{j=1}^{n} a_{ij} \cdot \text{cof}(a_{kj})$. When $i = k$, this sum equals $\det A$ by the cofactor expansion. When $i \neq k$, this sum equals the determinant of a matrix with two identical rows (the $i$-th and $k$-th rows both equal to the $i$-th row of $A$), which is 0. Thus $A(\text{cof } A)^t = (\det A)I$. \qed

\section{Cramer's Rule}

\begin{theorem}{Cramer's Rule}
If a system of $n$ linear equations in $n$ unknowns
\[
\sum_{j=1}^{n} a_{ij}x_j = b_i \quad (i = 1, 2, \ldots, n)
\]
has a nonsingular coefficient matrix $A = (a_{ij})$, then there is a unique solution for the system given by the formulas:
\[
x_j = \frac{1}{\det A} \sum_{k=1}^{n} b_k \cdot \text{cof}(a_{kj}) \quad \text{for } j = 1, 2, \ldots, n
\]
Equivalently:
\[
x_j = \frac{\det C_j}{\det A}
\]
where $C_j$ is the matrix obtained from $A$ by replacing the $j$-th column of $A$ by the column matrix $B = (b_1, \ldots, b_n)^t$.
\end{theorem}

\textit{Proof.} The system can be written in matrix form:
\[
AX = B
\]
where $X$ and $B$ are column matrices. Since $A$ is nonsingular, there is a unique solution given by:
\[
X = A^{-1}B = \frac{1}{\det A}(\text{cof } A)^t B
\]
The formula follows by equating components.

To derive the equivalent form, note that:
\[
\sum_{k=1}^{n} b_k \cdot \text{cof}(a_{kj}) = \det C_j
\]
by expanding $\det C_j$ along the $j$-th column, where $C_j$ is obtained from $A$ by replacing column $j$ with $B$. \qed



\chapter{Eigenvalues and Operators}

\section{Eigenvalues and Inner Products}

\begin{definition}{Hermitian and Skew-Hermitian Operators}
Let $V$ be a Euclidean space with inner product $(\cdot, \cdot)$. A linear operator $T: V \to V$ is called:
\begin{enumerate}
    \item[(a)] \textbf{Hermitian} (or self-adjoint) if $(T(x), y) = (x, T(y))$ for all $x, y \in V$
    \item[(b)] \textbf{Skew-Hermitian} if $(T(x), y) = -(x, T(y))$ for all $x, y \in V$
\end{enumerate}
\end{definition}

\begin{theorem}{Eigenvalues of Hermitian Operators}
Let $T: V \to V$ be a Hermitian operator on a finite-dimensional Euclidean space $V$. Then:
\begin{enumerate}
    \item[(a)] All eigenvalues of $T$ are real
    \item[(b)] Eigenvectors corresponding to distinct eigenvalues are orthogonal
\end{enumerate}
\end{theorem}

\textit{Proof.} (a) Let $T(x) = \lambda x$ with $x \neq O$. Then:
\[
\lambda(x, x) = (\lambda x, x) = (T(x), x) = (x, T(x)) = (x, \lambda x) = \bar{\lambda}(x, x)
\]
Since $(x, x) > 0$, we have $\lambda = \bar{\lambda}$, so $\lambda$ is real.

(b) Let $T(x) = \lambda x$ and $T(y) = \mu y$ with $\lambda \neq \mu$. Then:
\[
\lambda(x, y) = (T(x), y) = (x, T(y)) = (x, \mu y) = \mu(x, y)
\]
Since $\lambda \neq \mu$, we have $(x, y) = 0$. \qed

\section{Orthogonal and Unitary Transformations}

\begin{definition}{Orthogonal Transformation}
A linear transformation $T: V \to V$ on a real Euclidean space is called orthogonal if it preserves inner products:
\[
(T(x), T(y)) = (x, y) \quad \text{for all } x, y \in V
\]
Equivalently, $T$ is orthogonal if $\|T(x)\| = \|x\|$ for all $x \in V$.
\end{definition}

\begin{definition}{Unitary Transformation}
A linear transformation $T: V \to V$ on a complex Euclidean space (unitary space) is called unitary if:
\[
(T(x), T(y)) = (x, y) \quad \text{for all } x, y \in V
\]
\end{definition}

\begin{theorem}{Properties of Orthogonal Transformations}
Let $T: V \to V$ be an orthogonal transformation on a finite-dimensional real Euclidean space. Then:
\begin{enumerate}
    \item[(a)] $T$ is invertible and $T^{-1}$ is orthogonal
    \item[(b)] The eigenvalues of $T$ have absolute value $1$
    \item[(c)] If $\lambda$ is an eigenvalue of $T$, then $\lambda = \pm 1$
    \item[(d)] $\det T = \pm 1$
\end{enumerate}
\end{theorem}

\begin{theorem}{Orthogonal and Unitary Matrices}
\begin{enumerate}
    \item[(a)] A real matrix $A$ is orthogonal if and only if $A^TA = AA^T = I$
    \item[(b)] A complex matrix $A$ is unitary if and only if $A^*A = AA^* = I$, where $A^*$ is the conjugate transpose
\end{enumerate}
\end{theorem}

\section{Spectral Theorem}

\begin{theorem}{Spectral Theorem for Hermitian Operators}
Let $T: V \to V$ be a Hermitian operator on a finite-dimensional complex Euclidean space. Then there exists an orthonormal basis for $V$ consisting of eigenvectors of $T$. Equivalently, there exists a unitary matrix $U$ such that:
\[
U^*TU = D
\]
where $D$ is a real diagonal matrix.
\end{theorem}

\textit{Proof.} By induction on $\dim V$. For $\dim V = 1$, the result is trivial. Assume true for spaces of dimension less than $n$. For $\dim V = n$, let $\lambda_1$ be an eigenvalue of $T$ with unit eigenvector $u_1$. Let $W = \text{span}\{u_1\}$ and $W^\perp$ be the orthogonal complement. Then $T$ maps $W^\perp$ to $W^\perp$, so by induction, $W^\perp$ has an orthonormal basis of eigenvectors. Adding $u_1$ gives an orthonormal basis for $V$. \qed

\begin{theorem}{Spectral Theorem for Real Symmetric Matrices}
Let $A$ be a real symmetric $n \times n$ matrix. Then there exists an orthogonal matrix $Q$ such that:
\[
Q^TAQ = D
\]
where $D$ is a real diagonal matrix containing the eigenvalues of $A$.
\end{theorem}

\textit{Proof.} A real symmetric matrix represents a Hermitian operator on $\mathbb{R}^n$. By the spectral theorem, there exists an orthonormal basis of eigenvectors. The matrix $Q$ with these eigenvectors as columns is orthogonal, and $Q^TAQ = D$. \qed

\section{Normal Operators}

\begin{definition}{Normal Operator}
A linear operator $T: V \to V$ on a Euclidean space is called normal if it commutes with its adjoint:
\[
TT^* = T^*T
\]
where $T^*$ is the adjoint of $T$ defined by $(T(x), y) = (x, T^*(y))$.
\end{definition}

\begin{theorem}{Properties of Normal Operators}
Let $T: V \to V$ be a normal operator on a finite-dimensional complex Euclidean space. Then:
\begin{enumerate}
    \item[(a)] $\|T(x)\| = \|T^*(x)\|$ for all $x \in V$
    \item[(b)] $T$ and $T^*$ have the same eigenvectors; if $T(x) = \lambda x$, then $T^*(x) = \bar{\lambda} x$
    \item[(c)] Eigenvectors corresponding to distinct eigenvalues are orthogonal
\end{enumerate}
\end{theorem}

\begin{theorem}{Spectral Theorem for Normal Operators}
A linear operator $T$ on a finite-dimensional complex Euclidean space is normal if and only if there exists an orthonormal basis for $V$ consisting of eigenvectors of $T$.
\end{theorem}




% ====================
% Additional Chapter: Linear Differential Equations (Advanced)
% ====================
\chapter{Linear Differential Equations}

\section{Linear Differential Operators}

\begin{definition}{Linear Differential Operator}
Let $C^{(n)}(I)$ denote the linear space of all real-valued functions with $n$ continuous derivatives on an interval $I$. A linear differential operator $L$ of order $n$ is a linear transformation $L: C^{(n)}(I) \to C(I)$ of the form:
\[
L = D^n + a_1D^{n-1} + \cdots + a_{n-1}D + a_n
\]
where $D = \frac{d}{dx}$ is the differentiation operator and $a_1, \ldots, a_n$ are functions in $C(I)$.
\end{definition}

\begin{theorem}{Linearity of $L$}
The operator $L$ is a linear transformation.
\end{theorem}

\textit{Proof.} Direct verification using the linearity of differentiation. \qed

\section{Existence and Uniqueness for Linear Equations}

\begin{theorem}{Existence-Uniqueness Theorem}
Let $a_1, \ldots, a_n$ and $R$ be continuous on an open interval $I$, and let $x_0 \in I$. Then for any initial values $b_0, b_1, \ldots, b_{n-1}$, there exists exactly one solution $y = f(x)$ of:
\[
L(y) = R(x), \quad f(x_0) = b_0, \quad f'(x_0) = b_1, \quad \ldots, \quad f^{(n-1)}(x_0) = b_{n-1}
\]
\end{theorem}

\section{The Dimension of the Solution Space}

\begin{definition}{Solution Space}
The set of all solutions of $L(y) = 0$ is called the solution space of $L$, denoted by $N(L)$.
\end{definition}

\begin{theorem}{Dimension of Solution Space}
The solution space $N(L)$ of an $n$-th order linear homogeneous equation has dimension $n$.
\end{theorem}

\begin{definition}{Wronskian}
Let $y_1, \ldots, y_n$ be functions with $n-1$ derivatives on $I$. The Wronskian is:
\[
W(y_1, \ldots, y_n) = \det\begin{pmatrix}
y_1 & y_2 & \cdots & y_n \\
y_1' & y_2' & \cdots & y_n' \\
\vdots & \vdots & \ddots & \vdots \\
y_1^{(n-1)} & y_2^{(n-1)} & \cdots & y_n^{(n-1)}
\end{pmatrix}
\]
\end{definition}

\begin{theorem}{Wronskian Criterion}
Solutions $y_1, \ldots, y_n$ of $L(y) = 0$ are linearly independent iff $W \neq 0$.
\end{theorem}

\section{Constant-Coefficient Operators}

\begin{theorem}{Factorization}
If $p(r) = (r - r_1)\cdots(r - r_n)$ is the characteristic polynomial, then:
\[
L = (D - r_1I)\cdots(D - r_nI)
\]
where the factors commute.
\end{theorem}

\begin{theorem}{Basis Solutions}
Let $p(r)$ have roots $r_1, \ldots, r_n$ (with multiplicities). Then:
\begin{enumerate}
    \item[(a)] If $r$ is real with multiplicity $m$: $e^{rx}, xe^{rx}, \ldots, x^{m-1}e^{rx}$ are solutions.
    \item[(b)] If $\alpha \pm i\beta$ have multiplicity $m$: $e^{\alpha x}\cos\beta x, e^{\alpha x}\sin\beta x, \ldots, x^{m-1}e^{\alpha x}\sin\beta x$ are solutions.
\end{enumerate}
\end{theorem}

\section{The Nonhomogeneous Equation}

\begin{theorem}{General Solution}
Let $y_p$ be a particular solution of $L(y) = R(x)$. Then every solution has the form $y = y_p + y_h$ where $L(y_h) = 0$.
\end{theorem}

\section{The Method of Undetermined Coefficients}

\begin{theorem}{Annihilator Method}
If $R(x)$ satisfies $M(R) = 0$ for some operator $M$, then every solution of $L(y) = R$ satisfies $(M \circ L)(y) = 0$.
\end{theorem}



\chapter{Systems of Differential Equations}

\section{Introduction to Systems}

\begin{definition}{System of First-Order Equations}
A system of $n$ first-order differential equations in $n$ unknown functions $y_1, \ldots, y_n$ has the form:
\[
y_i' = f_i(x, y_1, \ldots, y_n), \quad i = 1, \ldots, n
\]
In vector notation: $\mathbf{y}' = \mathbf{f}(x, \mathbf{y})$ where $\mathbf{y} = (y_1, \ldots, y_n)$.
\end{definition}

\begin{definition}{Linear System}
A system is called linear if it has the form:
\[
\mathbf{y}' = A(x)\mathbf{y} + \mathbf{Q}(x)
\]
where $A(x)$ is an $n \times n$ matrix and $\mathbf{Q}(x)$ is an $n$-dimensional vector. The system is homogeneous if $\mathbf{Q}(x) = \mathbf{0}$.
\end{definition}

\begin{theorem}{Existence and Uniqueness for Systems}
If $A(x)$ and $\mathbf{Q}(x)$ are continuous on an interval $I$, then for any $x_0 \in I$ and any initial vector $\mathbf{c}$, there exists a unique solution $\mathbf{y}$ of:
\[
\mathbf{y}' = A(x)\mathbf{y} + \mathbf{Q}(x), \quad \mathbf{y}(x_0) = \mathbf{c}
\]
on the entire interval $I$.
\end{theorem}

\section{Linear Systems with Constant Coefficients}

\begin{definition}{Homogeneous Constant-Coefficient System}
A system of the form $\mathbf{y}' = A\mathbf{y}$ where $A$ is a constant $n \times n$ matrix.
\end{definition}

\begin{theorem}{Solution via Eigenvalues}
If $\mathbf{v}$ is an eigenvector of $A$ with eigenvalue $\lambda$, then $\mathbf{y}(x) = e^{\lambda x}\mathbf{v}$ is a solution of $\mathbf{y}' = A\mathbf{y}$.
\end{theorem}

\textit{Proof.} $\mathbf{y}' = \lambda e^{\lambda x}\mathbf{v} = e^{\lambda x}A\mathbf{v} = A(e^{\lambda x}\mathbf{v}) = A\mathbf{y}$. \qed

\begin{theorem}{General Solution of Diagonalizable System}
If $A$ has $n$ linearly independent eigenvectors $\mathbf{v}_1, \ldots, \mathbf{v}_n$ with eigenvalues $\lambda_1, \ldots, \lambda_n$, then the general solution of $\mathbf{y}' = A\mathbf{y}$ is:
\[
\mathbf{y}(x) = c_1e^{\lambda_1 x}\mathbf{v}_1 + \cdots + c_ne^{\lambda_n x}\mathbf{v}_n
\]
\end{theorem}

\section{The Matrix Exponential}

\begin{definition}{Matrix Exponential}
For any $n \times n$ matrix $A$, the matrix exponential $e^A$ is defined by the series:
\[
e^A = \sum_{k=0}^{\infty} \frac{A^k}{k!} = I + A + \frac{A^2}{2!} + \frac{A^3}{3!} + \cdots
\]
\end{definition}

\begin{theorem}{Convergence of Matrix Exponential}
The series for $e^A$ converges for all matrices $A$.
\end{theorem}

\begin{theorem}{Properties of Matrix Exponential}
\begin{enumerate}
    \item[(a)] $e^{0} = I$
    \item[(b)] $e^{A}e^{B} = e^{A+B}$ if $AB = BA$
    \item[(c)] $(e^A)^{-1} = e^{-A}$
    \item[(d)] $\frac{d}{dx}e^{xA} = Ae^{xA} = e^{xA}A$
\end{enumerate}
\end{theorem}

\begin{theorem}{Solution via Matrix Exponential}
The unique solution of $\mathbf{y}' = A\mathbf{y}$ with $\mathbf{y}(0) = \mathbf{c}$ is:
\[
\mathbf{y}(x) = e^{xA}\mathbf{c}
\]
\end{theorem}

\textit{Proof.} Differentiating: $\frac{d}{dx}e^{xA}\mathbf{c} = Ae^{xA}\mathbf{c} = A\mathbf{y}$. At $x=0$: $\mathbf{y}(0) = e^{0}\mathbf{c} = \mathbf{c}$. \qed

\section{Nonhomogeneous Linear Systems}

\begin{theorem}{Variation of Parameters for Systems}
The general solution of $\mathbf{y}' = A(x)\mathbf{y} + \mathbf{Q}(x)$ is:
\[
\mathbf{y}(x) = \Phi(x)\mathbf{c} + \Phi(x)\int \Phi(x)^{-1}\mathbf{Q}(x)\,dx
\]
where $\Phi(x)$ is a fundamental matrix (matrix whose columns are linearly independent solutions of the homogeneous system).
\end{theorem}

\textit{Proof.} Using the variation of parameters ansatz $\mathbf{y} = \Phi\mathbf{v}$ and substituting yields $\mathbf{v}' = \Phi^{-1}\mathbf{Q}$. \qed

\section{Phase Plane Analysis}

\begin{definition}{Phase Plane}
For a 2D system $\mathbf{y}' = A\mathbf{y}$, the phase plane is the $(y_1, y_2)$-plane where solution curves (trajectories) are plotted.
\end{definition}

\begin{definition}{Critical Point}
A point $\mathbf{y}_0$ where $A\mathbf{y}_0 = \mathbf{0}$ is called a critical point (or equilibrium point).
\end{definition}

\begin{theorem}{Classification of Critical Points}
For $\mathbf{y}' = A\mathbf{y}$ with $\det A \neq 0$, the origin is a critical point classified by eigenvalues:
\begin{enumerate}
    \item[(a)] Real, distinct, same sign: Node (attracting if negative, repelling if positive)
    \item[(b)] Real, opposite signs: Saddle point
    \item[(c)] Complex with nonzero real part: Spiral point
    \item[(d)] Purely imaginary: Center
    \item[(e)] Repeated real: Improper node
\end{enumerate}
\end{theorem}


\end{document}

\chapter{Differential Calculus of Scalar and Vector Fields}

\section{Functions from $\mathbb{R}^n$ to $\mathbb{R}^m$}

\begin{definition}{Scalar and Vector Fields}
Let $S$ be a subset of $\mathbb{R}^n$. A function $f: S \to \mathbb{R}$ is called a scalar field. A function $\mathbf{f}: S \to \mathbb{R}^m$ with $m > 1$ is called a vector field.
\end{definition}

\begin{definition}{Limit of a Function}
Let $\mathbf{f}: S \to \mathbb{R}^m$ where $S \subseteq \mathbb{R}^n$. We say $\mathbf{f}(\mathbf{x}) \to \mathbf{b}$ as $\mathbf{x} \to \mathbf{a}$, written:
\[
\lim_{\mathbf{x} \to \mathbf{a}} \mathbf{f}(\mathbf{x}) = \mathbf{b}
\]
if for every $\epsilon > 0$ there exists $\delta > 0$ such that $\|\mathbf{f}(\mathbf{x}) - \mathbf{b}\| < \epsilon$ whenever $0 < \|\mathbf{x} - \mathbf{a}\| < \delta$ and $\mathbf{x} \in S$.
\end{definition}

\begin{definition}{Continuity}
A function $\mathbf{f}: S \to \mathbb{R}^m$ is continuous at $\mathbf{a} \in S$ if $\lim_{\mathbf{x} \to \mathbf{a}} \mathbf{f}(\mathbf{x}) = \mathbf{f}(\mathbf{a})$.
\end{definition}

\section{Directional Derivatives and Partial Derivatives}

\begin{definition}{Directional Derivative}
Let $f: S \to \mathbb{R}$ be a scalar field and $\mathbf{u}$ a unit vector in $\mathbb{R}^n$. The directional derivative of $f$ at $\mathbf{a}$ in the direction $\mathbf{u}$ is:
\[
f'(\mathbf{a}; \mathbf{u}) = \lim_{h \to 0} \frac{f(\mathbf{a} + h\mathbf{u}) - f(\mathbf{a})}{h}
\]
when this limit exists.
\end{definition}

\begin{definition}{Partial Derivative}
The partial derivative of $f$ with respect to the $j$-th coordinate at $\mathbf{a} = (a_1, \ldots, a_n)$ is:
\[
D_j f(\mathbf{a}) = \frac{\partial f}{\partial x_j}(\mathbf{a}) = \lim_{h \to 0} \frac{f(a_1, \ldots, a_j + h, \ldots, a_n) - f(\mathbf{a})}{h}
\]
\end{definition}

\begin{theorem}{Relation Between Directional and Partial Derivatives}
If all partial derivatives of $f$ exist and are continuous in a neighborhood of $\mathbf{a}$, then for any unit vector $\mathbf{u} = (u_1, \ldots, u_n)$:
\[
f'(\mathbf{a}; \mathbf{u}) = \sum_{j=1}^{n} D_j f(\mathbf{a}) \cdot u_j = \nabla f(\mathbf{a}) \cdot \mathbf{u}
\]
\end{theorem}

\section{Total Derivative}

\begin{definition}{Differentiability}
A function $\mathbf{f}: S \to \mathbb{R}^m$ is differentiable at an interior point $\mathbf{a}$ of $S$ if there exists a linear transformation $T: \mathbb{R}^n \to \mathbb{R}^m$ such that:
\[
\mathbf{f}(\mathbf{a} + \mathbf{v}) = \mathbf{f}(\mathbf{a}) + T(\mathbf{v}) + \|\mathbf{v}\| \mathbf{E}(\mathbf{a}, \mathbf{v})
\]
where $\mathbf{E}(\mathbf{a}, \mathbf{v}) \to \mathbf{0}$ as $\mathbf{v} \to \mathbf{0}$. The linear transformation $T$ is called the total derivative of $\mathbf{f}$ at $\mathbf{a}$, denoted $\mathbf{f}'(\mathbf{a})$.
\end{definition}

\begin{theorem}{Uniqueness of Total Derivative}
If $\mathbf{f}$ is differentiable at $\mathbf{a}$, the total derivative is unique.
\end{theorem}

\begin{theorem}{Differentiability Implies Continuity}
If $\mathbf{f}$ is differentiable at $\mathbf{a}$, then $\mathbf{f}$ is continuous at $\mathbf{a}$.
\end{theorem}

\textit{Proof.} As $\mathbf{v} \to \mathbf{0}$, we have $\mathbf{f}(\mathbf{a} + \mathbf{v}) - \mathbf{f}(\mathbf{a}) = T(\mathbf{v}) + \|\mathbf{v}\|\mathbf{E}(\mathbf{a}, \mathbf{v}) \to \mathbf{0}$. \qed

\begin{theorem}{Jacobian Matrix}
If $\mathbf{f} = (f_1, \ldots, f_m): S \to \mathbb{R}^m$ is differentiable at $\mathbf{a}$, then all partial derivatives exist and the matrix of $T = \mathbf{f}'(\mathbf{a})$ is the Jacobian matrix:
\[
D\mathbf{f}(\mathbf{a}) = \begin{pmatrix}
D_1 f_1(\mathbf{a}) & D_2 f_1(\mathbf{a}) & \cdots & D_n f_1(\mathbf{a}) \\
D_1 f_2(\mathbf{a}) & D_2 f_2(\mathbf{a}) & \cdots & D_n f_2(\mathbf{a}) \\
\vdots & \vdots & \ddots & \vdots \\
D_1 f_m(\mathbf{a}) & D_2 f_m(\mathbf{a}) & \cdots & D_n f_m(\mathbf{a})
\end{pmatrix}
\]
\end{theorem}

\section{The Gradient Vector}

\begin{definition}{Gradient}
For a scalar field $f: \mathbb{R}^n \to \mathbb{R}$ with continuous partial derivatives, the gradient is the vector field:
\[
\nabla f(\mathbf{x}) = (D_1 f(\mathbf{x}), D_2 f(\mathbf{x}), \ldots, D_n f(\mathbf{x})) = \sum_{j=1}^{n} D_j f(\mathbf{x}) \mathbf{e}_j
\]
\end{definition}

\begin{theorem}{Direction of Maximum Increase}
The gradient $\nabla f(\mathbf{a})$ points in the direction of maximum rate of increase of $f$ at $\mathbf{a}$, and its magnitude equals that maximum rate:
\[
\|\nabla f(\mathbf{a})\| = \max_{\|\mathbf{u}\|=1} f'(\mathbf{a}; \mathbf{u})
\]
\end{theorem}

\textit{Proof.} By Cauchy-Schwarz, $f'(\mathbf{a}; \mathbf{u}) = \nabla f(\mathbf{a}) \cdot \mathbf{u} \leq \|\nabla f(\mathbf{a})\| \|\mathbf{u}\| = \|\nabla f(\mathbf{a})\|$, with equality when $\mathbf{u}$ is parallel to $\nabla f(\mathbf{a})$. \qed

\section{The Chain Rule}

\begin{theorem}{Chain Rule for Vector Fields}
Let $\mathbf{f}: S \to \mathbb{R}^m$ be differentiable at $\mathbf{a} \in S \subseteq \mathbb{R}^n$, and $\mathbf{g}: T \to \mathbb{R}^p$ be differentiable at $\mathbf{b} = \mathbf{f}(\mathbf{a}) \in T \subseteq \mathbb{R}^m$. Then the composition $\mathbf{h} = \mathbf{g} \circ \mathbf{f}$ is differentiable at $\mathbf{a}$ and:
\[
\mathbf{h}'(\mathbf{a}) = \mathbf{g}'(\mathbf{f}(\mathbf{a})) \circ \mathbf{f}'(\mathbf{a})
\]
In matrix form: $D\mathbf{h}(\mathbf{a}) = D\mathbf{g}(\mathbf{f}(\mathbf{a})) \cdot D\mathbf{f}(\mathbf{a})$.
\end{theorem}

\section{Higher-Order Partial Derivatives}

\begin{definition}{Mixed Partial Derivatives}
For a scalar field $f$, the mixed partial derivative of order $k$ is:
\[
D_{i_1 i_2 \cdots i_k} f = D_{i_1}(D_{i_2}(\cdots(D_{i_k} f)\cdots))
\]
\end{definition}

\begin{theorem}{Equality of Mixed Partials}
If $f$ has continuous second partial derivatives in an open set containing $\mathbf{a}$, then:
\[
D_{i,j} f(\mathbf{a}) = D_{j,i} f(\mathbf{a})
\]
\end{theorem}

\section{Mean-Value Theorems}

\begin{theorem}{Mean-Value Theorem for Scalar Fields}
Let $f$ be differentiable on an open set containing the line segment from $\mathbf{a}$ to $\mathbf{a} + \mathbf{h}$. Then there exists $\theta \in (0, 1)$ such that:
\[
f(\mathbf{a} + \mathbf{h}) - f(\mathbf{a}) = \nabla f(\mathbf{a} + \theta \mathbf{h}) \cdot \mathbf{h}
\]
\end{theorem}

\begin{theorem}{Mean-Value Theorem for Vector Fields}
Let $\mathbf{f}$ be differentiable on an open set containing the line segment from $\mathbf{a}$ to $\mathbf{a} + \mathbf{h}$. Then for each component $f_i$:
\[
f_i(\mathbf{a} + \mathbf{h}) - f_i(\mathbf{a}) = \nabla f_i(\mathbf{a} + \theta_i \mathbf{h}) \cdot \mathbf{h}
\]
for some $\theta_i \in (0, 1)$.
\end{theorem}

\section{Taylor's Formula}

\begin{theorem}{Taylor's Theorem with Remainder}
Let $f$ have continuous partial derivatives up to order $k+1$ in a neighborhood of $\mathbf{a}$. Then:
\[
f(\mathbf{a} + \mathbf{h}) = \sum_{m=0}^{k} \frac{1}{m!}(\mathbf{h} \cdot \nabla)^m f(\mathbf{a}) + R_k(\mathbf{h})
\]
where the remainder $R_k(\mathbf{h}) = \frac{1}{(k+1)!}(\mathbf{h} \cdot \nabla)^{k+1} f(\mathbf{a} + \theta \mathbf{h})$ for some $\theta \in (0, 1)$.
\end{theorem}


\chapter{Applications of Differential Calculus}

\section{Partial Differential Equations}

\begin{definition}{Partial Differential Equation}
A partial differential equation (PDE) is an equation involving an unknown function of several variables and its partial derivatives. The order of a PDE is the order of the highest derivative appearing in the equation.
\end{definition}

\begin{theorem}{First-Order Linear PDE with Constant Coefficients}
The general solution of the first-order linear PDE:
\[
a\frac{\partial u}{\partial x} + b\frac{\partial u}{\partial y} = 0
\]
where $a$ and $b$ are constants (not both zero), is:
\[
u(x,y) = f(bx - ay)
\]
where $f$ is an arbitrary differentiable function.
\end{theorem}

\textit{Proof.} Introduce new variables $\xi = ax + by$ and $\eta = bx - ay$. By the chain rule:
\[
\frac{\partial u}{\partial x} = a\frac{\partial u}{\partial \xi} + b\frac{\partial u}{\partial \eta}, \quad \frac{\partial u}{\partial y} = b\frac{\partial u}{\partial \xi} - a\frac{\partial u}{\partial \eta}
\]
Substituting: $a(a\frac{\partial u}{\partial \xi} + b\frac{\partial u}{\partial \eta}) + b(b\frac{\partial u}{\partial \xi} - a\frac{\partial u}{\partial \eta}) = (a^2+b^2)\frac{\partial u}{\partial \xi} = 0$.

Thus $\frac{\partial u}{\partial \xi} = 0$, so $u$ depends only on $\eta = bx - ay$. \qed

\section{The One-Dimensional Wave Equation}

\begin{definition}{Wave Equation}
The one-dimensional wave equation is:
\[
\frac{\partial^2 u}{\partial t^2} = c^2 \frac{\partial^2 u}{\partial x^2}
\]
where $c > 0$ is the wave speed.
\end{definition}

\begin{theorem}{d'Alembert's Solution}
The general solution of the wave equation is:
\[
u(x,t) = f(x+ct) + g(x-ct)
\]
where $f$ and $g$ are arbitrary twice-differentiable functions.
\end{theorem}

\textit{Proof.} Using the factorization:
\[
\frac{\partial^2}{\partial t^2} - c^2\frac{\partial^2}{\partial x^2} = \left(\frac{\partial}{\partial t} - c\frac{\partial}{\partial x}\right)\left(\frac{\partial}{\partial t} + c\frac{\partial}{\partial x}\right)
\]

Let $v = \frac{\partial u}{\partial t} + c\frac{\partial u}{\partial x}$. Then $\frac{\partial v}{\partial t} - c\frac{\partial v}{\partial x} = 0$, giving $v = h(x+ct)$ for some function $h$.

Solving $\frac{\partial u}{\partial t} + c\frac{\partial u}{\partial x} = h(x+ct)$ yields $u = f(x+ct) + g(x-ct)$. \qed

\section{The Method of Characteristics}

\begin{theorem}{Characteristic Curves}
For the first-order PDE $a(x,y)\frac{\partial u}{\partial x} + b(x,y)\frac{\partial u}{\partial y} = c(x,y,u)$, the characteristic curves are solutions of:
\[
\frac{dx}{a(x,y)} = \frac{dy}{b(x,y)} = \frac{du}{c(x,y,u)}
\]
\end{theorem}

\section{Extrema of Scalar Fields}

\begin{definition}{Local Extrema}
A scalar field $f$ has a local maximum at $\mathbf{a}$ if $f(\mathbf{a}) \geq f(\mathbf{x})$ for all $\mathbf{x}$ in some neighborhood of $\mathbf{a}$. Similarly for local minimum.
\end{definition}

\begin{theorem}{Critical Points}
If $f$ has a local extremum at $\mathbf{a}$ and $\nabla f(\mathbf{a})$ exists, then $\nabla f(\mathbf{a}) = \mathbf{0}$.
\end{theorem}

\textit{Proof.} For any direction $\mathbf{u}$, the directional derivative $f'(\mathbf{a}; \mathbf{u}) = \nabla f(\mathbf{a}) \cdot \mathbf{u} = 0$ at an extremum. Thus $\nabla f(\mathbf{a}) = \mathbf{0}$. \qed

\begin{definition}{Hessian Matrix}
For a scalar field $f$ with continuous second partial derivatives, the Hessian matrix at $\mathbf{a}$ is:
\[
H(\mathbf{a}) = \begin{pmatrix}
D_{1,1}f(\mathbf{a}) & D_{1,2}f(\mathbf{a}) & \cdots & D_{1,n}f(\mathbf{a}) \\
D_{2,1}f(\mathbf{a}) & D_{2,2}f(\mathbf{a}) & \cdots & D_{2,n}f(\mathbf{a}) \\
\vdots & \vdots & \ddots & \vdots \\
D_{n,1}f(\mathbf{a}) & D_{n,2}f(\mathbf{a}) & \cdots & D_{n,n}f(\mathbf{a})
\end{pmatrix}
\]
\end{definition}

\begin{theorem}{Second Derivative Test}
Let $f$ have continuous second partial derivatives near a critical point $\mathbf{a}$ (where $\nabla f(\mathbf{a}) = \mathbf{0}$).
\begin{enumerate}
    \item[(a)] If $H(\mathbf{a})$ is positive definite, $f$ has a local minimum at $\mathbf{a}$.
    \item[(b)] If $H(\mathbf{a})$ is negative definite, $f$ has a local maximum at $\mathbf{a}$.
    \item[(c)] If $H(\mathbf{a})$ is indefinite, $f$ has a saddle point at $\mathbf{a}$.
\end{enumerate}
\end{theorem}

\section{Lagrange Multipliers}

\begin{theorem}{Lagrange Multipliers}
Let $f$ and $g$ be scalar fields with continuous partial derivatives. If $f$ has a local extremum at $\mathbf{a}$ subject to the constraint $g(\mathbf{a}) = 0$, and $\nabla g(\mathbf{a}) \neq \mathbf{0}$, then there exists a scalar $\lambda$ (called a Lagrange multiplier) such that:
\[
\nabla f(\mathbf{a}) = \lambda \nabla g(\mathbf{a})
\]
\end{theorem}

\textit{Proof.} Since $\nabla g(\mathbf{a}) \neq \mathbf{0}$, the level surface $g = 0$ has a tangent plane at $\mathbf{a}$. Any curve on this surface through $\mathbf{a}$ has tangent vector $\mathbf{v}$ satisfying $\nabla g(\mathbf{a}) \cdot \mathbf{v} = 0$. For $f$ to have an extremum on the constraint surface, its directional derivative must vanish in all tangent directions, so $\nabla f(\mathbf{a}) \cdot \mathbf{v} = 0$. Thus $\nabla f(\mathbf{a})$ is parallel to $\nabla g(\mathbf{a})$. \qed

\section{The Implicit Function Theorem}

\begin{theorem}{Implicit Function Theorem}
Let $F$ be a scalar field with continuous partial derivatives in a neighborhood of $(\mathbf{a}, b) \in \mathbb{R}^{n+1}$, where $F(\mathbf{a}, b) = 0$ and $D_{n+1}F(\mathbf{a}, b) \neq 0$. Then there exists a neighborhood of $\mathbf{a}$ and a unique function $f$ defined on this neighborhood such that $f(\mathbf{a}) = b$ and $F(\mathbf{x}, f(\mathbf{x})) = 0$ for all $\mathbf{x}$ in this neighborhood. Moreover, $f$ has continuous partial derivatives given by:
\[
D_j f(\mathbf{x}) = -\frac{D_j F(\mathbf{x}, f(\mathbf{x}))}{D_{n+1} F(\mathbf{x}, f(\mathbf{x}))}
\]
\end{theorem}

\section{Extreme-Value Theorem for Scalar Fields}

\begin{theorem}{Extreme-Value Theorem}
If $f$ is continuous on a closed interval $[a, b]$ in $\mathbb{R}^n$, then there exist points $c$ and $d$ in $[a, b]$ such that:
\[
f(c) = \sup f \quad \text{and} \quad f(d) = \inf f
\]
\end{theorem}

\textit{Proof.} It suffices to prove $f$ attains its supremum. Let $M = \sup f$. If $f(x) < M$ for all $x$, define $g(x) = \frac{1}{M-f(x)}$. Then $g$ is continuous and unbounded, contradicting the boundedness of continuous functions on compact sets. Thus $f(c) = M$ for some $c$. \qed

\begin{theorem}{Uniform Continuity (Small-Span Theorem)}
Let $f$ be continuous on a closed interval $[a, b]$ in $\mathbb{R}^n$. Then for every $\epsilon > 0$ there is a partition of $[a, b]$ into a finite number of subintervals such that the span (oscillation) of $f$ in each subinterval is less than $\epsilon$.
\end{theorem}


\chapter{Line Integrals}

\section{Introduction}

\begin{definition}{Line Integral of a Scalar Field}
Let $f$ be a scalar field defined on a curve $C$ parameterized by $\mathbf{r}(t)$ for $a \leq t \leq b$. The line integral of $f$ along $C$ with respect to arc length is:
\[
\int_C f\, ds = \int_a^b f(\mathbf{r}(t)) \|\mathbf{r}'(t)\|\, dt
\]
\end{definition}

\begin{definition}{Line Integral of a Vector Field}
Let $\mathbf{F}$ be a vector field defined on a curve $C$ parameterized by $\mathbf{r}(t)$ for $a \leq t \leq b$. The line integral of $\mathbf{F}$ along $C$ is:
\[
\int_C \mathbf{F} \cdot d\mathbf{r} = \int_a^b \mathbf{F}(\mathbf{r}(t)) \cdot \mathbf{r}'(t)\, dt
\]
\end{definition}

\section{Basic Properties of Line Integrals}

\begin{theorem}{Additivity}
If $C$ is the union of two curves $C_1$ and $C_2$, then:
\[
\int_C \mathbf{F} \cdot d\mathbf{r} = \int_{C_1} \mathbf{F} \cdot d\mathbf{r} + \int_{C_2} \mathbf{F} \cdot d\mathbf{r}
\]
\end{theorem}

\begin{theorem}{Reversal of Direction}
If $-C$ denotes the curve $C$ traversed in the opposite direction, then:
\[
\int_{-C} \mathbf{F} \cdot d\mathbf{r} = -\int_C \mathbf{F} \cdot d\mathbf{r}
\]
\end{theorem}

\begin{theorem}{Bounds for Line Integrals}
If $|\mathbf{F}(\mathbf{x})| \leq M$ for all $\mathbf{x}$ on $C$, then:
\[
\left|\int_C \mathbf{F} \cdot d\mathbf{r}\right| \leq ML
\]
where $L$ is the length of $C$.
\end{theorem}

\section{Independence of Path}

\begin{definition}{Conservative Vector Field}
A vector field $\mathbf{F}$ is called conservative if $\int_C \mathbf{F} \cdot d\mathbf{r}$ depends only on the endpoints of $C$, not on the path taken.
\end{definition}

\begin{theorem}{Conservative Fields and Potential Functions}
A continuous vector field $\mathbf{F}$ on an open connected set $S$ is conservative if and only if there exists a scalar field $\varphi$ (called a potential function) such that $\mathbf{F} = \nabla \varphi$ on $S$.
\end{theorem}

\textit{Proof.} $(\Rightarrow)$ If $\mathbf{F} = \nabla \varphi$, then by the fundamental theorem:
\[
\int_C \mathbf{F} \cdot d\mathbf{r} = \varphi(\mathbf{b}) - \varphi(\mathbf{a})
\]
which depends only on endpoints.

$(\Leftarrow)$ Define $\varphi(\mathbf{x}) = \int_{\mathbf{a}}^{\mathbf{x}} \mathbf{F} \cdot d\mathbf{r}$ along any path from fixed $\mathbf{a}$ to $\mathbf{x}$. By path independence, $\varphi$ is well-defined, and one can show $\nabla \varphi = \mathbf{F}$. \qed

\section{Green's Theorem}

\begin{theorem}{Green's Theorem}
Let $C$ be a positively oriented, piecewise smooth, simple closed curve in the plane bounding a region $D$. If $P$ and $Q$ have continuous partial derivatives on an open region containing $D$, then:
\[
\oint_C (P\, dx + Q\, dy) = \iint_D \left(\frac{\partial Q}{\partial x} - \frac{\partial P}{\partial y}\right) dA
\]
\end{theorem}

\textit{Proof.} The proof involves computing both sides for simple regions. For a region $D$ described as $a \leq x \leq b$, $g_1(x) \leq y \leq g_2(x)$:
\[
\iint_D \frac{\partial P}{\partial y}\, dA = \int_a^b [P(x, g_2(x)) - P(x, g_1(x))]\, dx = -\oint_C P\, dx
\]
Similarly for $Q$, establishing the theorem. \qed

\section{Potential Functions on Convex Sets}

\begin{theorem}{Differentiation Under the Integral Sign}
Let $\psi(\mathbf{x}, t)$ and $\frac{\partial \psi}{\partial x_k}$ be continuous on $S \times J$ where $S$ is an open set and $J = [a, b]$. Then:
\[
\frac{\partial}{\partial x_k} \int_a^b \psi(\mathbf{x}, t)\, dt = \int_a^b \frac{\partial \psi}{\partial x_k}(\mathbf{x}, t)\, dt
\]
\end{theorem}

\begin{theorem}{Condition for a Gradient Field}
Let $\mathbf{f} = (f_1, \ldots, f_n)$ be a continuously differentiable vector field on an open convex set $S$ in $\mathbb{R}^n$. Then $\mathbf{f}$ is a gradient on $S$ if and only if:
\[
D_k f_j(\mathbf{x}) = D_j f_k(\mathbf{x})
\]
for all $\mathbf{x} \in S$ and all $k, j = 1, 2, \ldots, n$.
\end{theorem}

\textit{Proof.} If $\mathbf{f} = \nabla \varphi$, then $D_k f_j = D_k D_j \varphi = D_j D_k \varphi = D_j f_k$ by equality of mixed partials.

Conversely, define $\varphi(\mathbf{x}) = \int_0^1 \mathbf{f}(t\mathbf{x}) \cdot \mathbf{x}\, dt$. Then:
\[
D_k \varphi(\mathbf{x}) = \int_0^1 [f_k(t\mathbf{x}) + t\nabla f_k(t\mathbf{x}) \cdot \mathbf{x}]\, dt
\]
Integrating by parts shows $D_k \varphi = f_k$, so $\nabla \varphi = \mathbf{f}$. \qed



\chapter{Multiple Integrals}

\section{Introduction}

\begin{definition}{Double Integral}
Let $f$ be a scalar field defined on a rectangle $R = [a, b] \times [c, d]$ in $\mathbb{R}^2$. The double integral of $f$ over $R$ is defined as the limit of Riemann sums:
\[
\iint_R f(x, y)\, dA = \lim_{m, n \to \infty} \sum_{i=1}^{m} \sum_{j=1}^{n} f(x_i^*, y_j^*) \Delta A_{ij}
\]
where $\Delta A_{ij} = \Delta x_i \Delta y_j$ is the area of the subrectangle, and $(x_i^*, y_j^*)$ is a sample point.
\end{definition}

\begin{definition}{Iterated Integral}
For a function $f$ continuous on $R = [a, b] \times [c, d]$, the iterated integral is:
\[
\int_a^b \left[ \int_c^d f(x, y)\, dy \right] dx
\]
where the inner integral is computed first (treating $x$ as constant), followed by the outer integral.
\end{definition}

\section{Fubini's Theorem}

\begin{theorem}{Fubini's Theorem}
If $f$ is continuous on the rectangle $R = [a, b] \times [c, d]$, then:
\[
\iint_R f(x, y)\, dA = \int_a^b \int_c^d f(x, y)\, dy\, dx = \int_c^d \int_a^b f(x, y)\, dx\, dy
\]
\end{theorem}

\textit{Proof.} The proof involves showing that both iterated integrals equal the double integral. This follows from the uniform continuity of $f$ on the compact set $R$ and the fact that upper and lower sums converge to the same limit. \qed

\section{Integration Over General Regions}

\begin{definition}{Type I and Type II Regions}
A region $D$ is of Type I if it can be described as:
\[
D = \{(x, y) : a \leq x \leq b, \ g_1(x) \leq y \leq g_2(x)\}
\]
A region $D$ is of Type II if it can be described as:
\[
D = \{(x, y) : c \leq y \leq d, \ h_1(y) \leq x \leq h_2(y)\}
\]
\end{definition}

\begin{theorem}{Double Integral Over General Regions}
If $f$ is continuous on a Type I region $D$, then:
\[
\iint_D f(x, y)\, dA = \int_a^b \int_{g_1(x)}^{g_2(x)} f(x, y)\, dy\, dx
\]
Similarly for Type II regions with the order of integration reversed.
\end{theorem}

\section{Change of Variables}

\begin{theorem}{Change of Variables Formula}
Let $T$ be a one-to-one transformation from a region $S$ in the $uv$-plane to a region $R$ in the $xy$-plane, given by $x = x(u, v)$ and $y = y(u, v)$. If $f$ is continuous on $R$ and the Jacobian determinant:
\[
J(u, v) = \frac{\partial(x, y)}{\partial(u, v)} = \begin{vmatrix} \frac{\partial x}{\partial u} & \frac{\partial x}{\partial v} \\ \frac{\partial y}{\partial u} & \frac{\partial y}{\partial v} \end{vmatrix}
\]
is nonzero and continuous, then:
\[
\iint_R f(x, y)\, dx\, dy = \iint_S f(x(u, v), y(u, v)) |J(u, v)|\, du\, dv
\]
\end{theorem}

\begin{theorem}{Polar Coordinates}
For polar coordinates $x = r\cos\theta$, $y = r\sin\theta$, the Jacobian is:
\[
J(r, \theta) = \begin{vmatrix} \cos\theta & -r\sin\theta \\ \sin\theta & r\cos\theta \end{vmatrix} = r
\]
Therefore:
\[
\iint_R f(x, y)\, dA = \iint_S f(r\cos\theta, r\sin\theta) \cdot r\, dr\, d\theta
\]
\end{theorem}

\section{Triple Integrals}

\begin{definition}{Triple Integral}
For a function $f$ defined on a rectangular box $B = [a, b] \times [c, d] \times [r, s]$:
\[
\iiint_B f(x, y, z)\, dV = \lim_{l, m, n \to \infty} \sum_{i=1}^{l} \sum_{j=1}^{m} \sum_{k=1}^{n} f(x_i^*, y_j^*, z_k^*) \Delta V_{ijk}
\]
where $\Delta V_{ijk} = \Delta x_i \Delta y_j \Delta z_k$.
\end{definition}

\begin{theorem}{Fubini's Theorem for Triple Integrals}
If $f$ is continuous on $B$, then:
\[
\iiint_B f(x, y, z)\, dV = \int_a^b \int_c^d \int_r^s f(x, y, z)\, dz\, dy\, dx
\]
and similarly for any other order of integration.
\end{theorem}

\section{Spherical and Cylindrical Coordinates}

\begin{theorem}{Cylindrical Coordinates}
For $x = r\cos\theta$, $y = r\sin\theta$, $z = z$, the volume element is:
\[
dV = r\, dr\, d\theta\, dz
\]
\end{theorem}

\begin{theorem}{Spherical Coordinates}
For $x = \rho\sin\phi\cos\theta$, $y = \rho\sin\phi\sin\theta$, $z = \rho\cos\phi$, the volume element is:
\[
dV = \rho^2 \sin\phi\, d\rho\, d\phi\, d\theta
\]
\end{theorem}

\section{Applications}

\begin{theorem}{Volume}
The volume of a solid region $E$ is given by:
\[
V = \iiint_E dV
\]
\end{theorem}

\begin{theorem}{Center of Mass}
For a solid region $E$ with density function $\rho(x, y, z)$, the center of mass is:
\[
(\bar{x}, \bar{y}, \bar{z}) = \left(\frac{M_{yz}}{m}, \frac{M_{xz}}{m}, \frac{M_{xy}}{m}\right)
\]
where $m = \iiint_E \rho\, dV$ and $M_{yz} = \iiint_E x\rho\, dV$, etc.
\end{theorem}



\chapter{Surface Integrals}

\section{Parametric Representation of a Surface}

\begin{definition}{Parametric Surface}
A parametric surface is defined by a vector function $\mathbf{r}(u, v)$ where $(u, v)$ varies over a region $D$ in the $uv$-plane:
\[
\mathbf{r}(u, v) = X(u, v)\mathbf{i} + Y(u, v)\mathbf{j} + Z(u, v)\mathbf{k}, \quad (u, v) \in D
\]
\end{definition}

\section{The Fundamental Vector Product}

\begin{definition}{Tangent Vectors}
For a parametric surface $\mathbf{r}(u, v)$, the tangent vectors are:
\[
\frac{\partial \mathbf{r}}{\partial u} = \frac{\partial X}{\partial u}\mathbf{i} + \frac{\partial Y}{\partial u}\mathbf{j} + \frac{\partial Z}{\partial u}\mathbf{k}
\]
\[
\frac{\partial \mathbf{r}}{\partial v} = \frac{\partial X}{\partial v}\mathbf{i} + \frac{\partial Y}{\partial v}\mathbf{j} + \frac{\partial Z}{\partial v}\mathbf{k}
\]
\end{definition}

\begin{definition}{Fundamental Vector Product}
The fundamental vector product (or normal vector) is:
\[
\frac{\partial \mathbf{r}}{\partial u} \times \frac{\partial \mathbf{r}}{\partial v} = \begin{vmatrix} \mathbf{i} & \mathbf{j} & \mathbf{k} \\ \frac{\partial X}{\partial u} & \frac{\partial Y}{\partial u} & \frac{\partial Z}{\partial u} \\ \frac{\partial X}{\partial v} & \frac{\partial Y}{\partial v} & \frac{\partial Z}{\partial v} \end{vmatrix}
\]
Its magnitude is: $\left\|\frac{\partial \mathbf{r}}{\partial u} \times \frac{\partial \mathbf{r}}{\partial v}\right\| = \sqrt{EG - F^2}$
where $E = \|\frac{\partial \mathbf{r}}{\partial u}\|^2$, $F = \frac{\partial \mathbf{r}}{\partial u} \cdot \frac{\partial \mathbf{r}}{\partial v}$, $G = \|\frac{\partial \mathbf{r}}{\partial v}\|^2$.
\end{definition}

\section{Surface Area}

\begin{definition}{Surface Area}
The surface area of a parametric surface $S$ is:
\[
A(S) = \iint_D \left\|\frac{\partial \mathbf{r}}{\partial u} \times \frac{\partial \mathbf{r}}{\partial v}\right\| \, du\, dv
\]
\end{definition}

\begin{theorem}{Surface Area Formula}
For a surface $z = f(x, y)$ over a region $R$ in the $xy$-plane:
\[
A(S) = \iint_R \sqrt{1 + \left(\frac{\partial f}{\partial x}\right)^2 + \left(\frac{\partial f}{\partial y}\right)^2} \, dA
\]
\end{theorem}

\section{Surface Integrals of Scalar Fields}

\begin{definition}{Surface Integral of a Scalar Field}
Let $f$ be a scalar field defined on a surface $S$ parameterized by $\mathbf{r}(u, v)$ for $(u, v) \in D$. The surface integral is:
\[
\iint_S f\, dS = \iint_D f(\mathbf{r}(u, v)) \left\|\frac{\partial \mathbf{r}}{\partial u} \times \frac{\partial \mathbf{r}}{\partial v}\right\| \, du\, dv
\]
\end{definition}

\section{Surface Integrals of Vector Fields}

\begin{definition}{Surface Integral of a Vector Field (Flux)}
Let $\mathbf{F}$ be a vector field defined on an orientable surface $S$ with unit normal $\mathbf{n}$. The surface integral (flux) is:
\[
\iint_S \mathbf{F} \cdot d\mathbf{S} = \iint_S \mathbf{F} \cdot \mathbf{n}\, dS = \iint_D \mathbf{F}(\mathbf{r}(u, v)) \cdot \left(\frac{\partial \mathbf{r}}{\partial u} \times \frac{\partial \mathbf{r}}{\partial v}\right) \, du\, dv
\]
\end{definition}

\section{Stokes' Theorem}

\begin{theorem}{Stokes' Theorem}
Let $S$ be an oriented piecewise smooth surface bounded by a simple, closed, piecewise smooth boundary curve $\partial S$ with positive orientation. If $\mathbf{F}$ is a vector field with continuous partial derivatives on an open region containing $S$, then:
\[
\oint_{\partial S} \mathbf{F} \cdot d\mathbf{r} = \iint_S (\nabla \times \mathbf{F}) \cdot d\mathbf{S}
\]
\end{theorem}

\textit{Proof.} (Outline) Parameterize $S$ by $\mathbf{r}(u, v)$ and use Green's theorem in the $uv$-plane. The line integral around $\partial S$ transforms to a line integral around $\partial D$, which equals the double integral of the curl over $D$, which maps back to the surface integral of $\nabla \times \mathbf{F}$ over $S$. \qed

\section{The Divergence Theorem}

\begin{theorem}{Divergence Theorem (Gauss's Theorem)}
Let $V$ be a solid region in $\mathbb{R}^3$ bounded by a closed surface $\partial V$ with outward orientation. If $\mathbf{F}$ is a vector field with continuous partial derivatives on an open region containing $V$, then:
\[
\iiint_V (\nabla \cdot \mathbf{F})\, dV = \iint_{\partial V} \mathbf{F} \cdot d\mathbf{S}
\]
\end{theorem}

\textit{Proof.} The proof involves showing that the volume integral of each component equals the corresponding surface integral. For the $z$-component:
\[
\iiint_V \frac{\partial R}{\partial z}\, dV = \iint_D [R(x, y, h_2(x, y)) - R(x, y, h_1(x, y))]\, dA = \iint_{\partial V} R\mathbf{k} \cdot \mathbf{n}\, dS
\]
where $z = h_1(x, y)$ and $z = h_2(x, y)$ describe the bottom and top surfaces. Similar arguments hold for the other components. \qed

\section{Conservative Vector Fields}

\begin{theorem}{Equivalence of Conservative Conditions}
For a vector field $\mathbf{F}$ with continuous partial derivatives on a simply connected domain in $\mathbb{R}^3$, the following are equivalent:
\begin{enumerate}
    \item[(a)] $\mathbf{F}$ is conservative (i.e., $\mathbf{F} = \nabla \varphi$ for some potential $\varphi$)
    \item[(b)] $\oint_C \mathbf{F} \cdot d\mathbf{r} = 0$ for every closed curve $C$
    \item[(c)] $\nabla \times \mathbf{F} = \mathbf{0}$
\end{enumerate}
\end{theorem}



\chapter{Set Functions and Elementary Probability}

\section{Introduction}

\begin{definition}{Set Function}
A set function is a function defined on a collection of sets. If $\mathcal{A}$ is a collection of subsets of a set $S$, a set function $f$ assigns to each set $A \in \mathcal{A}$ a real number $f(A)$.
\end{definition}

\section{Finitely Additive Set Functions}

\begin{definition}{Finitely Additive}
A set function $f$ defined on a collection $\mathcal{A}$ of sets is finitely additive if:
\[
f(A \cup B) = f(A) + f(B)
\]
whenever $A, B \in \mathcal{A}$, $A \cap B = \emptyset$, and $A \cup B \in \mathcal{A}$.
\end{definition}

\begin{theorem}{Finite Additivity}
If $f$ is finitely additive, then for any finite collection of disjoint sets $A_1, \ldots, A_n$ in $\mathcal{A}$ with $\bigcup_{i=1}^{n} A_i \in \mathcal{A}$:
\[
f\left(\bigcup_{i=1}^{n} A_i\right) = \sum_{i=1}^{n} f(A_i)
\]
\end{theorem}

\section{Probability Measures}

\begin{definition}{Probability Space}
A probability space is a triple $(\Omega, \mathcal{F}, P)$ where:
\begin{enumerate}
    \item $\Omega$ is a sample space (set of all possible outcomes)
    \item $\mathcal{F}$ is a $\sigma$-algebra of subsets of $\Omega$ (events)
    \item $P: \mathcal{F} \to [0, 1]$ is a probability measure satisfying:
    \begin{enumerate}
        \item[(a)] $P(\Omega) = 1$
        \item[(b)] Countable additivity: For disjoint events $A_1, A_2, \ldots$:
        \[
        P\left(\bigcup_{i=1}^{\infty} A_i\right) = \sum_{i=1}^{\infty} P(A_i)
        \]
    \end{enumerate}
\end{enumerate}
\end{definition}

\begin{theorem}{Properties of Probability}
For any events $A$ and $B$:
\begin{enumerate}
    \item[(a)] $P(\emptyset) = 0$
    \item[(b)] $P(A^c) = 1 - P(A)$
    \item[(c)] If $A \subseteq B$, then $P(A) \leq P(B)$
    \item[(d)] $P(A \cup B) = P(A) + P(B) - P(A \cap B)$
\end{enumerate}
\end{theorem}

\textit{Proof.} (a) Since $\Omega = \Omega \cup \emptyset$ and $\Omega \cap \emptyset = \emptyset$, we have $P(\Omega) = P(\Omega) + P(\emptyset)$, so $P(\emptyset) = 0$.

(b) Since $\Omega = A \cup A^c$ and $A \cap A^c = \emptyset$, we have $1 = P(A) + P(A^c)$. \qed

\section{Countably Additive Set Functions}

\begin{definition}{Countably Additive}
A set function $f$ is countably additive (or $\sigma$-additive) if for any countable collection of disjoint sets $\{A_i\}$ with $\bigcup_{i=1}^{\infty} A_i \in \mathcal{A}$:
\[
f\left(\bigcup_{i=1}^{\infty} A_i\right) = \sum_{i=1}^{\infty} f(A_i)
\]
\end{definition}

\begin{theorem}{Continuity of Probability}
If $\{A_n\}$ is an increasing sequence of events ($A_n \subseteq A_{n+1}$), then:
\[
P\left(\bigcup_{n=1}^{\infty} A_n\right) = \lim_{n \to \infty} P(A_n)
\]
If $\{A_n\}$ is a decreasing sequence ($A_n \supseteq A_{n+1}$), then:
\[
P\left(\bigcap_{n=1}^{\infty} A_n\right) = \lim_{n \to \infty} P(A_n)
\]
\end{theorem}

\section{The Monte Carlo Method}

\begin{theorem}{Monte Carlo Integration}
Let $f$ be integrable on $[0, 1]$. If $X_1, X_2, \ldots, X_n$ are independent uniform random variables on $[0, 1]$, then:
\[
\frac{1}{n}\sum_{i=1}^{n} f(X_i) \to \int_0^1 f(x)\, dx \quad \text{(in probability as $n \to \infty$)}
\]
\end{theorem}

\textit{Proof.} By the law of large numbers, the sample mean converges to the expected value $E[f(X)] = \int_0^1 f(x)\, dx$. \qed

\section{The Central Limit Theorem}

\begin{definition}{Independence}
Events $A_1, A_2, \ldots, A_n$ are independent if for any subset $\{i_1, \ldots, i_k\} \subseteq \{1, \ldots, n\}$:
\[
P(A_{i_1} \cap \cdots \cap A_{i_k}) = P(A_{i_1}) \cdots P(A_{i_k})
\]
\end{definition}

\begin{definition}{Random Variable}
A random variable $X$ on a probability space $(\Omega, \mathcal{F}, P)$ is a measurable function $X: \Omega \to \mathbb{R}$.
\end{definition}

\begin{definition}{Expected Value and Variance}
For a random variable $X$ with probability density function $f$:
\[
E[X] = \int_{-\infty}^{\infty} x f(x)\, dx, \quad \text{Var}(X) = E[(X - E[X])^2] = E[X^2] - (E[X])^2
\]
\end{definition}

\begin{theorem}{Central Limit Theorem}
Let $X_1, X_2, \ldots$ be independent and identically distributed random variables with mean $\mu$ and variance $\sigma^2 < \infty$. Then:
\[
\frac{\sum_{i=1}^{n} X_i - n\mu}{\sigma\sqrt{n}} \xrightarrow{d} N(0, 1) \quad \text{as } n \to \infty
\]
That is, the standardized sum converges in distribution to the standard normal distribution.
\end{theorem}

\textit{Proof.} (Sketch) Using characteristic functions, one shows that the characteristic function of the standardized sum converges to $e^{-t^2/2}$, which is the characteristic function of the standard normal distribution. By the continuity theorem for characteristic functions, this implies convergence in distribution. \qed

\section{Conditional Probability}

\begin{definition}{Conditional Probability}
For events $A$ and $B$ with $P(B) > 0$, the conditional probability of $A$ given $B$ is:
\[
P(A | B) = \frac{P(A \cap B)}{P(B)}
\]
\end{definition}

\begin{theorem}{Law of Total Probability}
If $\{B_i\}$ is a partition of $\Omega$ (i.e., disjoint and $\bigcup_i B_i = \Omega$) with $P(B_i) > 0$, then:
\[
P(A) = \sum_i P(A | B_i) P(B_i)
\]
\end{theorem}

\begin{theorem}{Bayes' Theorem}
If $\{B_i\}$ is a partition of $\Omega$ with $P(B_i) > 0$, then:
\[
P(B_j | A) = \frac{P(A | B_j) P(B_j)}{\sum_i P(A | B_i) P(B_i)}
\]
\end{theorem}



\chapter{Introduction to Numerical Analysis}

\section{Introduction}

\begin{definition}{Numerical Analysis}
Numerical analysis is the study of algorithms for solving problems in continuous mathematics, including approximation of functions, numerical integration, and solution of differential equations.
\end{definition}

\section{Polynomial Interpolation}

\begin{definition}{Lagrange Interpolation}
Given $n+1$ distinct points $(x_0, y_0), (x_1, y_1), \ldots, (x_n, y_n)$, the Lagrange interpolating polynomial is:
\[
P_n(x) = \sum_{i=0}^{n} y_i L_i(x)
\]
where $L_i(x) = \prod_{j \neq i} \frac{x - x_j}{x_i - x_j}$ are the Lagrange basis polynomials.
\end{definition}

\begin{theorem}{Existence and Uniqueness of Interpolating Polynomial}
There exists a unique polynomial of degree at most $n$ that passes through $n+1$ distinct points with distinct $x$-coordinates.
\end{theorem}

\textit{Proof.} The Lagrange formula constructs such a polynomial. For uniqueness, suppose $P$ and $Q$ both interpolate the points. Then $P-Q$ is a polynomial of degree at most $n$ with $n+1$ roots, so $P-Q = 0$. \qed

\section{Numerical Integration}

\begin{theorem}{Trapezoidal Rule}
For a function $f$ with continuous second derivative on $[a, b]$:
\[
\int_a^b f(x)\, dx \approx \frac{b-a}{2}[f(a) + f(b)]
\]
with error term $-\frac{(b-a)^3}{12}f''(\xi)$ for some $\xi \in (a, b)$.
\end{theorem}

\begin{theorem}{Simpson's Rule}
For a function $f$ with continuous fourth derivative on $[a, b]$:
\[
\int_a^b f(x)\, dx \approx \frac{b-a}{6}\left[f(a) + 4f\left(\frac{a+b}{2}\right) + f(b)\right]
\]
with error term $-\frac{(b-a)^5}{2880}f^{(4)}(\xi)$ for some $\xi \in (a, b)$.
\end{theorem}

\section{Newton's Method}

\begin{definition}{Newton's Method}
For finding roots of $f(x) = 0$, Newton's method iterates:
\[
x_{n+1} = x_n - \frac{f(x_n)}{f'(x_n)}
\]
\end{definition}

\begin{theorem}{Convergence of Newton's Method}
If $f$ has a continuous second derivative near a simple root $\alpha$ (where $f(\alpha) = 0$ and $f'(\alpha) \neq 0$), and if the initial guess $x_0$ is sufficiently close to $\alpha$, then Newton's method converges quadratically:
\[
|x_{n+1} - \alpha| \leq C|x_n - \alpha|^2
\]
for some constant $C$.
\end{theorem}

\textit{Proof.} Using Taylor expansion about $x_n$:
\[
0 = f(\alpha) = f(x_n) + f'(x_n)(\alpha - x_n) + \frac{f''(\xi)}{2}(\alpha - x_n)^2
\]
Rearranging and using the iteration formula gives the quadratic convergence. \qed

\section{Numerical Solution of Differential Equations}

\begin{definition}{Euler's Method}
For the initial value problem $y' = f(x, y)$, $y(x_0) = y_0$, Euler's method approximates:
\[
y_{n+1} = y_n + h f(x_n, y_n)
\]
where $h$ is the step size.
\end{definition}

\begin{theorem}{Error of Euler's Method}
If $f$ satisfies a Lipschitz condition and the exact solution $y(x)$ has continuous second derivative, then the error of Euler's method satisfies:
\[
|y(x_n) - y_n| \leq \frac{hM}{2L}(e^{L(x_n-x_0)} - 1)
\]
where $M$ bounds $|y''|$ and $L$ is the Lipschitz constant.
\end{theorem}

\begin{definition}{Runge-Kutta Methods}
The fourth-order Runge-Kutta method (RK4) for $y' = f(x, y)$ is:
\[
y_{n+1} = y_n + \frac{h}{6}(k_1 + 2k_2 + 2k_3 + k_4)
\]
where:
\begin{align*}
k_1 &= f(x_n, y_n) \\
k_2 &= f(x_n + \frac{h}{2}, y_n + \frac{h}{2}k_1) \\
k_3 &= f(x_n + \frac{h}{2}, y_n + \frac{h}{2}k_2) \\
k_4 &= f(x_n + h, y_n + hk_3)
\end{align*}
\end{definition}

\begin{theorem}{Convergence of RK4}
The fourth-order Runge-Kutta method has local truncation error $O(h^5)$ and global error $O(h^4)$.
\end{theorem}

\section{Fixed-Point Iteration}

\begin{definition}{Fixed Point}
A point $\alpha$ is a fixed point of a function $g$ if $g(\alpha) = \alpha$.
\end{definition}

\begin{theorem}{Contraction Mapping Theorem}
If $g$ maps $[a, b]$ into itself and satisfies $|g(x) - g(y)| \leq L|x - y|$ for all $x, y \in [a, b]$ with $L < 1$ (Lipschitz condition with constant $L < 1$), then:
\begin{enumerate}
    \item[(a)] There exists a unique fixed point $\alpha \in [a, b]$
    \item[(b)] For any $x_0 \in [a, b]$, the iteration $x_{n+1} = g(x_n)$ converges to $\alpha$
    \item[(c)] $|x_n - \alpha| \leq \frac{L^n}{1-L}|x_1 - x_0|$
\end{enumerate}
\end{theorem}

\textit{Proof.} (a) Existence follows from the intermediate value theorem applied to $g(x) - x$. Uniqueness: if $\alpha$ and $\beta$ are fixed points, then $|\alpha - \beta| = |g(\alpha) - g(\beta)| \leq L|\alpha - \beta|$, so $\alpha = \beta$ since $L < 1$.

(b) $|x_{n+1} - \alpha| = |g(x_n) - g(\alpha)| \leq L|x_n - \alpha| \leq \cdots \leq L^{n+1}|x_0 - \alpha| \to 0$. \qed



\chapter{Fourier Series and Fourier Integrals}

\section{Introduction}

\begin{definition}{Fourier Series}
For a periodic function $f$ with period $2\pi$, the Fourier series is:
\[
f(x) \sim \frac{a_0}{2} + \sum_{n=1}^{\infty} (a_n \cos nx + b_n \sin nx)
\]
where the Fourier coefficients are:
\[
a_n = \frac{1}{\pi} \int_{-\pi}^{\pi} f(x) \cos nx \, dx, \quad b_n = \frac{1}{\pi} \int_{-\pi}^{\pi} f(x) \sin nx \, dx
\]
\end{definition}

\section{Orthogonality Relations}

\begin{theorem}{Orthogonality of Trigonometric Functions}
\begin{align*}
\int_{-\pi}^{\pi} \cos mx \cos nx \, dx &= \begin{cases} 0 & m \neq n \\ \pi & m = n \neq 0 \\ 2\pi & m = n = 0 \end{cases} \\
\int_{-\pi}^{\pi} \sin mx \sin nx \, dx &= \begin{cases} 0 & m \neq n \\ \pi & m = n \end{cases} \\
\int_{-\pi}^{\pi} \cos mx \sin nx \, dx &= 0 \quad \text{for all } m, n
\end{align*}
\end{theorem}

\textit{Proof.} Using product-to-sum formulas:
$\cos mx \cos nx = \frac{1}{2}[\cos(m+n)x + \cos(m-n)x]$. Integrating over $[-\pi, \pi]$ gives the result. \qed

\section{Convergence of Fourier Series}

\begin{theorem}{Pointwise Convergence}
If $f$ is piecewise smooth on $[-\pi, \pi]$ with period $2\pi$, then the Fourier series converges to:
\[
\frac{f(x^+) + f(x^-)}{2}
\]
at every point $x$, where $f(x^+)$ and $f(x^-)$ are the right and left limits.
\end{theorem}

\begin{theorem}{Uniform Convergence}
If $f$ is continuous, piecewise smooth, and $f(-\pi) = f(\pi)$, then the Fourier series converges uniformly to $f$.
\end{theorem}

\section{Parseval's Identity}

\begin{theorem}{Parseval's Identity}
If $f$ has Fourier coefficients $a_n$ and $b_n$, then:
\[
\frac{1}{\pi} \int_{-\pi}^{\pi} |f(x)|^2 \, dx = \frac{a_0^2}{2} + \sum_{n=1}^{\infty} (a_n^2 + b_n^2)
\]
\end{theorem}

\textit{Proof.} Multiply the Fourier series by $f(x)$ and integrate term by term using orthogonality. \qed

\section{Complex Form of Fourier Series}

\begin{definition}{Complex Fourier Series}
The complex Fourier series is:
\[
f(x) \sim \sum_{n=-\infty}^{\infty} c_n e^{inx}
\]
where $c_n = \frac{1}{2\pi} \int_{-\pi}^{\pi} f(x) e^{-inx} \, dx$.
\end{definition}

\begin{theorem}{Relation Between Real and Complex Coefficients}
\[
c_0 = \frac{a_0}{2}, \quad c_n = \frac{a_n - ib_n}{2}, \quad c_{-n} = \frac{a_n + ib_n}{2} \quad (n \geq 1)
\]
\end{theorem}

\section{Fourier Integral}

\begin{definition}{Fourier Transform}
For a non-periodic function $f$ with $\int_{-\infty}^{\infty} |f(x)| \, dx < \infty$, the Fourier transform is:
\[
\hat{f}(\omega) = \frac{1}{\sqrt{2\pi}} \int_{-\infty}^{\infty} f(x) e^{-i\omega x} \, dx
\]
\end{definition}

\begin{theorem}{Fourier Inversion Theorem}
If $f$ and $\hat{f}$ are integrable, then:
\[
f(x) = \frac{1}{\sqrt{2\pi}} \int_{-\infty}^{\infty} \hat{f}(\omega) e^{i\omega x} \, d\omega
\]
\end{theorem}

\section{Properties of Fourier Transforms}

\begin{theorem}{Linearity}
$\widehat{(af + bg)} = a\hat{f} + b\hat{g}$
\end{theorem}

\begin{theorem}{Differentiation}
If $f$ and $f'$ are integrable and $f(x) \to 0$ as $|x| \to \infty$, then:
\[
\widehat{f'(\omega)} = i\omega \hat{f}(\omega)
\]
\end{theorem}

\begin{theorem}{Convolution Theorem}
The Fourier transform of the convolution $(f * g)(x) = \int_{-\infty}^{\infty} f(x-t)g(t) \, dt$ is:
\[
\widehat{f * g} = \sqrt{2\pi} \cdot \hat{f} \cdot \hat{g}
\]
\end{theorem}

\end{document}
